% Book cover, front (page 1) and back (page 2), A4.
% Build: cd book/cover && python make_cover_art.py && latexmk -pdf cover.tex
\documentclass{article}
\usepackage[paperwidth=210mm,paperheight=297mm,margin=0pt]{geometry}
\usepackage[T1]{fontenc}
\usepackage{libertinus}
\usepackage[type1,lining]{raleway}
\usepackage[scale=0.9]{plex-mono}
\usepackage{xcolor,tikz}
\usetikzlibrary{fadings}
\pagestyle{empty}
\hyphenation{MATLAB}
% ---------------------------------------------------------------------------
% fuzzyclust visual identity: colors, display font and the logo mark.
% Shared by the book (preamble.tex), the cover, the banner and the software
% documentation. Colors: Pannon blue for structure, the Okabe--Ito triple
% (blue, vermilion, bluish green) of the logo and of every data figure.
% ---------------------------------------------------------------------------
\definecolor{fcnavy}{HTML}{14304A}   % cover and dark backgrounds
\definecolor{fcbrand}{HTML}{285780}  % headings, rules, links
\definecolor{fcink}{HTML}{1F2A36}    % running text in headings
\definecolor{fcblue}{HTML}{0072B2}
\definecolor{fcverm}{HTML}{D55E00}
\definecolor{fcgreen}{HTML}{009E73}
\definecolor{fcpurple}{HTML}{CC79A7}
\definecolor{fcorange}{HTML}{E69F00}
\definecolor{fcsky}{HTML}{56B4E9}
\definecolor{fcgray}{HTML}{6B7785}
\definecolor{fcpaper}{HTML}{F4F6F8}
\definecolor{fclight}{HTML}{9CC3E6}  % light accent on dark backgrounds

% display face: Raleway (Type 1 via the raleway package)
\providecommand{\fcdisplay}{\fontfamily{Raleway-TLF}\selectfont}
\providecommand{\fcblack}{\fontfamily{Raleway-TLF}\fontseries{k}\selectfont}
\providecommand{\fcbold}{\fontfamily{Raleway-TLF}\fontseries{b}\selectfont}
\providecommand{\fcsemi}{\fontfamily{Raleway-TLF}\fontseries{sb}\selectfont}
\providecommand{\fcmedium}{\fontfamily{Raleway-TLF}\fontseries{mb}\selectfont}
\providecommand{\fclightface}{\fontfamily{Raleway-TLF}\fontseries{l}\selectfont}

% logo mark: three membership clouds; \fclogo[<scale>]{<opacity factor>}
% drawn with its center at the origin, unit = 1cm at scale 1
\newcommand{\fclogocloud}[3]{% x, y, color
  \foreach \r/\o in {1.25/0.18,1.0/0.32,0.68/0.55,0.36/0.9}{%
    \fill[#3, fill opacity=\o] (#1,#2) circle (\r);}}
\newcommand{\fclogo}[1][1]{%
  \begin{tikzpicture}[scale=#1, every node/.style={transform shape}]
    \fclogocloud{-0.9}{-0.52}{fcblue}
    \fclogocloud{0.9}{-0.52}{fcverm}
    \fclogocloud{0}{1.04}{fcgreen}
  \end{tikzpicture}}

\tikzfading[name=artfade, top color=transparent!0, bottom color=transparent!100]
\begin{document}
% ---------------------------------------------------------------- front cover
\noindent\begin{tikzpicture}[remember picture, overlay, x=1cm, y=1cm]
  \begin{scope}[shift={(current page.south west)}]
    \fill[fcnavy] (0,0) rectangle (21,29.7);
    % artwork: a Gustafson--Kessel partition computed by fuzzyclust
    \begin{scope}[shift={(-0.3,10.2)}]
      \clip (0.3,0) rectangle (21.3,19.5);
      % generated by book/cover/make_cover_art.py -- do not edit
\draw[line width=0.45pt, draw={rgb,255:red,58;green,154;blue,217}, draw opacity=0.30] (4.71,-2.49) -- (4.74,-2.28) -- (4.78,-2.00) -- (4.81,-1.78) -- (4.85,-1.50) -- (4.88,-1.28) -- (4.91,-1.01) -- (4.95,-0.76) -- (4.98,-0.52) -- (5.02,-0.26) -- (5.06,-0.02) -- (5.09,0.21) -- (5.14,0.47) -- (5.17,0.68) -- (5.23,0.96) -- (5.27,1.17) -- (5.32,1.39) -- (5.38,1.63) -- (5.45,1.88) -- (5.51,2.09) -- (5.58,2.30) -- (5.65,2.51) -- (5.74,2.72) -- (5.81,2.90) -- (5.89,3.08) -- (6.00,3.29) -- (6.10,3.47) -- (6.19,3.64) -- (6.31,3.83) -- (6.42,3.99) -- (6.53,4.14) -- (6.67,4.33) -- (6.80,4.49) -- (6.92,4.63) -- (7.05,4.77) -- (7.18,4.91) -- (7.32,5.05) -- (7.47,5.19) -- (7.63,5.33) -- (7.79,5.47) -- (7.96,5.61) -- (8.11,5.73) -- (8.25,5.84) -- (8.41,5.97) -- (8.60,6.11) -- (8.75,6.21) -- (8.91,6.32) -- (9.11,6.46) -- (9.26,6.55) -- (9.44,6.67) -- (9.62,6.78) -- (9.77,6.88) -- (9.97,7.01) -- (10.12,7.10) -- (10.32,7.23) -- (10.48,7.34) -- (10.64,7.44) -- (10.84,7.58) -- (10.98,7.68) -- (11.13,7.80) -- (11.30,7.94) -- (11.46,8.08) -- (11.60,8.22) -- (11.73,8.36) -- (11.84,8.50) -- (11.98,8.71) -- (12.07,8.85) -- (12.17,9.06) -- (12.26,9.28) -- (12.33,9.49) -- (12.40,9.70) -- (12.45,9.91) -- (12.49,10.12) -- (12.55,10.40) -- (12.58,10.61) -- (12.62,10.90) -- (12.65,11.11) -- (12.69,11.39) -- (12.72,11.60) -- (12.76,11.88) -- (12.78,12.09) -- (12.82,12.38) -- (12.85,12.63) -- (12.87,12.87) -- (12.91,13.15) -- (12.93,13.36) -- (12.95,13.64) -- (12.97,13.93) -- (12.98,14.21) -- (12.97,14.49) -- (12.92,14.73) -- (12.84,14.91) -- (12.70,15.07) -- (12.53,15.19) -- (12.34,15.28) -- (12.13,15.36) -- (11.91,15.43) -- (11.70,15.48) -- (11.41,15.54) -- (11.19,15.58) -- (10.97,15.62) -- (10.69,15.66) -- (10.48,15.69) -- (10.19,15.72) -- (9.90,15.76) -- (9.69,15.78) -- (9.40,15.80) -- (9.11,15.82) -- (8.90,15.84) -- (8.61,15.86) -- (8.32,15.87) -- (8.04,15.88) -- (7.75,15.89) -- (7.46,15.89) -- (7.17,15.89) -- (6.89,15.89) -- (6.60,15.88) -- (6.31,15.87) -- (6.03,15.86) -- (5.74,15.85) -- (5.45,15.83) -- (5.24,15.81) -- (4.95,15.79) -- (4.66,15.76) -- (4.45,15.74) -- (4.16,15.70) -- (3.94,15.68) -- (3.66,15.64) -- (3.44,15.61) -- (3.15,15.56) -- (2.94,15.52) -- (2.67,15.48) -- (2.44,15.43) -- (2.22,15.39) -- (1.97,15.33) -- (1.72,15.28) -- (1.50,15.23) -- (1.29,15.17) -- (1.07,15.12) -- (0.82,15.05) -- (0.57,14.98) -- (0.35,14.92) -- (0.14,14.86) -- (-0.08,14.79) -- (-0.29,14.72) -- (-0.51,14.65) -- (-0.72,14.58) -- (-0.94,14.50) -- (-1.15,14.42) -- (-1.35,14.35) -- (-1.53,14.28) -- (-1.73,14.20) -- (-1.94,14.11);
\draw[line width=0.45pt, draw={rgb,255:red,58;green,154;blue,217}, draw opacity=0.30] (14.63,20.69) -- (14.62,20.62) -- (14.60,20.55) -- (14.58,20.48) -- (14.57,20.41) -- (14.57,20.41) -- (14.55,20.34) -- (14.54,20.27) -- (14.52,20.20) -- (14.50,20.12) -- (14.50,20.10) -- (14.49,20.05) -- (14.47,19.98) -- (14.46,19.91) -- (14.44,19.84) -- (14.43,19.77) -- (14.43,19.77) -- (14.41,19.70) -- (14.39,19.63) -- (14.38,19.56) -- (14.36,19.49) -- (14.35,19.45) -- (14.35,19.42) -- (14.33,19.35) -- (14.32,19.28) -- (14.30,19.21) -- (14.29,19.14) -- (14.28,19.11) -- (14.27,19.07) -- (14.26,19.00) -- (14.24,18.93) -- (14.23,18.86) -- (14.21,18.79) -- (14.21,18.77) -- (14.20,18.72) -- (14.19,18.65) -- (14.17,18.57) -- (14.16,18.50) -- (14.14,18.43) -- (14.14,18.41) -- (14.13,18.36) -- (14.12,18.29) -- (14.10,18.22) -- (14.09,18.15) -- (14.08,18.08) -- (14.07,18.03) -- (14.06,18.01) -- (14.05,17.94) -- (14.04,17.87) -- (14.03,17.80) -- (14.01,17.73) -- (14.00,17.66) -- (14.00,17.62) -- (13.99,17.59) -- (13.98,17.52) -- (13.97,17.45) -- (13.96,17.38) -- (13.95,17.31) -- (13.94,17.24) -- (13.93,17.17) -- (13.92,17.14) -- (13.92,17.10) -- (13.91,17.03) -- (13.90,16.95) -- (13.89,16.88) -- (13.89,16.81) -- (13.88,16.74) -- (13.88,16.67) -- (13.87,16.60) -- (13.87,16.53) -- (13.87,16.46) -- (13.87,16.39) -- (13.87,16.32) -- (13.88,16.25) -- (13.89,16.18) -- (13.90,16.11) -- (13.91,16.04) -- (13.92,16.01) -- (13.94,15.97) -- (13.96,15.90) -- (14.00,15.84) -- (14.00,15.83) -- (14.05,15.76) -- (14.07,15.74) -- (14.11,15.69) -- (14.14,15.66) -- (14.19,15.62) -- (14.21,15.61) -- (14.28,15.56) -- (14.30,15.55) -- (14.35,15.51) -- (14.43,15.48) -- (14.43,15.48) -- (14.50,15.44) -- (14.57,15.41) -- (14.59,15.40) -- (14.64,15.39) -- (14.71,15.36) -- (14.78,15.34) -- (14.79,15.33) -- (14.86,15.31) -- (14.93,15.29) -- (15.00,15.27) -- (15.01,15.26) -- (15.07,15.25) -- (15.14,15.23) -- (15.22,15.21) -- (15.27,15.19) -- (15.29,15.19) -- (15.36,15.17) -- (15.43,15.15) -- (15.50,15.14) -- (15.56,15.12) -- (15.57,15.12) -- (15.65,15.10) -- (15.72,15.09) -- (15.79,15.07) -- (15.86,15.05) -- (15.87,15.05) -- (15.93,15.04) -- (16.01,15.02) -- (16.08,15.01) -- (16.15,14.99) -- (16.21,14.98) -- (16.22,14.98) -- (16.29,14.97) -- (16.36,14.95) -- (16.44,14.94) -- (16.51,14.92) -- (16.58,14.91) -- (16.58,14.91) -- (16.65,14.90) -- (16.72,14.89) -- (16.80,14.87) -- (16.87,14.86) -- (16.94,14.85) -- (16.98,14.84) -- (17.01,14.84) -- (17.08,14.82) -- (17.15,14.81) -- (17.23,14.80) -- (17.30,14.79) -- (17.37,14.78) -- (17.43,14.77) -- (17.44,14.77) -- (17.51,14.76) -- (17.59,14.75) -- (17.66,14.74) -- (17.73,14.73) -- (17.80,14.72) -- (17.87,14.71) -- (17.94,14.70) -- (17.96,14.70) -- (18.02,14.69) -- (18.09,14.69) -- (18.16,14.68) -- (18.23,14.67) -- (18.30,14.66) -- (18.38,14.66) -- (18.45,14.65) -- (18.52,14.64) -- (18.59,14.64) -- (18.65,14.63) -- (18.66,14.63) -- (18.73,14.62) -- (18.81,14.62) -- (18.88,14.61) -- (18.95,14.61) -- (19.02,14.60) -- (19.09,14.60) -- (19.16,14.59) -- (19.24,14.59) -- (19.31,14.59) -- (19.38,14.58) -- (19.45,14.58) -- (19.52,14.58) -- (19.60,14.58) -- (19.67,14.57) -- (19.74,14.57) -- (19.81,14.57) -- (19.88,14.57) -- (19.95,14.57) -- (20.03,14.57) -- (20.10,14.57) -- (20.17,14.57) -- (20.24,14.57) -- (20.31,14.57) -- (20.39,14.57) -- (20.46,14.57) -- (20.53,14.57) -- (20.60,14.57) -- (20.67,14.58) -- (20.74,14.58) -- (20.82,14.58) -- (20.89,14.58) -- (20.96,14.59) -- (21.03,14.59) -- (21.10,14.59) -- (21.18,14.60) -- (21.25,14.60) -- (21.32,14.61) -- (21.39,14.61) -- (21.46,14.61) -- (21.53,14.62) -- (21.61,14.63) -- (21.67,14.63) -- (21.68,14.63) -- (21.75,14.64) -- (21.82,14.64) -- (21.89,14.65) -- (21.97,14.66) -- (22.04,14.66) -- (22.11,14.67) -- (22.18,14.68) -- (22.25,14.68) -- (22.32,14.69) -- (22.40,14.70) -- (22.42,14.70) -- (22.47,14.71) -- (22.54,14.71) -- (22.61,14.72) -- (22.68,14.73) -- (22.75,14.74) -- (22.83,14.75) -- (22.90,14.76) -- (22.97,14.77) -- (23.00,14.77) -- (23.04,14.78) -- (23.11,14.79) -- (23.19,14.80) -- (23.26,14.81) -- (23.33,14.82) -- (23.40,14.83) -- (23.47,14.84) -- (23.51,14.84) -- (23.54,14.85) -- (23.62,14.86) -- (23.69,14.87);
\draw[line width=0.45pt, draw={rgb,255:red,58;green,154;blue,217}, draw opacity=0.44] (0.38,-2.49) -- (0.50,-2.30) -- (0.59,-2.14) -- (0.71,-1.93) -- (0.80,-1.78) -- (0.92,-1.57) -- (1.01,-1.43) -- (1.13,-1.22) -- (1.22,-1.08) -- (1.35,-0.87) -- (1.43,-0.73) -- (1.56,-0.52) -- (1.65,-0.37) -- (1.78,-0.16) -- (1.87,-0.02) -- (2.00,0.18) -- (2.10,0.33) -- (2.22,0.51) -- (2.33,0.68) -- (2.44,0.84) -- (2.57,1.03) -- (2.67,1.17) -- (2.79,1.35) -- (2.92,1.53) -- (3.02,1.67) -- (3.15,1.84) -- (3.29,2.02) -- (3.40,2.16) -- (3.51,2.30) -- (3.66,2.48) -- (3.80,2.65) -- (3.92,2.80) -- (4.05,2.94) -- (4.17,3.08) -- (4.30,3.22) -- (4.45,3.37) -- (4.59,3.52) -- (4.73,3.66) -- (4.88,3.80) -- (5.02,3.94) -- (5.16,4.07) -- (5.31,4.21) -- (5.46,4.34) -- (5.62,4.49) -- (5.79,4.63) -- (5.95,4.76) -- (6.10,4.88) -- (6.24,4.99) -- (6.40,5.12) -- (6.59,5.26) -- (6.74,5.38) -- (6.89,5.48) -- (7.06,5.61) -- (7.25,5.74) -- (7.39,5.84) -- (7.56,5.97) -- (7.75,6.09) -- (7.89,6.19) -- (8.09,6.32) -- (8.25,6.43) -- (8.41,6.53) -- (8.61,6.66) -- (8.75,6.76) -- (8.95,6.88) -- (9.11,6.99) -- (9.27,7.09) -- (9.47,7.23) -- (9.62,7.32) -- (9.79,7.44) -- (9.97,7.57) -- (10.12,7.68) -- (10.28,7.80) -- (10.46,7.94) -- (10.62,8.07) -- (10.76,8.20) -- (10.91,8.34) -- (11.05,8.48) -- (11.19,8.64) -- (11.31,8.78) -- (11.41,8.93) -- (11.55,9.14) -- (11.63,9.28) -- (11.73,9.49) -- (11.82,9.70) -- (11.90,9.91) -- (11.97,10.12) -- (12.02,10.33) -- (12.07,10.54) -- (12.11,10.83) -- (12.14,11.04) -- (12.15,11.32) -- (12.14,11.60) -- (12.12,11.81) -- (12.06,12.08) -- (11.98,12.28) -- (11.90,12.45) -- (11.77,12.65) -- (11.63,12.80) -- (11.48,12.93) -- (11.34,13.03) -- (11.12,13.14) -- (10.93,13.22) -- (10.72,13.29) -- (10.48,13.35) -- (10.26,13.40) -- (9.99,13.43) -- (9.76,13.45) -- (9.47,13.47) -- (9.18,13.47) -- (8.90,13.47) -- (8.61,13.45) -- (8.36,13.43) -- (8.11,13.41) -- (7.82,13.37) -- (7.60,13.34) -- (7.33,13.29) -- (7.10,13.25) -- (6.89,13.20) -- (6.65,13.15) -- (6.38,13.08) -- (6.17,13.03) -- (5.95,12.97) -- (5.74,12.90) -- (5.52,12.83) -- (5.31,12.76) -- (5.09,12.69) -- (4.88,12.61) -- (4.66,12.52) -- (4.47,12.45) -- (4.30,12.37) -- (4.09,12.28) -- (3.87,12.18) -- (3.69,12.09) -- (3.51,12.01) -- (3.30,11.90) -- (3.14,11.81) -- (2.94,11.70) -- (2.75,11.60) -- (2.58,11.50) -- (2.40,11.39) -- (2.22,11.28) -- (2.06,11.18) -- (1.86,11.05) -- (1.72,10.95) -- (1.54,10.83) -- (1.36,10.70) -- (1.21,10.59) -- (1.05,10.47) -- (0.87,10.33) -- (0.71,10.21) -- (0.57,10.10) -- (0.42,9.98) -- (0.26,9.84) -- (0.09,9.70) -- (-0.08,9.56) -- (-0.22,9.43) -- (-0.37,9.30) -- (-0.51,9.17) -- (-0.65,9.04) -- (-0.80,8.91) -- (-0.94,8.78) -- (-1.09,8.64) -- (-1.24,8.50) -- (-1.38,8.36) -- (-1.53,8.22) -- (-1.68,8.08) -- (-1.83,7.94) -- (-1.97,7.80);
\draw[line width=0.45pt, draw={rgb,255:red,58;green,154;blue,217}, draw opacity=0.44] (20.59,20.69) -- (20.57,20.62) -- (20.56,20.55) -- (20.55,20.48) -- (20.53,20.41) -- (20.53,20.37) -- (20.52,20.34) -- (20.51,20.27) -- (20.51,20.20) -- (20.50,20.12) -- (20.50,20.05) -- (20.49,19.98) -- (20.49,19.91) -- (20.49,19.84) -- (20.50,19.77) -- (20.50,19.70) -- (20.51,19.63) -- (20.52,19.56) -- (20.53,19.53) -- (20.54,19.49) -- (20.56,19.42) -- (20.58,19.35) -- (20.60,19.30) -- (20.61,19.28) -- (20.65,19.21) -- (20.67,19.16) -- (20.69,19.14) -- (20.74,19.07) -- (20.74,19.06) -- (20.80,19.00) -- (20.82,18.98) -- (20.88,18.93) -- (20.89,18.92) -- (20.96,18.86) -- (20.97,18.86) -- (21.03,18.82) -- (21.09,18.79) -- (21.10,18.78) -- (21.18,18.74) -- (21.24,18.72) -- (21.25,18.72) -- (21.32,18.69) -- (21.39,18.67) -- (21.46,18.65) -- (21.48,18.65) -- (21.53,18.63) -- (21.61,18.62) -- (21.68,18.61) -- (21.75,18.60) -- (21.82,18.59) -- (21.89,18.59) -- (21.97,18.58) -- (22.04,18.58) -- (22.11,18.58) -- (22.18,18.58) -- (22.25,18.58) -- (22.32,18.58) -- (22.40,18.58) -- (22.47,18.59) -- (22.54,18.59) -- (22.61,18.60) -- (22.68,18.60) -- (22.75,18.61) -- (22.83,18.62) -- (22.90,18.63) -- (22.97,18.64) -- (23.04,18.65) -- (23.04,18.65) -- (23.11,18.66) -- (23.19,18.67) -- (23.26,18.68) -- (23.33,18.69) -- (23.40,18.70) -- (23.47,18.72) -- (23.48,18.72) -- (23.54,18.73) -- (23.62,18.74) -- (23.69,18.76);
\draw[line width=0.45pt, draw={rgb,255:red,58;green,154;blue,217}, draw opacity=0.58] (0.50,3.35) -- (0.78,3.34) -- (1.00,3.36) -- (1.29,3.41) -- (1.50,3.47) -- (1.72,3.53) -- (1.93,3.60) -- (2.15,3.68) -- (2.36,3.76) -- (2.58,3.85) -- (2.75,3.92) -- (2.94,4.00) -- (3.15,4.10) -- (3.37,4.20) -- (3.52,4.27) -- (3.73,4.38) -- (3.94,4.48) -- (4.09,4.56) -- (4.30,4.67) -- (4.49,4.77) -- (4.66,4.86) -- (4.88,4.97) -- (5.02,5.05) -- (5.24,5.17) -- (5.41,5.26) -- (5.59,5.36) -- (5.79,5.47) -- (5.95,5.57) -- (6.16,5.68) -- (6.31,5.77) -- (6.53,5.89) -- (6.67,5.98) -- (6.89,6.10) -- (7.03,6.19) -- (7.25,6.31) -- (7.39,6.40) -- (7.60,6.53) -- (7.75,6.62) -- (7.95,6.74) -- (8.11,6.84) -- (8.29,6.95) -- (8.47,7.06) -- (8.62,7.16) -- (8.83,7.30) -- (8.97,7.39) -- (9.15,7.51) -- (9.33,7.64) -- (9.47,7.75) -- (9.63,7.87) -- (9.82,8.01) -- (9.97,8.14) -- (10.12,8.26) -- (10.26,8.39) -- (10.41,8.53) -- (10.55,8.68) -- (10.69,8.84) -- (10.82,8.99) -- (10.93,9.14) -- (11.05,9.32) -- (11.15,9.49) -- (11.27,9.70) -- (11.34,9.86) -- (11.41,10.05) -- (11.48,10.27) -- (11.54,10.54) -- (11.57,10.76) -- (11.57,11.04) -- (11.54,11.25) -- (11.47,11.46) -- (11.36,11.67) -- (11.24,11.81) -- (11.07,11.95) -- (10.91,12.05) -- (10.69,12.14) -- (10.48,12.20) -- (10.26,12.24) -- (9.97,12.27) -- (9.69,12.28) -- (9.40,12.26) -- (9.11,12.24) -- (8.90,12.21) -- (8.66,12.16) -- (8.39,12.11) -- (8.18,12.06) -- (7.96,12.00) -- (7.75,11.94) -- (7.53,11.87) -- (7.32,11.79) -- (7.10,11.72) -- (6.89,11.63) -- (6.67,11.54) -- (6.49,11.46) -- (6.31,11.38) -- (6.10,11.27) -- (5.91,11.18) -- (5.74,11.09) -- (5.52,10.97) -- (5.38,10.88) -- (5.17,10.76) -- (5.02,10.66) -- (4.84,10.54) -- (4.66,10.42) -- (4.52,10.32) -- (4.33,10.19) -- (4.16,10.06) -- (4.01,9.95) -- (3.87,9.83) -- (3.71,9.70) -- (3.54,9.56) -- (3.39,9.42) -- (3.23,9.28) -- (3.08,9.14) -- (2.94,8.99) -- (2.80,8.85) -- (2.66,8.71) -- (2.53,8.57) -- (2.40,8.43) -- (2.28,8.29) -- (2.15,8.14) -- (2.00,7.97) -- (1.86,7.80) -- (1.75,7.66) -- (1.65,7.51) -- (1.50,7.32) -- (1.39,7.16) -- (1.29,7.02) -- (1.15,6.81) -- (1.05,6.67) -- (0.93,6.47) -- (0.84,6.32) -- (0.71,6.11) -- (0.63,5.97) -- (0.52,5.75) -- (0.42,5.57) -- (0.35,5.40) -- (0.25,5.19) -- (0.17,4.98) -- (0.09,4.77) -- (0.03,4.56) -- (-0.02,4.34) -- (-0.05,4.06) -- (-0.02,3.78) -- (0.07,3.60) -- (0.21,3.46) -- (0.42,3.37);
\draw[line width=0.45pt, draw={rgb,255:red,58;green,154;blue,217}, draw opacity=0.72] (4.45,6.03) -- (4.59,6.02) -- (4.73,6.03) -- (4.83,6.04) -- (4.95,6.05) -- (5.09,6.07) -- (5.24,6.10) -- (5.31,6.12) -- (5.45,6.15) -- (5.53,6.18) -- (5.67,6.22) -- (5.76,6.25) -- (5.88,6.29) -- (5.97,6.32) -- (6.10,6.36) -- (6.17,6.39) -- (6.31,6.45) -- (6.38,6.48) -- (6.50,6.53) -- (6.60,6.57) -- (6.67,6.60) -- (6.82,6.67) -- (6.89,6.70) -- (6.97,6.74) -- (7.10,6.81) -- (7.17,6.84) -- (7.25,6.88) -- (7.38,6.95) -- (7.46,6.99) -- (7.53,7.03) -- (7.64,7.09) -- (7.75,7.15) -- (7.82,7.19) -- (7.89,7.23) -- (8.01,7.30) -- (8.11,7.36) -- (8.18,7.41) -- (8.25,7.45) -- (8.35,7.51) -- (8.46,7.59) -- (8.54,7.64) -- (8.61,7.69) -- (8.68,7.73) -- (8.77,7.80) -- (8.87,7.87) -- (8.97,7.94) -- (9.04,7.99) -- (9.11,8.05) -- (9.18,8.10) -- (9.26,8.16) -- (9.33,8.22) -- (9.42,8.29) -- (9.50,8.36) -- (9.58,8.43) -- (9.66,8.50) -- (9.74,8.57) -- (9.81,8.64) -- (9.88,8.71) -- (9.95,8.78) -- (10.02,8.85) -- (10.08,8.92) -- (10.14,8.99) -- (10.20,9.06) -- (10.26,9.14) -- (10.33,9.23) -- (10.41,9.33) -- (10.47,9.42) -- (10.51,9.49) -- (10.55,9.56) -- (10.62,9.68) -- (10.67,9.77) -- (10.70,9.84) -- (10.76,9.98) -- (10.78,10.05) -- (10.82,10.19) -- (10.84,10.26) -- (10.86,10.40) -- (10.86,10.54) -- (10.84,10.69) -- (10.83,10.76) -- (10.77,10.90) -- (10.72,10.97) -- (10.67,11.04) -- (10.59,11.11) -- (10.48,11.18) -- (10.41,11.22) -- (10.32,11.25) -- (10.19,11.29) -- (10.05,11.31) -- (9.97,11.32) -- (9.83,11.33) -- (9.69,11.33) -- (9.54,11.32) -- (9.47,11.31) -- (9.33,11.30) -- (9.18,11.27) -- (9.06,11.25) -- (8.97,11.23) -- (8.83,11.19) -- (8.75,11.17) -- (8.61,11.13) -- (8.53,11.11) -- (8.39,11.06) -- (8.32,11.04) -- (8.18,10.98) -- (8.11,10.96) -- (7.96,10.90) -- (7.89,10.87) -- (7.80,10.83) -- (7.68,10.77) -- (7.60,10.74) -- (7.50,10.69) -- (7.39,10.63) -- (7.32,10.59) -- (7.22,10.54) -- (7.10,10.48) -- (7.03,10.44) -- (6.96,10.40) -- (6.85,10.33) -- (6.74,10.27) -- (6.67,10.22) -- (6.60,10.18) -- (6.51,10.12) -- (6.41,10.05) -- (6.31,9.98) -- (6.24,9.93) -- (6.17,9.88) -- (6.10,9.83) -- (6.01,9.77) -- (5.92,9.70) -- (5.83,9.63) -- (5.75,9.56) -- (5.67,9.49) -- (5.59,9.43) -- (5.52,9.37) -- (5.45,9.31) -- (5.38,9.24) -- (5.31,9.18) -- (5.24,9.11) -- (5.16,9.04) -- (5.09,8.97) -- (5.02,8.90) -- (4.95,8.82) -- (4.88,8.75) -- (4.80,8.67) -- (4.73,8.59) -- (4.66,8.51) -- (4.60,8.43) -- (4.54,8.36) -- (4.48,8.29) -- (4.43,8.22) -- (4.37,8.15) -- (4.30,8.05) -- (4.23,7.94) -- (4.18,7.87) -- (4.14,7.80) -- (4.09,7.71) -- (4.02,7.59) -- (3.98,7.51) -- (3.94,7.44) -- (3.88,7.30) -- (3.86,7.23) -- (3.81,7.09) -- (3.79,7.02) -- (3.76,6.88) -- (3.74,6.74) -- (3.75,6.60) -- (3.78,6.46) -- (3.81,6.39) -- (3.87,6.29) -- (3.94,6.21) -- (4.01,6.16) -- (4.12,6.11) -- (4.23,6.07) -- (4.37,6.04) -- (4.45,6.03);
\draw[line width=0.45pt, draw={rgb,255:red,58;green,154;blue,217}, draw opacity=0.86] (5.95,7.29) -- (6.03,7.27) -- (6.10,7.25) -- (6.17,7.24) -- (6.24,7.24) -- (6.31,7.24) -- (6.38,7.24) -- (6.46,7.25) -- (6.53,7.26) -- (6.60,7.27) -- (6.67,7.29) -- (6.73,7.30) -- (6.74,7.31) -- (6.82,7.33) -- (6.89,7.35) -- (6.96,7.37) -- (6.97,7.37) -- (7.03,7.39) -- (7.10,7.42) -- (7.17,7.44) -- (7.17,7.45) -- (7.25,7.48) -- (7.32,7.51) -- (7.34,7.51) -- (7.39,7.54) -- (7.46,7.57) -- (7.49,7.59) -- (7.53,7.60) -- (7.60,7.64) -- (7.64,7.66) -- (7.68,7.67) -- (7.75,7.71) -- (7.78,7.73) -- (7.82,7.75) -- (7.89,7.79) -- (7.90,7.80) -- (7.96,7.83) -- (8.03,7.87) -- (8.04,7.87) -- (8.11,7.92) -- (8.14,7.94) -- (8.18,7.96) -- (8.25,8.01) -- (8.25,8.01) -- (8.32,8.05) -- (8.36,8.08) -- (8.39,8.10) -- (8.46,8.15) -- (8.47,8.15) -- (8.54,8.20) -- (8.56,8.22) -- (8.61,8.25) -- (8.66,8.29) -- (8.68,8.31) -- (8.75,8.36) -- (8.75,8.36) -- (8.83,8.42) -- (8.84,8.43) -- (8.90,8.48) -- (8.92,8.50) -- (8.97,8.54) -- (9.01,8.57) -- (9.04,8.60) -- (9.09,8.64) -- (9.11,8.67) -- (9.16,8.71) -- (9.18,8.73) -- (9.24,8.78) -- (9.26,8.80) -- (9.31,8.85) -- (9.33,8.88) -- (9.37,8.92) -- (9.40,8.95) -- (9.44,8.99) -- (9.47,9.03) -- (9.50,9.06) -- (9.54,9.12) -- (9.56,9.14) -- (9.61,9.21) -- (9.62,9.21) -- (9.66,9.28) -- (9.69,9.31) -- (9.71,9.35) -- (9.76,9.42) -- (9.76,9.42) -- (9.80,9.49) -- (9.83,9.55) -- (9.84,9.56) -- (9.87,9.63) -- (9.90,9.70) -- (9.90,9.71) -- (9.92,9.77) -- (9.94,9.84) -- (9.95,9.91) -- (9.96,9.98) -- (9.96,10.05) -- (9.94,10.12) -- (9.92,10.19) -- (9.90,10.22) -- (9.87,10.26) -- (9.83,10.30) -- (9.79,10.33) -- (9.76,10.35) -- (9.69,10.39) -- (9.64,10.40) -- (9.62,10.41) -- (9.54,10.42) -- (9.47,10.43) -- (9.40,10.43) -- (9.33,10.43) -- (9.26,10.43) -- (9.18,10.42) -- (9.11,10.41) -- (9.08,10.40) -- (9.04,10.40) -- (8.97,10.38) -- (8.90,10.36) -- (8.83,10.34) -- (8.79,10.33) -- (8.75,10.32) -- (8.68,10.30) -- (8.61,10.28) -- (8.57,10.26) -- (8.54,10.25) -- (8.47,10.22) -- (8.39,10.19) -- (8.39,10.19) -- (8.32,10.16) -- (8.25,10.13) -- (8.23,10.12) -- (8.18,10.10) -- (8.11,10.06) -- (8.08,10.05) -- (8.04,10.03) -- (7.96,9.99) -- (7.95,9.98) -- (7.89,9.95) -- (7.82,9.91) -- (7.82,9.91) -- (7.75,9.87) -- (7.70,9.84) -- (7.68,9.83) -- (7.60,9.78) -- (7.58,9.77) -- (7.53,9.74) -- (7.47,9.70) -- (7.46,9.69) -- (7.39,9.65) -- (7.36,9.63) -- (7.32,9.60) -- (7.26,9.56) -- (7.25,9.55) -- (7.17,9.49) -- (7.16,9.49) -- (7.10,9.44) -- (7.07,9.42) -- (7.03,9.39) -- (6.98,9.35) -- (6.96,9.33) -- (6.89,9.28) -- (6.89,9.27) -- (6.82,9.21) -- (6.81,9.21) -- (6.74,9.15) -- (6.73,9.14) -- (6.67,9.08) -- (6.65,9.06) -- (6.60,9.02) -- (6.57,8.99) -- (6.53,8.95) -- (6.50,8.92) -- (6.46,8.88) -- (6.43,8.85) -- (6.38,8.80) -- (6.36,8.78) -- (6.31,8.73) -- (6.30,8.71) -- (6.24,8.64) -- (6.24,8.64) -- (6.18,8.57) -- (6.17,8.56) -- (6.12,8.50) -- (6.10,8.47) -- (6.07,8.43) -- (6.03,8.37) -- (6.02,8.36) -- (5.97,8.29) -- (5.95,8.26) -- (5.93,8.22) -- (5.89,8.15) -- (5.88,8.14) -- (5.85,8.08) -- (5.82,8.01) -- (5.81,7.99) -- (5.79,7.94) -- (5.77,7.87) -- (5.75,7.80) -- (5.74,7.74) -- (5.74,7.73) -- (5.73,7.66) -- (5.73,7.59) -- (5.74,7.56) -- (5.75,7.51) -- (5.78,7.44) -- (5.81,7.40) -- (5.83,7.37) -- (5.88,7.33) -- (5.93,7.30) -- (5.95,7.29);
\draw[line width=0.45pt, draw={rgb,255:red,240;green,123;blue,63}, draw opacity=0.30] (-2.09,5.99) -- (-2.06,6.04) -- (-2.02,6.11) -- (-2.02,6.12) -- (-1.98,6.18) -- (-1.95,6.25) -- (-1.94,6.25) -- (-1.91,6.32) -- (-1.87,6.39) -- (-1.87,6.39) -- (-1.84,6.46) -- (-1.80,6.53) -- (-1.80,6.54) -- (-1.77,6.60) -- (-1.74,6.67) -- (-1.73,6.69) -- (-1.71,6.74) -- (-1.68,6.81) -- (-1.66,6.86) -- (-1.65,6.88) -- (-1.62,6.95) -- (-1.59,7.02) -- (-1.59,7.03) -- (-1.56,7.09) -- (-1.54,7.16) -- (-1.51,7.23) -- (-1.51,7.23) -- (-1.49,7.30) -- (-1.47,7.37) -- (-1.44,7.44) -- (-1.44,7.45) -- (-1.42,7.51) -- (-1.40,7.59) -- (-1.39,7.66) -- (-1.37,7.71) -- (-1.37,7.73) -- (-1.35,7.80) -- (-1.33,7.87) -- (-1.32,7.94) -- (-1.31,8.01) -- (-1.30,8.05) -- (-1.29,8.08) -- (-1.28,8.15) -- (-1.27,8.22) -- (-1.26,8.29) -- (-1.26,8.36) -- (-1.25,8.43) -- (-1.24,8.50) -- (-1.24,8.57) -- (-1.24,8.64) -- (-1.23,8.71) -- (-1.23,8.78) -- (-1.23,8.85) -- (-1.23,8.92) -- (-1.24,8.99) -- (-1.24,9.06) -- (-1.25,9.14) -- (-1.25,9.21) -- (-1.26,9.28) -- (-1.27,9.35) -- (-1.28,9.42) -- (-1.29,9.49) -- (-1.30,9.51) -- (-1.31,9.56) -- (-1.32,9.63) -- (-1.34,9.70) -- (-1.36,9.77) -- (-1.37,9.82) -- (-1.38,9.84) -- (-1.40,9.91) -- (-1.42,9.98) -- (-1.44,10.05) -- (-1.44,10.06) -- (-1.46,10.12) -- (-1.49,10.19) -- (-1.51,10.25) -- (-1.52,10.26) -- (-1.55,10.33) -- (-1.58,10.40) -- (-1.59,10.42) -- (-1.61,10.47) -- (-1.64,10.54) -- (-1.66,10.58) -- (-1.67,10.61) -- (-1.71,10.69) -- (-1.73,10.72) -- (-1.75,10.76) -- (-1.78,10.83) -- (-1.80,10.86) -- (-1.82,10.90) -- (-1.86,10.97) -- (-1.87,10.98) -- (-1.91,11.04) -- (-1.94,11.10) -- (-1.95,11.11) -- (-2.00,11.18) -- (-2.02,11.21) -- (-2.04,11.25) -- (-2.09,11.32) -- (-2.09,11.32);
\draw[line width=0.45pt, draw={rgb,255:red,240;green,123;blue,63}, draw opacity=0.30] (23.69,11.34) -- (23.55,11.25) -- (23.45,11.18) -- (23.33,11.09) -- (23.19,10.99) -- (23.05,10.90) -- (22.95,10.83) -- (22.83,10.74) -- (22.68,10.63) -- (22.56,10.54) -- (22.46,10.47) -- (22.32,10.37) -- (22.18,10.27) -- (22.08,10.19) -- (21.97,10.11) -- (21.82,10.01) -- (21.69,9.91) -- (21.59,9.84) -- (21.46,9.74) -- (21.32,9.64) -- (21.21,9.56) -- (21.10,9.48) -- (20.96,9.38) -- (20.82,9.28) -- (20.71,9.21) -- (20.60,9.13) -- (20.46,9.03) -- (20.31,8.94) -- (20.18,8.85) -- (20.07,8.78) -- (19.95,8.71) -- (19.81,8.62) -- (19.67,8.54) -- (19.52,8.46) -- (19.38,8.38) -- (19.24,8.30) -- (19.09,8.23) -- (18.95,8.16) -- (18.81,8.10) -- (18.66,8.04) -- (18.52,7.98) -- (18.38,7.93) -- (18.20,7.87) -- (18.02,7.81) -- (17.87,7.77) -- (17.71,7.73) -- (17.51,7.68) -- (17.37,7.66) -- (17.15,7.62) -- (16.94,7.59) -- (16.80,7.58) -- (16.58,7.57) -- (16.36,7.56) -- (16.15,7.56) -- (15.93,7.57) -- (15.72,7.58) -- (15.57,7.59) -- (15.36,7.61) -- (15.14,7.64) -- (15.00,7.66) -- (14.78,7.70) -- (14.64,7.73) -- (14.43,7.77) -- (14.28,7.80) -- (14.07,7.84) -- (13.92,7.88) -- (13.71,7.92) -- (13.56,7.96) -- (13.35,8.00) -- (13.21,8.03) -- (12.99,8.07) -- (12.85,8.09) -- (12.63,8.12) -- (12.42,8.14) -- (12.20,8.14) -- (11.98,8.14) -- (11.77,8.12) -- (11.55,8.09) -- (11.41,8.07) -- (11.19,8.02) -- (11.05,7.98) -- (10.91,7.94) -- (10.70,7.87) -- (10.55,7.81) -- (10.41,7.76) -- (10.26,7.70) -- (10.12,7.63) -- (9.97,7.57) -- (9.83,7.50) -- (9.69,7.43) -- (9.54,7.36) -- (9.40,7.29) -- (9.26,7.22) -- (9.11,7.14) -- (8.97,7.07) -- (8.83,6.99) -- (8.68,6.91) -- (8.54,6.83) -- (8.39,6.75) -- (8.25,6.67) -- (8.13,6.60) -- (8.00,6.53) -- (7.88,6.46) -- (7.75,6.38) -- (7.60,6.30) -- (7.46,6.21) -- (7.32,6.13) -- (7.17,6.04) -- (7.05,5.97) -- (6.93,5.89) -- (6.82,5.82) -- (6.67,5.74) -- (6.53,5.65) -- (6.38,5.56) -- (6.24,5.47) -- (6.13,5.40) -- (6.02,5.33) -- (5.88,5.24) -- (5.74,5.15) -- (5.59,5.06) -- (5.47,4.98) -- (5.37,4.91) -- (5.24,4.82) -- (5.09,4.73) -- (4.95,4.63) -- (4.84,4.56) -- (4.73,4.48) -- (4.59,4.38) -- (4.45,4.28) -- (4.34,4.20) -- (4.23,4.13) -- (4.09,4.02) -- (3.95,3.92) -- (3.86,3.85) -- (3.73,3.75) -- (3.58,3.64) -- (3.49,3.57) -- (3.37,3.48) -- (3.22,3.36) -- (3.14,3.29) -- (3.01,3.18) -- (2.88,3.08) -- (2.79,3.00) -- (2.65,2.87) -- (2.56,2.80) -- (2.44,2.68) -- (2.34,2.58) -- (2.22,2.47) -- (2.12,2.37) -- (2.00,2.26) -- (1.91,2.16) -- (1.79,2.04) -- (1.71,1.95) -- (1.58,1.81) -- (1.50,1.72) -- (1.40,1.60) -- (1.29,1.46) -- (1.21,1.38) -- (1.11,1.25) -- (1.00,1.10) -- (0.93,1.00) -- (0.85,0.89) -- (0.75,0.75) -- (0.66,0.61) -- (0.57,0.48) -- (0.50,0.36) -- (0.42,0.24) -- (0.35,0.12) -- (0.27,-0.02) -- (0.19,-0.16) -- (0.12,-0.30) -- (0.04,-0.45) -- (-0.03,-0.59) -- (-0.09,-0.73) -- (-0.16,-0.87) -- (-0.22,-1.01) -- (-0.29,-1.18) -- (-0.37,-1.36) -- (-0.42,-1.50) -- (-0.47,-1.64) -- (-0.52,-1.78) -- (-0.58,-1.95) -- (-0.64,-2.14) -- (-0.68,-2.28) -- (-0.73,-2.42);
\draw[line width=0.45pt, draw={rgb,255:red,240;green,123;blue,63}, draw opacity=0.44] (11.14,-2.49) -- (11.19,-2.48) -- (11.27,-2.47) -- (11.34,-2.45) -- (11.41,-2.44) -- (11.48,-2.42) -- (11.51,-2.42) -- (11.55,-2.41) -- (11.63,-2.39) -- (11.70,-2.38) -- (11.77,-2.36) -- (11.83,-2.35) -- (11.84,-2.35) -- (11.91,-2.33) -- (11.98,-2.31) -- (12.06,-2.29) -- (12.13,-2.28) -- (12.13,-2.28) -- (12.20,-2.26) -- (12.27,-2.24) -- (12.34,-2.22) -- (12.40,-2.21) -- (12.42,-2.20) -- (12.49,-2.18) -- (12.56,-2.16) -- (12.63,-2.14) -- (12.65,-2.14) -- (12.70,-2.12) -- (12.77,-2.10) -- (12.85,-2.08) -- (12.89,-2.07) -- (12.92,-2.06) -- (12.99,-2.04) -- (13.06,-2.01) -- (13.12,-2.00) -- (13.13,-1.99) -- (13.21,-1.97) -- (13.28,-1.94) -- (13.34,-1.92) -- (13.35,-1.92) -- (13.42,-1.90) -- (13.49,-1.87) -- (13.55,-1.85) -- (13.56,-1.85) -- (13.64,-1.82) -- (13.71,-1.80) -- (13.75,-1.78) -- (13.78,-1.77) -- (13.85,-1.75) -- (13.92,-1.72) -- (13.95,-1.71) -- (14.00,-1.70) -- (14.07,-1.67) -- (14.14,-1.65) -- (14.15,-1.64) -- (14.21,-1.62) -- (14.28,-1.59) -- (14.34,-1.57) -- (14.35,-1.57) -- (14.43,-1.54) -- (14.50,-1.51) -- (14.52,-1.50) -- (14.57,-1.48) -- (14.64,-1.46) -- (14.70,-1.43) -- (14.71,-1.43) -- (14.78,-1.40) -- (14.86,-1.37) -- (14.88,-1.36) -- (14.93,-1.34) -- (15.00,-1.31) -- (15.06,-1.29) -- (15.07,-1.29) -- (15.14,-1.26) -- (15.22,-1.23) -- (15.23,-1.22) -- (15.29,-1.20) -- (15.36,-1.17) -- (15.41,-1.15) -- (15.43,-1.14) -- (15.50,-1.11) -- (15.57,-1.08) -- (15.58,-1.08) -- (15.65,-1.05) -- (15.72,-1.02) -- (15.75,-1.01) -- (15.79,-0.99) -- (15.86,-0.96) -- (15.92,-0.94) -- (15.93,-0.93) -- (16.01,-0.90) -- (16.08,-0.87) -- (16.08,-0.87) -- (16.15,-0.84) -- (16.22,-0.81) -- (16.25,-0.80) -- (16.29,-0.78) -- (16.36,-0.75) -- (16.42,-0.73) -- (16.44,-0.72) -- (16.51,-0.69) -- (16.58,-0.66) -- (16.59,-0.66) -- (16.65,-0.63) -- (16.72,-0.60) -- (16.75,-0.59) -- (16.80,-0.57) -- (16.87,-0.54) -- (16.92,-0.52) -- (16.94,-0.51) -- (17.01,-0.48) -- (17.08,-0.45) -- (17.09,-0.45) -- (17.15,-0.42) -- (17.23,-0.39) -- (17.26,-0.37) -- (17.30,-0.36) -- (17.37,-0.33) -- (17.43,-0.30) -- (17.44,-0.30) -- (17.51,-0.27) -- (17.59,-0.24) -- (17.60,-0.23) -- (17.66,-0.21) -- (17.73,-0.18) -- (17.77,-0.16) -- (17.80,-0.15) -- (17.87,-0.12) -- (17.94,-0.10) -- (17.95,-0.09) -- (18.02,-0.07) -- (18.09,-0.04) -- (18.13,-0.02) -- (18.16,-0.01) -- (18.23,0.02) -- (18.30,0.04) -- (18.31,0.05) -- (18.38,0.07) -- (18.45,0.10) -- (18.50,0.12) -- (18.52,0.12) -- (18.59,0.15) -- (18.66,0.18) -- (18.70,0.19) -- (18.73,0.20) -- (18.81,0.23) -- (18.88,0.25) -- (18.90,0.26) -- (18.95,0.28) -- (19.02,0.30) -- (19.09,0.33) -- (19.10,0.33) -- (19.16,0.35) -- (19.24,0.37) -- (19.31,0.40) -- (19.32,0.40) -- (19.38,0.42) -- (19.45,0.44) -- (19.52,0.46) -- (19.54,0.47) -- (19.60,0.49) -- (19.67,0.51) -- (19.74,0.53) -- (19.79,0.54) -- (19.81,0.55) -- (19.88,0.57) -- (19.95,0.59) -- (20.03,0.61) -- (20.04,0.61) -- (20.10,0.63) -- (20.17,0.64) -- (20.24,0.66) -- (20.31,0.68) -- (20.32,0.68) -- (20.39,0.70) -- (20.46,0.71) -- (20.53,0.73) -- (20.60,0.75) -- (20.63,0.75) -- (20.67,0.76) -- (20.74,0.78) -- (20.82,0.79) -- (20.89,0.80) -- (20.96,0.82) -- (20.99,0.82) -- (21.03,0.83) -- (21.10,0.84) -- (21.18,0.86) -- (21.25,0.87) -- (21.32,0.88) -- (21.39,0.89) -- (21.41,0.89) -- (21.46,0.90) -- (21.53,0.91) -- (21.61,0.92) -- (21.68,0.93) -- (21.75,0.94) -- (21.82,0.95) -- (21.89,0.95) -- (21.97,0.96) -- (21.99,0.96) -- (22.04,0.97) -- (22.11,0.97) -- (22.18,0.98) -- (22.25,0.99) -- (22.32,0.99) -- (22.40,1.00) -- (22.47,1.00) -- (22.54,1.01) -- (22.61,1.01) -- (22.68,1.01) -- (22.75,1.02) -- (22.83,1.02) -- (22.90,1.02) -- (22.97,1.02) -- (23.04,1.03) -- (23.11,1.03) -- (23.19,1.03) -- (23.26,1.03) -- (23.33,1.03) -- (23.40,1.03) -- (23.47,1.03) -- (23.54,1.03) -- (23.62,1.03) -- (23.69,1.03);
\draw[line width=0.45pt, draw={rgb,255:red,240;green,123;blue,63}, draw opacity=0.44] (23.69,8.26) -- (23.54,8.21) -- (23.37,8.15) -- (23.19,8.08) -- (23.04,8.03) -- (22.90,7.98) -- (22.75,7.93) -- (22.58,7.87) -- (22.40,7.80) -- (22.25,7.75) -- (22.11,7.71) -- (21.95,7.66) -- (21.75,7.59) -- (21.61,7.55) -- (21.46,7.50) -- (21.25,7.44) -- (21.10,7.40) -- (20.96,7.36) -- (20.74,7.31) -- (20.60,7.27) -- (20.42,7.23) -- (20.24,7.19) -- (20.09,7.16) -- (19.88,7.12) -- (19.70,7.09) -- (19.52,7.06) -- (19.31,7.03) -- (19.16,7.02) -- (18.95,6.99) -- (18.73,6.98) -- (18.52,6.96) -- (18.33,6.95) -- (18.16,6.95) -- (17.94,6.94) -- (17.73,6.94) -- (17.51,6.95) -- (17.35,6.95) -- (17.15,6.96) -- (16.94,6.97) -- (16.72,6.99) -- (16.51,7.01) -- (16.36,7.03) -- (16.15,7.05) -- (15.93,7.08) -- (15.79,7.10) -- (15.57,7.13) -- (15.36,7.16) -- (15.22,7.18) -- (15.00,7.22) -- (14.86,7.25) -- (14.64,7.28) -- (14.50,7.31) -- (14.28,7.35) -- (14.12,7.37) -- (13.92,7.41) -- (13.71,7.44) -- (13.56,7.47) -- (13.35,7.50) -- (13.21,7.52) -- (12.99,7.54) -- (12.77,7.56) -- (12.56,7.57) -- (12.34,7.58) -- (12.13,7.57) -- (11.91,7.56) -- (11.70,7.54) -- (11.51,7.51) -- (11.34,7.49) -- (11.15,7.44) -- (10.98,7.40) -- (10.84,7.36) -- (10.64,7.30) -- (10.48,7.25) -- (10.33,7.19) -- (10.19,7.14) -- (10.05,7.08) -- (9.90,7.02) -- (9.75,6.95) -- (9.60,6.88) -- (9.46,6.81) -- (9.32,6.74) -- (9.18,6.67) -- (9.04,6.60) -- (8.90,6.52) -- (8.75,6.44) -- (8.61,6.35) -- (8.47,6.27) -- (8.32,6.18) -- (8.21,6.11) -- (8.10,6.04) -- (7.96,5.95) -- (7.82,5.85) -- (7.68,5.75) -- (7.58,5.68) -- (7.46,5.60) -- (7.32,5.49) -- (7.20,5.40) -- (7.10,5.32) -- (6.96,5.21) -- (6.85,5.12) -- (6.74,5.02) -- (6.61,4.91) -- (6.53,4.83) -- (6.38,4.70) -- (6.31,4.63) -- (6.17,4.49) -- (6.10,4.41) -- (5.97,4.27) -- (5.88,4.17) -- (5.79,4.06) -- (5.68,3.92) -- (5.59,3.82) -- (5.52,3.71) -- (5.42,3.57) -- (5.32,3.43) -- (5.24,3.29) -- (5.16,3.16) -- (5.09,3.03) -- (5.02,2.88) -- (4.95,2.73) -- (4.89,2.58) -- (4.83,2.44) -- (4.79,2.30) -- (4.73,2.12) -- (4.69,1.95) -- (4.66,1.81) -- (4.63,1.60) -- (4.60,1.39) -- (4.59,1.17) -- (4.60,0.96) -- (4.61,0.75) -- (4.64,0.54) -- (4.67,0.40) -- (4.72,0.19) -- (4.76,0.05) -- (4.81,-0.09) -- (4.88,-0.27) -- (4.95,-0.42) -- (5.02,-0.56) -- (5.09,-0.68) -- (5.17,-0.80) -- (5.27,-0.94) -- (5.38,-1.08) -- (5.45,-1.17) -- (5.57,-1.29) -- (5.67,-1.39) -- (5.79,-1.50) -- (5.88,-1.58) -- (6.03,-1.69) -- (6.16,-1.78) -- (6.26,-1.85) -- (6.38,-1.93) -- (6.53,-2.01) -- (6.67,-2.08) -- (6.82,-2.15) -- (6.96,-2.21) -- (7.12,-2.28) -- (7.32,-2.35) -- (7.46,-2.39) -- (7.60,-2.43) -- (7.82,-2.49);
\draw[line width=0.45pt, draw={rgb,255:red,240;green,123;blue,63}, draw opacity=0.58] (10.12,0.82) -- (10.33,0.81) -- (10.55,0.80) -- (10.76,0.80) -- (10.98,0.81) -- (11.19,0.82) -- (11.34,0.83) -- (11.55,0.85) -- (11.77,0.88) -- (11.91,0.90) -- (12.13,0.93) -- (12.32,0.96) -- (12.49,1.00) -- (12.67,1.03) -- (12.85,1.07) -- (12.99,1.11) -- (13.21,1.16) -- (13.35,1.20) -- (13.51,1.25) -- (13.71,1.30) -- (13.85,1.35) -- (14.00,1.39) -- (14.18,1.46) -- (14.35,1.52) -- (14.50,1.57) -- (14.64,1.62) -- (14.78,1.68) -- (14.93,1.74) -- (15.10,1.81) -- (15.27,1.88) -- (15.43,1.95) -- (15.57,2.02) -- (15.72,2.08) -- (15.86,2.15) -- (16.01,2.22) -- (16.15,2.30) -- (16.29,2.37) -- (16.43,2.44) -- (16.56,2.51) -- (16.68,2.58) -- (16.80,2.65) -- (16.94,2.74) -- (17.08,2.83) -- (17.23,2.92) -- (17.36,3.01) -- (17.46,3.08) -- (17.59,3.17) -- (17.73,3.28) -- (17.83,3.36) -- (17.94,3.45) -- (18.08,3.57) -- (18.16,3.64) -- (18.30,3.78) -- (18.38,3.87) -- (18.48,3.99) -- (18.58,4.13) -- (18.66,4.27) -- (18.73,4.42) -- (18.77,4.56) -- (18.80,4.77) -- (18.77,4.98) -- (18.72,5.12) -- (18.64,5.26) -- (18.53,5.40) -- (18.45,5.48) -- (18.30,5.61) -- (18.20,5.68) -- (18.09,5.76) -- (17.94,5.84) -- (17.80,5.92) -- (17.66,5.99) -- (17.51,6.05) -- (17.37,6.11) -- (17.20,6.18) -- (17.01,6.25) -- (16.87,6.30) -- (16.72,6.34) -- (16.58,6.39) -- (16.36,6.45) -- (16.22,6.49) -- (16.08,6.53) -- (15.86,6.59) -- (15.72,6.63) -- (15.54,6.67) -- (15.36,6.71) -- (15.22,6.74) -- (15.00,6.79) -- (14.86,6.82) -- (14.64,6.86) -- (14.50,6.89) -- (14.28,6.93) -- (14.12,6.95) -- (13.92,6.98) -- (13.71,7.01) -- (13.56,7.03) -- (13.35,7.05) -- (13.13,7.07) -- (12.92,7.08) -- (12.70,7.09) -- (12.49,7.09) -- (12.27,7.08) -- (12.06,7.07) -- (11.84,7.05) -- (11.63,7.02) -- (11.48,7.00) -- (11.27,6.95) -- (11.12,6.92) -- (10.96,6.88) -- (10.76,6.82) -- (10.62,6.78) -- (10.48,6.73) -- (10.32,6.67) -- (10.14,6.60) -- (9.98,6.53) -- (9.83,6.46) -- (9.69,6.39) -- (9.54,6.32) -- (9.41,6.25) -- (9.29,6.18) -- (9.17,6.11) -- (9.04,6.03) -- (8.90,5.94) -- (8.75,5.84) -- (8.63,5.75) -- (8.54,5.68) -- (8.39,5.57) -- (8.27,5.47) -- (8.18,5.39) -- (8.04,5.26) -- (7.96,5.19) -- (7.82,5.05) -- (7.75,4.97) -- (7.63,4.84) -- (7.53,4.71) -- (7.46,4.61) -- (7.37,4.49) -- (7.29,4.34) -- (7.21,4.20) -- (7.14,4.06) -- (7.09,3.92) -- (7.03,3.75) -- (6.99,3.57) -- (6.96,3.38) -- (6.95,3.22) -- (6.96,3.01) -- (6.97,2.87) -- (7.02,2.65) -- (7.07,2.51) -- (7.14,2.37) -- (7.21,2.23) -- (7.30,2.09) -- (7.39,1.98) -- (7.48,1.88) -- (7.60,1.75) -- (7.70,1.67) -- (7.82,1.57) -- (7.96,1.47) -- (8.10,1.39) -- (8.23,1.32) -- (8.37,1.25) -- (8.53,1.17) -- (8.68,1.11) -- (8.83,1.06) -- (8.97,1.02) -- (9.18,0.96) -- (9.33,0.93) -- (9.53,0.89) -- (9.69,0.87) -- (9.90,0.84) -- (10.10,0.82);
\draw[line width=0.45pt, draw={rgb,255:red,240;green,123;blue,63}, draw opacity=0.72] (11.34,2.72) -- (11.48,2.72) -- (11.63,2.72) -- (11.70,2.72) -- (11.84,2.73) -- (11.98,2.74) -- (12.13,2.75) -- (12.27,2.76) -- (12.42,2.78) -- (12.53,2.80) -- (12.63,2.81) -- (12.77,2.83) -- (12.92,2.86) -- (12.99,2.87) -- (13.13,2.90) -- (13.28,2.94) -- (13.35,2.95) -- (13.49,2.99) -- (13.56,3.01) -- (13.71,3.05) -- (13.80,3.08) -- (13.92,3.11) -- (14.02,3.15) -- (14.14,3.19) -- (14.22,3.22) -- (14.35,3.27) -- (14.43,3.30) -- (14.57,3.35) -- (14.64,3.39) -- (14.74,3.43) -- (14.86,3.49) -- (14.93,3.52) -- (15.03,3.57) -- (15.14,3.63) -- (15.22,3.67) -- (15.29,3.72) -- (15.39,3.78) -- (15.50,3.85) -- (15.57,3.90) -- (15.65,3.96) -- (15.72,4.01) -- (15.79,4.07) -- (15.86,4.13) -- (15.94,4.20) -- (16.01,4.27) -- (16.08,4.35) -- (16.15,4.45) -- (16.22,4.55) -- (16.26,4.63) -- (16.30,4.70) -- (16.35,4.84) -- (16.36,4.91) -- (16.38,5.05) -- (16.36,5.19) -- (16.35,5.26) -- (16.29,5.40) -- (16.25,5.47) -- (16.20,5.54) -- (16.15,5.61) -- (16.08,5.69) -- (16.00,5.75) -- (15.91,5.82) -- (15.81,5.89) -- (15.72,5.95) -- (15.65,6.00) -- (15.57,6.04) -- (15.43,6.10) -- (15.36,6.14) -- (15.26,6.18) -- (15.14,6.22) -- (15.07,6.25) -- (14.93,6.29) -- (14.85,6.32) -- (14.71,6.36) -- (14.60,6.39) -- (14.50,6.41) -- (14.35,6.44) -- (14.28,6.46) -- (14.14,6.48) -- (14.00,6.51) -- (13.86,6.53) -- (13.78,6.54) -- (13.64,6.55) -- (13.49,6.57) -- (13.35,6.58) -- (13.21,6.59) -- (13.06,6.59) -- (12.92,6.60) -- (12.77,6.60) -- (12.63,6.59) -- (12.49,6.59) -- (12.34,6.58) -- (12.20,6.56) -- (12.06,6.55) -- (11.92,6.53) -- (11.84,6.52) -- (11.70,6.49) -- (11.55,6.46) -- (11.48,6.45) -- (11.34,6.41) -- (11.24,6.39) -- (11.12,6.35) -- (11.00,6.32) -- (10.91,6.29) -- (10.80,6.25) -- (10.69,6.21) -- (10.62,6.18) -- (10.48,6.12) -- (10.41,6.08) -- (10.30,6.04) -- (10.19,5.98) -- (10.12,5.94) -- (10.03,5.89) -- (9.92,5.82) -- (9.83,5.77) -- (9.76,5.72) -- (9.69,5.67) -- (9.61,5.61) -- (9.52,5.54) -- (9.43,5.47) -- (9.35,5.40) -- (9.28,5.33) -- (9.21,5.26) -- (9.15,5.19) -- (9.09,5.12) -- (9.03,5.05) -- (8.97,4.95) -- (8.90,4.84) -- (8.86,4.77) -- (8.83,4.68) -- (8.78,4.56) -- (8.75,4.43) -- (8.74,4.34) -- (8.74,4.20) -- (8.75,4.06) -- (8.77,3.99) -- (8.81,3.85) -- (8.85,3.78) -- (8.90,3.69) -- (8.97,3.59) -- (9.04,3.50) -- (9.11,3.43) -- (9.18,3.36) -- (9.26,3.31) -- (9.33,3.26) -- (9.40,3.21) -- (9.50,3.15) -- (9.62,3.09) -- (9.69,3.06) -- (9.80,3.01) -- (9.90,2.97) -- (10.00,2.94) -- (10.12,2.90) -- (10.24,2.87) -- (10.33,2.84) -- (10.48,2.81) -- (10.58,2.80) -- (10.69,2.78) -- (10.84,2.76) -- (10.98,2.74) -- (11.12,2.73) -- (11.27,2.73) -- (11.34,2.72);
\draw[line width=0.45pt, draw={rgb,255:red,240;green,123;blue,63}, draw opacity=0.86] (11.55,3.98) -- (11.63,3.97) -- (11.70,3.96) -- (11.77,3.95) -- (11.84,3.95) -- (11.91,3.94) -- (11.98,3.94) -- (12.06,3.94) -- (12.13,3.93) -- (12.20,3.93) -- (12.27,3.93) -- (12.34,3.94) -- (12.42,3.94) -- (12.49,3.94) -- (12.56,3.95) -- (12.63,3.96) -- (12.70,3.96) -- (12.77,3.97) -- (12.85,3.98) -- (12.92,3.99) -- (12.93,3.99) -- (12.99,4.00) -- (13.06,4.01) -- (13.13,4.03) -- (13.21,4.04) -- (13.28,4.06) -- (13.29,4.06) -- (13.35,4.08) -- (13.42,4.09) -- (13.49,4.11) -- (13.56,4.13) -- (13.56,4.13) -- (13.64,4.16) -- (13.71,4.18) -- (13.78,4.20) -- (13.78,4.20) -- (13.85,4.23) -- (13.92,4.26) -- (13.96,4.27) -- (14.00,4.29) -- (14.07,4.32) -- (14.12,4.34) -- (14.14,4.35) -- (14.21,4.39) -- (14.26,4.42) -- (14.28,4.43) -- (14.35,4.47) -- (14.38,4.49) -- (14.43,4.52) -- (14.49,4.56) -- (14.50,4.56) -- (14.57,4.62) -- (14.58,4.63) -- (14.64,4.68) -- (14.66,4.70) -- (14.71,4.75) -- (14.73,4.77) -- (14.78,4.83) -- (14.80,4.84) -- (14.85,4.91) -- (14.86,4.93) -- (14.89,4.98) -- (14.92,5.05) -- (14.93,5.10) -- (14.94,5.12) -- (14.95,5.19) -- (14.94,5.26) -- (14.93,5.33) -- (14.93,5.35) -- (14.91,5.40) -- (14.87,5.47) -- (14.86,5.50) -- (14.83,5.54) -- (14.78,5.59) -- (14.76,5.61) -- (14.71,5.66) -- (14.68,5.68) -- (14.64,5.71) -- (14.58,5.75) -- (14.57,5.76) -- (14.50,5.80) -- (14.45,5.82) -- (14.43,5.84) -- (14.35,5.87) -- (14.29,5.89) -- (14.28,5.90) -- (14.21,5.92) -- (14.14,5.95) -- (14.07,5.97) -- (14.07,5.97) -- (14.00,5.99) -- (13.92,6.00) -- (13.85,6.02) -- (13.78,6.03) -- (13.76,6.04) -- (13.71,6.04) -- (13.64,6.05) -- (13.56,6.06) -- (13.49,6.07) -- (13.42,6.08) -- (13.35,6.08) -- (13.28,6.09) -- (13.21,6.09) -- (13.13,6.09) -- (13.06,6.10) -- (12.99,6.10) -- (12.92,6.09) -- (12.85,6.09) -- (12.77,6.09) -- (12.70,6.08) -- (12.63,6.08) -- (12.56,6.07) -- (12.49,6.06) -- (12.42,6.05) -- (12.34,6.05) -- (12.28,6.04) -- (12.27,6.03) -- (12.20,6.02) -- (12.13,6.01) -- (12.06,5.99) -- (11.98,5.98) -- (11.93,5.97) -- (11.91,5.96) -- (11.84,5.94) -- (11.77,5.92) -- (11.70,5.90) -- (11.68,5.89) -- (11.63,5.88) -- (11.55,5.85) -- (11.48,5.83) -- (11.47,5.82) -- (11.41,5.80) -- (11.34,5.77) -- (11.30,5.75) -- (11.27,5.74) -- (11.19,5.70) -- (11.15,5.68) -- (11.12,5.67) -- (11.05,5.63) -- (11.03,5.61) -- (10.98,5.58) -- (10.91,5.54) -- (10.91,5.54) -- (10.84,5.49) -- (10.81,5.47) -- (10.76,5.43) -- (10.73,5.40) -- (10.69,5.37) -- (10.65,5.33) -- (10.62,5.30) -- (10.59,5.26) -- (10.55,5.21) -- (10.53,5.19) -- (10.49,5.12) -- (10.48,5.10) -- (10.45,5.05) -- (10.42,4.98) -- (10.41,4.92) -- (10.40,4.91) -- (10.39,4.84) -- (10.39,4.77) -- (10.40,4.70) -- (10.41,4.69) -- (10.42,4.63) -- (10.46,4.56) -- (10.48,4.52) -- (10.50,4.49) -- (10.55,4.42) -- (10.56,4.42) -- (10.62,4.35) -- (10.63,4.34) -- (10.69,4.30) -- (10.72,4.27) -- (10.76,4.25) -- (10.84,4.21) -- (10.84,4.20) -- (10.91,4.17) -- (10.98,4.14) -- (10.99,4.13) -- (11.05,4.11) -- (11.12,4.08) -- (11.19,4.06) -- (11.19,4.06) -- (11.27,4.04) -- (11.34,4.02) -- (11.41,4.01) -- (11.48,3.99) -- (11.48,3.99) -- (11.55,3.98);
\draw[line width=0.45pt, draw={rgb,255:red,43;green,192;blue,142}, draw opacity=0.30] (-2.09,14.55) -- (-1.84,14.49) -- (-1.59,14.43) -- (-1.37,14.38) -- (-1.15,14.34) -- (-0.88,14.28) -- (-0.65,14.23) -- (-0.44,14.18) -- (-0.22,14.14) -- (0.07,14.08) -- (0.28,14.04) -- (0.50,13.99) -- (0.78,13.94) -- (1.00,13.90) -- (1.22,13.86) -- (1.50,13.80) -- (1.72,13.77) -- (2.00,13.72) -- (2.22,13.68) -- (2.44,13.64) -- (2.72,13.60) -- (2.94,13.56) -- (3.22,13.52) -- (3.44,13.49) -- (3.73,13.45) -- (3.94,13.42) -- (4.23,13.38) -- (4.45,13.36) -- (4.73,13.32) -- (5.02,13.29) -- (5.24,13.27) -- (5.52,13.24) -- (5.74,13.22) -- (6.03,13.19) -- (6.31,13.17) -- (6.53,13.15) -- (6.82,13.12) -- (7.10,13.10) -- (7.38,13.08) -- (7.60,13.06) -- (7.89,13.04) -- (8.18,13.02) -- (8.39,13.00) -- (8.68,12.97) -- (8.94,12.94) -- (9.18,12.91) -- (9.46,12.87) -- (9.69,12.83) -- (9.90,12.79) -- (10.15,12.73) -- (10.40,12.66) -- (10.59,12.59) -- (10.76,12.51) -- (10.98,12.40) -- (11.13,12.31) -- (11.30,12.16) -- (11.43,12.02) -- (11.55,11.84) -- (11.64,11.67) -- (11.71,11.46) -- (11.77,11.18) -- (11.78,10.97) -- (11.78,10.69) -- (11.77,10.44) -- (11.74,10.19) -- (11.70,9.91) -- (11.67,9.70) -- (11.64,9.42) -- (11.62,9.21) -- (11.60,8.92) -- (11.63,8.64) -- (11.68,8.43) -- (11.77,8.23) -- (11.88,8.08) -- (12.02,7.94) -- (12.20,7.81) -- (12.34,7.72) -- (12.56,7.62) -- (12.77,7.53) -- (12.99,7.45) -- (13.21,7.38) -- (13.42,7.32) -- (13.64,7.26) -- (13.85,7.20) -- (14.07,7.15) -- (14.29,7.09) -- (14.57,7.03) -- (14.78,6.98) -- (15.00,6.93) -- (15.22,6.88) -- (15.50,6.81) -- (15.72,6.76) -- (15.93,6.71) -- (16.15,6.66) -- (16.40,6.60) -- (16.65,6.54) -- (16.87,6.49) -- (17.08,6.43) -- (17.30,6.38) -- (17.55,6.32) -- (17.80,6.25) -- (18.02,6.20) -- (18.23,6.14) -- (18.45,6.08) -- (18.66,6.02) -- (18.88,5.96) -- (19.09,5.89) -- (19.32,5.82) -- (19.54,5.75) -- (19.75,5.68) -- (19.95,5.61) -- (20.17,5.52) -- (20.39,5.43) -- (20.58,5.33) -- (20.74,5.22) -- (20.89,5.10) -- (21.01,4.91) -- (20.98,4.63) -- (20.89,4.46) -- (20.75,4.27) -- (20.61,4.13) -- (20.47,3.99) -- (20.31,3.86) -- (20.17,3.74) -- (20.03,3.62) -- (19.87,3.50) -- (19.69,3.36) -- (19.52,3.23) -- (19.38,3.13) -- (19.21,3.01) -- (19.02,2.87) -- (18.88,2.76) -- (18.72,2.65) -- (18.52,2.51) -- (18.38,2.41) -- (18.23,2.30) -- (18.03,2.16) -- (17.87,2.05) -- (17.73,1.95) -- (17.53,1.81) -- (17.37,1.69) -- (17.23,1.59) -- (17.04,1.46) -- (16.87,1.33) -- (16.72,1.23) -- (16.56,1.10) -- (16.37,0.96) -- (16.22,0.85) -- (16.08,0.74) -- (15.90,0.61) -- (15.72,0.47) -- (15.57,0.35) -- (15.43,0.24) -- (15.28,0.12) -- (15.11,-0.02) -- (14.94,-0.16) -- (14.78,-0.30) -- (14.64,-0.42) -- (14.50,-0.55) -- (14.35,-0.67) -- (14.21,-0.80) -- (14.06,-0.94) -- (13.91,-1.08) -- (13.76,-1.22) -- (13.62,-1.36) -- (13.47,-1.50) -- (13.33,-1.64) -- (13.19,-1.78) -- (13.06,-1.92) -- (12.92,-2.07) -- (12.77,-2.23) -- (12.63,-2.38);
\draw[line width=0.45pt, draw={rgb,255:red,43;green,192;blue,142}, draw opacity=0.30] (23.69,18.69) -- (23.65,18.72) -- (23.62,18.74) -- (23.55,18.79) -- (23.54,18.79) -- (23.47,18.84) -- (23.45,18.86) -- (23.40,18.90) -- (23.37,18.93) -- (23.33,18.96) -- (23.28,19.00) -- (23.26,19.02) -- (23.21,19.07) -- (23.19,19.09) -- (23.13,19.14) -- (23.11,19.16) -- (23.07,19.21) -- (23.04,19.24) -- (23.00,19.28) -- (22.97,19.32) -- (22.95,19.35) -- (22.90,19.41) -- (22.89,19.42) -- (22.84,19.49) -- (22.83,19.51) -- (22.80,19.56) -- (22.75,19.63) -- (22.75,19.63) -- (22.72,19.70) -- (22.68,19.77) -- (22.68,19.77) -- (22.65,19.84) -- (22.62,19.91) -- (22.61,19.93) -- (22.59,19.98) -- (22.57,20.05) -- (22.54,20.12) -- (22.54,20.14) -- (22.52,20.20) -- (22.51,20.27) -- (22.49,20.34) -- (22.48,20.41) -- (22.47,20.46) -- (22.46,20.48) -- (22.45,20.55) -- (22.45,20.62) -- (22.44,20.69);
\draw[line width=0.45pt, draw={rgb,255:red,43;green,192;blue,142}, draw opacity=0.44] (-2.09,20.50) -- (-1.87,20.37) -- (-1.66,20.23) -- (-1.44,20.10) -- (-1.23,19.97) -- (-1.01,19.84) -- (-0.79,19.70) -- (-0.55,19.56) -- (-0.32,19.42) -- (-0.08,19.28) -- (0.14,19.15) -- (0.35,19.02) -- (0.57,18.90) -- (0.78,18.77) -- (1.00,18.65) -- (1.24,18.50) -- (1.49,18.36) -- (1.72,18.24) -- (1.93,18.12) -- (2.15,18.00) -- (2.38,17.87) -- (2.63,17.73) -- (2.87,17.60) -- (3.08,17.49) -- (3.30,17.37) -- (3.56,17.24) -- (3.80,17.11) -- (4.01,17.00) -- (4.25,16.88) -- (4.52,16.75) -- (4.73,16.64) -- (4.97,16.53) -- (5.24,16.40) -- (5.45,16.30) -- (5.72,16.18) -- (5.95,16.07) -- (6.19,15.97) -- (6.46,15.85) -- (6.67,15.76) -- (6.96,15.64) -- (7.18,15.55) -- (7.46,15.43) -- (7.70,15.33) -- (7.96,15.23) -- (8.23,15.12) -- (8.47,15.03) -- (8.75,14.92) -- (8.97,14.84) -- (9.26,14.73) -- (9.50,14.63) -- (9.76,14.53) -- (10.01,14.42) -- (10.26,14.31) -- (10.48,14.21) -- (10.76,14.07) -- (10.98,13.94) -- (11.19,13.80) -- (11.40,13.64) -- (11.57,13.50) -- (11.76,13.29) -- (11.91,13.09) -- (12.04,12.87) -- (12.13,12.66) -- (12.22,12.38) -- (12.27,12.09) -- (12.31,11.74) -- (12.32,11.39) -- (12.32,11.04) -- (12.30,10.69) -- (12.28,10.33) -- (12.26,10.05) -- (12.26,9.70) -- (12.27,9.42) -- (12.33,9.06) -- (12.42,8.81) -- (12.55,8.57) -- (12.70,8.38) -- (12.88,8.22) -- (13.08,8.08) -- (13.33,7.94) -- (13.56,7.83) -- (13.81,7.73) -- (14.07,7.64) -- (14.35,7.55) -- (14.64,7.47) -- (14.93,7.40) -- (15.22,7.34) -- (15.50,7.29) -- (15.86,7.23) -- (16.15,7.20) -- (16.51,7.17) -- (16.80,7.16) -- (17.15,7.16) -- (17.44,7.17) -- (17.80,7.20) -- (18.09,7.24) -- (18.45,7.30) -- (18.73,7.37) -- (18.99,7.44) -- (19.24,7.53) -- (19.52,7.65) -- (19.74,7.74) -- (19.98,7.87) -- (20.23,8.01) -- (20.45,8.15) -- (20.65,8.29) -- (20.83,8.43) -- (21.03,8.60) -- (21.23,8.78) -- (21.39,8.94) -- (21.57,9.14) -- (21.75,9.35) -- (21.89,9.54) -- (22.04,9.76) -- (22.17,9.98) -- (22.28,10.19) -- (22.40,10.44) -- (22.49,10.69) -- (22.57,10.97) -- (22.63,11.25) -- (22.66,11.60) -- (22.65,11.95) -- (22.60,12.23) -- (22.53,12.52) -- (22.42,12.80) -- (22.31,13.01) -- (22.18,13.22) -- (22.02,13.43) -- (21.84,13.64) -- (21.68,13.80) -- (21.46,13.99) -- (21.28,14.14) -- (21.09,14.28) -- (20.88,14.42) -- (20.65,14.56) -- (20.40,14.70) -- (20.17,14.83) -- (19.95,14.94) -- (19.72,15.05) -- (19.45,15.18) -- (19.24,15.28) -- (18.98,15.40) -- (18.73,15.53) -- (18.52,15.63) -- (18.29,15.76) -- (18.04,15.90) -- (17.82,16.04) -- (17.62,16.18) -- (17.44,16.33) -- (17.23,16.53) -- (17.08,16.71) -- (16.94,16.91) -- (16.80,17.17) -- (16.70,17.38) -- (16.61,17.66) -- (16.54,17.94) -- (16.49,18.22) -- (16.45,18.57) -- (16.43,18.86) -- (16.41,19.21) -- (16.41,19.56) -- (16.42,19.91) -- (16.43,20.27) -- (16.44,20.55);
\draw[line width=0.45pt, draw={rgb,255:red,43;green,192;blue,142}, draw opacity=0.58] (15.36,7.80) -- (15.50,7.78) -- (15.65,7.76) -- (15.79,7.75) -- (15.93,7.74) -- (16.08,7.74) -- (16.22,7.74) -- (16.36,7.74) -- (16.51,7.75) -- (16.65,7.77) -- (16.80,7.79) -- (16.87,7.80) -- (17.01,7.82) -- (17.15,7.86) -- (17.23,7.87) -- (17.37,7.92) -- (17.44,7.94) -- (17.59,7.99) -- (17.66,8.02) -- (17.80,8.07) -- (17.87,8.11) -- (17.96,8.15) -- (18.09,8.22) -- (18.16,8.26) -- (18.23,8.30) -- (18.33,8.36) -- (18.43,8.43) -- (18.52,8.49) -- (18.59,8.55) -- (18.66,8.61) -- (18.73,8.67) -- (18.81,8.73) -- (18.88,8.80) -- (18.95,8.88) -- (19.02,8.95) -- (19.09,9.04) -- (19.16,9.13) -- (19.23,9.21) -- (19.28,9.28) -- (19.32,9.35) -- (19.38,9.44) -- (19.45,9.56) -- (19.49,9.63) -- (19.52,9.70) -- (19.59,9.84) -- (19.62,9.91) -- (19.67,10.05) -- (19.69,10.12) -- (19.73,10.26) -- (19.74,10.33) -- (19.77,10.47) -- (19.79,10.61) -- (19.80,10.76) -- (19.80,10.90) -- (19.79,11.04) -- (19.78,11.18) -- (19.75,11.32) -- (19.74,11.39) -- (19.70,11.53) -- (19.67,11.64) -- (19.63,11.74) -- (19.60,11.82) -- (19.53,11.95) -- (19.50,12.02) -- (19.45,12.10) -- (19.38,12.22) -- (19.32,12.31) -- (19.27,12.38) -- (19.21,12.45) -- (19.15,12.52) -- (19.09,12.59) -- (19.02,12.66) -- (18.95,12.73) -- (18.88,12.80) -- (18.81,12.87) -- (18.72,12.94) -- (18.64,13.01) -- (18.55,13.08) -- (18.45,13.15) -- (18.38,13.20) -- (18.30,13.25) -- (18.23,13.30) -- (18.13,13.36) -- (18.02,13.43) -- (17.94,13.47) -- (17.87,13.51) -- (17.75,13.57) -- (17.66,13.62) -- (17.59,13.65) -- (17.45,13.71) -- (17.37,13.75) -- (17.29,13.78) -- (17.15,13.84) -- (17.08,13.87) -- (16.94,13.92) -- (16.87,13.94) -- (16.72,13.99) -- (16.65,14.02) -- (16.51,14.06) -- (16.44,14.08) -- (16.29,14.12) -- (16.22,14.14) -- (16.08,14.17) -- (15.93,14.20) -- (15.86,14.21) -- (15.72,14.24) -- (15.57,14.26) -- (15.43,14.28) -- (15.36,14.28) -- (15.22,14.29) -- (15.07,14.30) -- (14.93,14.29) -- (14.78,14.29) -- (14.71,14.28) -- (14.57,14.26) -- (14.43,14.23) -- (14.35,14.21) -- (14.21,14.16) -- (14.14,14.13) -- (14.01,14.07) -- (13.92,14.01) -- (13.85,13.96) -- (13.78,13.91) -- (13.71,13.84) -- (13.64,13.77) -- (13.56,13.69) -- (13.49,13.59) -- (13.43,13.50) -- (13.39,13.43) -- (13.35,13.36) -- (13.28,13.22) -- (13.25,13.15) -- (13.21,13.04) -- (13.17,12.94) -- (13.13,12.84) -- (13.10,12.73) -- (13.06,12.59) -- (13.04,12.52) -- (13.01,12.38) -- (12.99,12.30) -- (12.96,12.16) -- (12.94,12.02) -- (12.92,11.92) -- (12.90,11.81) -- (12.88,11.67) -- (12.86,11.53) -- (12.85,11.42) -- (12.83,11.32) -- (12.82,11.18) -- (12.81,11.04) -- (12.80,10.90) -- (12.79,10.76) -- (12.79,10.61) -- (12.78,10.47) -- (12.78,10.33) -- (12.79,10.19) -- (12.80,10.05) -- (12.81,9.91) -- (12.83,9.77) -- (12.85,9.68) -- (12.87,9.56) -- (12.91,9.42) -- (12.94,9.35) -- (12.99,9.22) -- (13.03,9.14) -- (13.07,9.06) -- (13.13,8.96) -- (13.21,8.86) -- (13.26,8.78) -- (13.33,8.71) -- (13.40,8.64) -- (13.47,8.57) -- (13.56,8.50) -- (13.64,8.44) -- (13.71,8.39) -- (13.78,8.35) -- (13.88,8.29) -- (14.00,8.23) -- (14.07,8.19) -- (14.15,8.15) -- (14.28,8.09) -- (14.35,8.06) -- (14.50,8.01) -- (14.57,7.99) -- (14.71,7.94) -- (14.78,7.92) -- (14.93,7.88) -- (15.00,7.87) -- (15.14,7.84) -- (15.29,7.81) -- (15.36,7.80);
\draw[line width=0.45pt, draw={rgb,255:red,43;green,192;blue,142}, draw opacity=0.72] (15.22,8.42) -- (15.29,8.41) -- (15.36,8.40) -- (15.43,8.39) -- (15.50,8.38) -- (15.57,8.37) -- (15.65,8.37) -- (15.72,8.37) -- (15.79,8.37) -- (15.86,8.37) -- (15.93,8.38) -- (16.01,8.38) -- (16.08,8.39) -- (16.15,8.40) -- (16.22,8.41) -- (16.29,8.43) -- (16.30,8.43) -- (16.36,8.45) -- (16.44,8.47) -- (16.51,8.49) -- (16.55,8.50) -- (16.58,8.51) -- (16.65,8.54) -- (16.72,8.57) -- (16.73,8.57) -- (16.80,8.60) -- (16.87,8.64) -- (16.87,8.64) -- (16.94,8.68) -- (17.00,8.71) -- (17.01,8.72) -- (17.08,8.77) -- (17.10,8.78) -- (17.15,8.82) -- (17.20,8.85) -- (17.23,8.88) -- (17.28,8.92) -- (17.30,8.94) -- (17.36,8.99) -- (17.37,9.01) -- (17.43,9.06) -- (17.44,9.08) -- (17.49,9.14) -- (17.51,9.16) -- (17.55,9.21) -- (17.59,9.26) -- (17.60,9.28) -- (17.65,9.35) -- (17.66,9.36) -- (17.69,9.42) -- (17.73,9.48) -- (17.73,9.49) -- (17.77,9.56) -- (17.80,9.63) -- (17.80,9.63) -- (17.83,9.70) -- (17.86,9.77) -- (17.87,9.81) -- (17.88,9.84) -- (17.90,9.91) -- (17.92,9.98) -- (17.93,10.05) -- (17.94,10.10) -- (17.95,10.12) -- (17.96,10.19) -- (17.96,10.26) -- (17.97,10.33) -- (17.97,10.40) -- (17.97,10.47) -- (17.97,10.54) -- (17.96,10.61) -- (17.95,10.69) -- (17.94,10.75) -- (17.94,10.76) -- (17.93,10.83) -- (17.92,10.90) -- (17.90,10.97) -- (17.88,11.04) -- (17.87,11.05) -- (17.85,11.11) -- (17.83,11.18) -- (17.80,11.24) -- (17.80,11.25) -- (17.77,11.32) -- (17.73,11.39) -- (17.73,11.40) -- (17.69,11.46) -- (17.66,11.52) -- (17.65,11.53) -- (17.61,11.60) -- (17.59,11.64) -- (17.56,11.67) -- (17.51,11.74) -- (17.51,11.74) -- (17.45,11.81) -- (17.44,11.83) -- (17.40,11.88) -- (17.37,11.91) -- (17.33,11.95) -- (17.30,11.99) -- (17.26,12.02) -- (17.23,12.06) -- (17.19,12.09) -- (17.15,12.13) -- (17.11,12.16) -- (17.08,12.19) -- (17.03,12.23) -- (17.01,12.25) -- (16.94,12.30) -- (16.94,12.31) -- (16.87,12.36) -- (16.84,12.38) -- (16.80,12.41) -- (16.73,12.45) -- (16.72,12.45) -- (16.65,12.50) -- (16.62,12.52) -- (16.58,12.54) -- (16.51,12.58) -- (16.49,12.59) -- (16.44,12.61) -- (16.36,12.65) -- (16.34,12.66) -- (16.29,12.68) -- (16.22,12.71) -- (16.17,12.73) -- (16.15,12.74) -- (16.08,12.76) -- (16.01,12.79) -- (15.96,12.80) -- (15.93,12.81) -- (15.86,12.83) -- (15.79,12.84) -- (15.72,12.86) -- (15.67,12.87) -- (15.65,12.87) -- (15.57,12.89) -- (15.50,12.90) -- (15.43,12.90) -- (15.36,12.91) -- (15.29,12.91) -- (15.22,12.91) -- (15.14,12.91) -- (15.07,12.91) -- (15.00,12.90) -- (14.93,12.90) -- (14.86,12.88) -- (14.78,12.87) -- (14.78,12.87) -- (14.71,12.85) -- (14.64,12.83) -- (14.57,12.81) -- (14.54,12.80) -- (14.50,12.78) -- (14.43,12.75) -- (14.38,12.73) -- (14.35,12.72) -- (14.28,12.68) -- (14.25,12.66) -- (14.21,12.63) -- (14.15,12.59) -- (14.14,12.58) -- (14.07,12.53) -- (14.06,12.52) -- (14.00,12.46) -- (13.98,12.45) -- (13.92,12.39) -- (13.91,12.38) -- (13.85,12.31) -- (13.85,12.31) -- (13.79,12.23) -- (13.78,12.22) -- (13.74,12.16) -- (13.71,12.11) -- (13.70,12.09) -- (13.65,12.02) -- (13.64,11.99) -- (13.62,11.95) -- (13.58,11.88) -- (13.56,11.84) -- (13.55,11.81) -- (13.52,11.74) -- (13.50,11.67) -- (13.49,11.66) -- (13.47,11.60) -- (13.45,11.53) -- (13.43,11.46) -- (13.42,11.42) -- (13.41,11.39) -- (13.40,11.32) -- (13.38,11.25) -- (13.37,11.18) -- (13.36,11.11) -- (13.35,11.04) -- (13.35,11.01) -- (13.34,10.97) -- (13.34,10.90) -- (13.34,10.83) -- (13.33,10.76) -- (13.33,10.69) -- (13.33,10.61) -- (13.34,10.54) -- (13.34,10.47) -- (13.35,10.40) -- (13.35,10.38) -- (13.35,10.33) -- (13.36,10.26) -- (13.38,10.19) -- (13.39,10.12) -- (13.41,10.05) -- (13.42,9.99) -- (13.42,9.98) -- (13.44,9.91) -- (13.47,9.84) -- (13.49,9.77) -- (13.49,9.77) -- (13.52,9.70) -- (13.55,9.63) -- (13.56,9.60) -- (13.59,9.56) -- (13.62,9.49) -- (13.64,9.46) -- (13.66,9.42) -- (13.71,9.35) -- (13.71,9.35) -- (13.76,9.28) -- (13.78,9.25) -- (13.82,9.21) -- (13.85,9.17) -- (13.88,9.14) -- (13.92,9.09) -- (13.95,9.06) -- (14.00,9.02) -- (14.02,8.99) -- (14.07,8.95) -- (14.10,8.92) -- (14.14,8.90) -- (14.19,8.85) -- (14.21,8.84) -- (14.28,8.79) -- (14.30,8.78) -- (14.35,8.75) -- (14.41,8.71) -- (14.43,8.71) -- (14.50,8.67) -- (14.55,8.64) -- (14.57,8.63) -- (14.64,8.60) -- (14.70,8.57) -- (14.71,8.57) -- (14.78,8.54) -- (14.86,8.51) -- (14.90,8.50) -- (14.93,8.49) -- (15.00,8.47) -- (15.07,8.45) -- (15.14,8.43) -- (15.16,8.43) -- (15.22,8.42);
\draw[line width=0.45pt, draw={rgb,255:red,43;green,192;blue,142}, draw opacity=0.86] (15.36,9.14) -- (15.43,9.13) -- (15.50,9.13) -- (15.57,9.13) -- (15.60,9.14) -- (15.65,9.14) -- (15.72,9.15) -- (15.79,9.17) -- (15.86,9.19) -- (15.91,9.21) -- (15.93,9.21) -- (16.01,9.24) -- (16.07,9.28) -- (16.08,9.28) -- (16.15,9.32) -- (16.18,9.35) -- (16.22,9.38) -- (16.27,9.42) -- (16.29,9.44) -- (16.34,9.49) -- (16.36,9.51) -- (16.40,9.56) -- (16.44,9.60) -- (16.46,9.63) -- (16.50,9.70) -- (16.51,9.71) -- (16.54,9.77) -- (16.57,9.84) -- (16.58,9.87) -- (16.60,9.91) -- (16.62,9.98) -- (16.63,10.05) -- (16.65,10.12) -- (16.65,10.16) -- (16.66,10.19) -- (16.66,10.26) -- (16.66,10.33) -- (16.66,10.40) -- (16.65,10.47) -- (16.65,10.48) -- (16.64,10.54) -- (16.63,10.61) -- (16.61,10.69) -- (16.59,10.76) -- (16.58,10.80) -- (16.57,10.83) -- (16.54,10.90) -- (16.51,10.97) -- (16.51,10.98) -- (16.48,11.04) -- (16.44,11.11) -- (16.44,11.11) -- (16.39,11.18) -- (16.36,11.22) -- (16.35,11.25) -- (16.29,11.32) -- (16.29,11.32) -- (16.23,11.39) -- (16.22,11.40) -- (16.17,11.46) -- (16.15,11.48) -- (16.09,11.53) -- (16.08,11.54) -- (16.01,11.60) -- (16.01,11.60) -- (15.93,11.66) -- (15.92,11.67) -- (15.86,11.71) -- (15.81,11.74) -- (15.79,11.75) -- (15.72,11.79) -- (15.67,11.81) -- (15.65,11.83) -- (15.57,11.86) -- (15.50,11.88) -- (15.50,11.88) -- (15.43,11.91) -- (15.36,11.92) -- (15.29,11.94) -- (15.22,11.95) -- (15.20,11.95) -- (15.14,11.96) -- (15.07,11.96) -- (15.00,11.96) -- (14.93,11.96) -- (14.90,11.95) -- (14.86,11.95) -- (14.78,11.93) -- (14.71,11.91) -- (14.64,11.89) -- (14.62,11.88) -- (14.57,11.86) -- (14.50,11.82) -- (14.48,11.81) -- (14.43,11.78) -- (14.38,11.74) -- (14.35,11.72) -- (14.30,11.67) -- (14.28,11.66) -- (14.23,11.60) -- (14.21,11.58) -- (14.17,11.53) -- (14.14,11.48) -- (14.13,11.46) -- (14.09,11.39) -- (14.07,11.35) -- (14.05,11.32) -- (14.02,11.25) -- (14.00,11.18) -- (14.00,11.17) -- (13.98,11.11) -- (13.96,11.04) -- (13.95,10.97) -- (13.94,10.90) -- (13.94,10.83) -- (13.94,10.76) -- (13.94,10.69) -- (13.94,10.61) -- (13.95,10.54) -- (13.96,10.47) -- (13.97,10.40) -- (13.99,10.33) -- (14.00,10.31) -- (14.01,10.26) -- (14.03,10.19) -- (14.06,10.12) -- (14.07,10.11) -- (14.09,10.05) -- (14.13,9.98) -- (14.14,9.96) -- (14.17,9.91) -- (14.21,9.84) -- (14.21,9.84) -- (14.26,9.77) -- (14.28,9.74) -- (14.31,9.70) -- (14.35,9.65) -- (14.38,9.63) -- (14.43,9.58) -- (14.45,9.56) -- (14.50,9.51) -- (14.53,9.49) -- (14.57,9.45) -- (14.62,9.42) -- (14.64,9.40) -- (14.71,9.35) -- (14.72,9.35) -- (14.78,9.31) -- (14.85,9.28) -- (14.86,9.27) -- (14.93,9.24) -- (15.00,9.21) -- (15.03,9.21) -- (15.07,9.19) -- (15.14,9.17) -- (15.22,9.16) -- (15.29,9.14) -- (15.36,9.14) -- (15.36,9.14);
\fill[fill={rgb,255:red,58;green,154;blue,216}, fill opacity=0.94] (7.92,9.42) circle (0.094);
\fill[fill={rgb,255:red,58;green,154;blue,216}, fill opacity=0.92] (9.35,9.00) circle (0.092);
\fill[fill={rgb,255:red,58;green,154;blue,216}, fill opacity=0.90] (9.92,9.27) circle (0.089);
\fill[fill={rgb,255:red,59;green,154;blue,215}, fill opacity=0.89] (7.19,10.14) circle (0.087);
\fill[fill={rgb,255:red,58;green,154;blue,216}, fill opacity=0.92] (9.82,9.64) circle (0.091);
\fill[fill={rgb,255:red,58;green,153;blue,216}, fill opacity=0.93] (6.51,8.52) circle (0.092);
\fill[fill={rgb,255:red,58;green,153;blue,216}, fill opacity=0.93] (8.30,8.24) circle (0.092);
\fill[fill={rgb,255:red,58;green,154;blue,216}, fill opacity=0.92] (7.79,9.80) circle (0.092);
\fill[fill={rgb,255:red,52;green,170;blue,183}, fill opacity=0.75] (12.14,11.41) circle (0.070);
\fill[fill={rgb,255:red,43;green,191;blue,142}, fill opacity=0.91] (14.17,11.82) circle (0.090);
\fill[fill={rgb,255:red,45;green,187;blue,150}, fill opacity=0.84] (13.42,12.55) circle (0.081);
\fill[fill={rgb,255:red,55;green,163;blue,198}, fill opacity=0.79] (11.52,11.84) circle (0.075);
\fill[fill={rgb,255:red,58;green,153;blue,216}, fill opacity=0.95] (7.75,8.73) circle (0.095);
\fill[fill={rgb,255:red,45;green,187;blue,149}, fill opacity=0.84] (15.50,13.61) circle (0.081);
\fill[fill={rgb,255:red,58;green,154;blue,216}, fill opacity=0.94] (8.17,9.37) circle (0.094);
\fill[fill={rgb,255:red,47;green,181;blue,162}, fill opacity=0.79] (12.68,11.75) circle (0.075);
\fill[fill={rgb,255:red,57;green,157;blue,208}, fill opacity=0.83] (11.30,10.42) circle (0.080);
\fill[fill={rgb,255:red,61;green,153;blue,213}, fill opacity=0.89] (5.53,6.49) circle (0.087);
\fill[fill={rgb,255:red,58;green,154;blue,216}, fill opacity=0.91] (7.70,9.95) circle (0.090);
\fill[fill={rgb,255:red,58;green,154;blue,216}, fill opacity=0.92] (9.80,10.22) circle (0.091);
\fill[fill={rgb,255:red,58;green,153;blue,216}, fill opacity=0.95] (7.74,9.02) circle (0.095);
\fill[fill={rgb,255:red,55;green,161;blue,200}, fill opacity=0.80] (11.51,11.61) circle (0.077);
\fill[fill={rgb,255:red,81;green,150;blue,196}, fill opacity=0.83] (5.09,5.32) circle (0.079);
\fill[fill={rgb,255:red,58;green,153;blue,216}, fill opacity=0.91] (5.60,7.62) circle (0.091);
\fill[fill={rgb,255:red,58;green,157;blue,206}, fill opacity=0.82] (8.81,12.02) circle (0.079);
\fill[fill={rgb,255:red,62;green,153;blue,213}, fill opacity=0.89] (6.59,6.75) circle (0.087);
\fill[fill={rgb,255:red,58;green,154;blue,216}, fill opacity=0.93] (7.56,9.50) circle (0.092);
\fill[fill={rgb,255:red,58;green,154;blue,216}, fill opacity=0.92] (8.25,10.09) circle (0.091);
\fill[fill={rgb,255:red,58;green,154;blue,216}, fill opacity=0.92] (7.91,9.85) circle (0.092);
\fill[fill={rgb,255:red,66;green,152;blue,209}, fill opacity=0.86] (4.39,5.87) circle (0.084);
\fill[fill={rgb,255:red,58;green,153;blue,216}, fill opacity=0.95] (7.77,8.42) circle (0.095);
\fill[fill={rgb,255:red,58;green,153;blue,216}, fill opacity=0.91] (8.38,7.99) circle (0.090);
\fill[fill={rgb,255:red,58;green,154;blue,216}, fill opacity=0.92] (9.83,10.14) circle (0.091);
\fill[fill={rgb,255:red,61;green,153;blue,213}, fill opacity=0.88] (4.92,8.37) circle (0.087);
\fill[fill={rgb,255:red,51;green,173;blue,177}, fill opacity=0.74] (12.26,11.10) circle (0.069);
\fill[fill={rgb,255:red,90;green,148;blue,189}, fill opacity=0.81] (7.35,6.34) circle (0.077);
\fill[fill={rgb,255:red,58;green,154;blue,216}, fill opacity=0.93] (9.45,10.00) circle (0.093);
\fill[fill={rgb,255:red,62;green,153;blue,212}, fill opacity=0.88] (4.16,7.31) circle (0.086);
\fill[fill={rgb,255:red,58;green,154;blue,216}, fill opacity=0.94] (9.22,9.52) circle (0.094);
\fill[fill={rgb,255:red,59;green,154;blue,213}, fill opacity=0.87] (7.97,10.88) circle (0.085);
\fill[fill={rgb,255:red,58;green,154;blue,216}, fill opacity=0.93] (9.47,9.77) circle (0.093);
\fill[fill={rgb,255:red,58;green,153;blue,216}, fill opacity=0.94] (7.16,8.41) circle (0.094);
\fill[fill={rgb,255:red,58;green,153;blue,216}, fill opacity=0.91] (9.26,8.67) circle (0.090);
\fill[fill={rgb,255:red,58;green,153;blue,216}, fill opacity=0.95] (8.36,8.85) circle (0.095);
\fill[fill={rgb,255:red,61;green,153;blue,213}, fill opacity=0.88] (4.44,7.46) circle (0.087);
\fill[fill={rgb,255:red,58;green,154;blue,216}, fill opacity=0.92] (7.79,9.76) circle (0.092);
\fill[fill={rgb,255:red,58;green,154;blue,215}, fill opacity=0.89] (8.48,10.68) circle (0.088);
\fill[fill={rgb,255:red,58;green,154;blue,216}, fill opacity=0.93] (7.79,9.66) circle (0.092);
\fill[fill={rgb,255:red,54;green,164;blue,196}, fill opacity=0.79] (11.55,11.95) circle (0.075);
\fill[fill={rgb,255:red,43;green,191;blue,142}, fill opacity=0.92] (13.97,10.81) circle (0.091);
\fill[fill={rgb,255:red,58;green,154;blue,216}, fill opacity=0.92] (9.44,9.05) circle (0.091);
\fill[fill={rgb,255:red,62;green,154;blue,210}, fill opacity=0.84] (6.40,10.66) circle (0.082);
\fill[fill={rgb,255:red,58;green,155;blue,214}, fill opacity=0.87] (10.78,10.29) circle (0.086);
\fill[fill={rgb,255:red,58;green,153;blue,216}, fill opacity=0.93] (7.24,9.18) circle (0.093);
\fill[fill={rgb,255:red,58;green,154;blue,216}, fill opacity=0.95] (8.70,9.36) circle (0.095);
\fill[fill={rgb,255:red,58;green,153;blue,216}, fill opacity=0.91] (5.82,8.24) circle (0.091);
\fill[fill={rgb,255:red,57;green,156;blue,210}, fill opacity=0.85] (11.05,11.11) circle (0.082);
\fill[fill={rgb,255:red,58;green,153;blue,216}, fill opacity=0.92] (8.57,8.27) circle (0.091);
\fill[fill={rgb,255:red,58;green,153;blue,216}, fill opacity=0.93] (7.82,8.00) circle (0.093);
\fill[fill={rgb,255:red,59;green,153;blue,215}, fill opacity=0.91] (5.24,7.46) circle (0.090);
\fill[fill={rgb,255:red,58;green,153;blue,216}, fill opacity=0.95] (8.16,8.54) circle (0.094);
\fill[fill={rgb,255:red,76;green,151;blue,199}, fill opacity=0.81] (9.52,7.81) circle (0.078);
\fill[fill={rgb,255:red,56;green,158;blue,206}, fill opacity=0.82] (11.09,11.66) circle (0.079);
\fill[fill={rgb,255:red,46;green,184;blue,156}, fill opacity=0.81] (13.64,13.97) circle (0.077);
\fill[fill={rgb,255:red,48;green,180;blue,164}, fill opacity=0.78] (12.61,13.16) circle (0.074);
\fill[fill={rgb,255:red,57;green,155;blue,213}, fill opacity=0.86] (10.05,11.44) circle (0.084);
\fill[fill={rgb,255:red,62;green,153;blue,212}, fill opacity=0.88] (8.56,7.70) circle (0.086);
\fill[fill={rgb,255:red,66;green,152;blue,209}, fill opacity=0.86] (3.77,7.93) circle (0.083);
\fill[fill={rgb,255:red,58;green,153;blue,216}, fill opacity=0.95] (8.43,9.09) circle (0.095);
\fill[fill={rgb,255:red,58;green,153;blue,216}, fill opacity=0.92] (9.07,8.71) circle (0.092);
\fill[fill={rgb,255:red,61;green,153;blue,214}, fill opacity=0.89] (5.28,6.63) circle (0.088);
\fill[fill={rgb,255:red,58;green,153;blue,216}, fill opacity=0.94] (8.67,8.65) circle (0.093);
\fill[fill={rgb,255:red,59;green,154;blue,212}, fill opacity=0.86] (10.56,9.37) circle (0.084);
\fill[fill={rgb,255:red,61;green,153;blue,213}, fill opacity=0.89] (4.62,6.63) circle (0.087);
\fill[fill={rgb,255:red,59;green,153;blue,215}, fill opacity=0.91] (5.35,7.33) circle (0.090);
\fill[fill={rgb,255:red,58;green,154;blue,216}, fill opacity=0.91] (10.23,10.26) circle (0.090);
\fill[fill={rgb,255:red,58;green,154;blue,216}, fill opacity=0.92] (9.67,10.27) circle (0.091);
\fill[fill={rgb,255:red,58;green,154;blue,216}, fill opacity=0.94] (9.31,9.67) circle (0.093);
\fill[fill={rgb,255:red,49;green,177;blue,170}, fill opacity=0.76] (12.43,11.20) circle (0.072);
\fill[fill={rgb,255:red,70;green,151;blue,205}, fill opacity=0.85] (3.78,5.46) circle (0.083);
\fill[fill={rgb,255:red,57;green,157;blue,209}, fill opacity=0.84] (10.88,11.49) circle (0.081);
\fill[fill={rgb,255:red,79;green,150;blue,198}, fill opacity=0.83] (4.91,5.31) circle (0.080);
\fill[fill={rgb,255:red,58;green,154;blue,216}, fill opacity=0.94] (9.02,9.08) circle (0.094);
\fill[fill={rgb,255:red,48;green,180;blue,163}, fill opacity=0.78] (12.67,13.17) circle (0.074);
\fill[fill={rgb,255:red,58;green,154;blue,216}, fill opacity=0.95] (8.12,9.29) circle (0.094);
\fill[fill={rgb,255:red,58;green,154;blue,215}, fill opacity=0.89] (9.97,11.00) circle (0.087);
\fill[fill={rgb,255:red,58;green,154;blue,216}, fill opacity=0.92] (9.69,10.11) circle (0.092);
\fill[fill={rgb,255:red,55;green,166;blue,190}, fill opacity=0.77] (11.89,10.35) circle (0.072);
\fill[fill={rgb,255:red,63;green,153;blue,211}, fill opacity=0.88] (4.59,6.20) circle (0.086);
\fill[fill={rgb,255:red,58;green,153;blue,216}, fill opacity=0.94] (7.28,8.59) circle (0.094);
\fill[fill={rgb,255:red,58;green,154;blue,216}, fill opacity=0.93] (9.61,9.63) circle (0.092);
\fill[fill={rgb,255:red,58;green,153;blue,216}, fill opacity=0.93] (6.47,7.73) circle (0.093);
\fill[fill={rgb,255:red,58;green,154;blue,216}, fill opacity=0.94] (8.51,9.49) circle (0.094);
\fill[fill={rgb,255:red,61;green,153;blue,213}, fill opacity=0.88] (4.39,7.46) circle (0.087);
\fill[fill={rgb,255:red,58;green,153;blue,216}, fill opacity=0.94] (7.32,7.99) circle (0.094);
\fill[fill={rgb,255:red,53;green,168;blue,187}, fill opacity=0.76] (11.48,12.61) circle (0.072);
\fill[fill={rgb,255:red,59;green,153;blue,215}, fill opacity=0.91] (5.44,7.21) circle (0.090);
\fill[fill={rgb,255:red,58;green,153;blue,216}, fill opacity=0.91] (6.13,8.76) circle (0.090);
\fill[fill={rgb,255:red,60;green,153;blue,214}, fill opacity=0.89] (5.24,8.32) circle (0.088);
\fill[fill={rgb,255:red,58;green,155;blue,212}, fill opacity=0.86] (8.55,11.27) circle (0.084);
\fill[fill={rgb,255:red,55;green,163;blue,197}, fill opacity=0.79] (11.14,12.28) circle (0.075);
\fill[fill={rgb,255:red,58;green,153;blue,216}, fill opacity=0.91] (6.22,8.93) circle (0.090);
\fill[fill={rgb,255:red,76;green,151;blue,200}, fill opacity=0.83] (1.97,6.75) circle (0.080);
\fill[fill={rgb,255:red,52;green,175;blue,170}, fill opacity=0.75] (12.25,10.04) circle (0.070);
\fill[fill={rgb,255:red,60;green,153;blue,214}, fill opacity=0.89] (6.34,6.80) circle (0.088);
\fill[fill={rgb,255:red,58;green,154;blue,216}, fill opacity=0.92] (7.67,9.88) circle (0.091);
\fill[fill={rgb,255:red,43;green,191;blue,142}, fill opacity=0.91] (13.89,10.67) circle (0.091);
\fill[fill={rgb,255:red,58;green,153;blue,216}, fill opacity=0.94] (7.33,8.72) circle (0.094);
\fill[fill={rgb,255:red,58;green,154;blue,216}, fill opacity=0.94] (8.76,9.69) circle (0.094);
\fill[fill={rgb,255:red,59;green,159;blue,201}, fill opacity=0.80] (11.35,9.63) circle (0.076);
\fill[fill={rgb,255:red,58;green,153;blue,216}, fill opacity=0.92] (9.09,8.73) circle (0.092);
\fill[fill={rgb,255:red,49;green,178;blue,168}, fill opacity=0.77] (12.46,12.76) circle (0.073);
\fill[fill={rgb,255:red,58;green,154;blue,216}, fill opacity=0.93] (8.04,9.63) circle (0.093);
\fill[fill={rgb,255:red,57;green,157;blue,209}, fill opacity=0.84] (11.25,10.56) circle (0.081);
\fill[fill={rgb,255:red,57;green,158;blue,206}, fill opacity=0.82] (11.40,10.40) circle (0.079);
\fill[fill={rgb,255:red,58;green,153;blue,216}, fill opacity=0.94] (7.97,8.18) circle (0.093);
\fill[fill={rgb,255:red,58;green,153;blue,216}, fill opacity=0.93] (6.90,8.77) circle (0.093);
\fill[fill={rgb,255:red,87;green,149;blue,191}, fill opacity=0.81] (3.13,4.16) circle (0.078);
\fill[fill={rgb,255:red,58;green,153;blue,216}, fill opacity=0.91] (5.63,7.40) circle (0.091);
\fill[fill={rgb,255:red,60;green,153;blue,214}, fill opacity=0.88] (8.92,8.05) circle (0.087);
\fill[fill={rgb,255:red,58;green,154;blue,216}, fill opacity=0.90] (8.98,10.67) circle (0.089);
\fill[fill={rgb,255:red,58;green,154;blue,216}, fill opacity=0.94] (8.28,9.41) circle (0.094);
\fill[fill={rgb,255:red,58;green,154;blue,215}, fill opacity=0.89] (8.28,10.67) circle (0.088);
\fill[fill={rgb,255:red,92;green,149;blue,185}, fill opacity=0.79] (9.35,7.46) circle (0.075);
\fill[fill={rgb,255:red,67;green,158;blue,194}, fill opacity=0.78] (11.12,8.94) circle (0.074);
\fill[fill={rgb,255:red,50;green,179;blue,163}, fill opacity=0.77] (17.57,15.03) circle (0.073);
\fill[fill={rgb,255:red,62;green,153;blue,212}, fill opacity=0.88] (4.30,6.65) circle (0.086);
\fill[fill={rgb,255:red,54;green,166;blue,191}, fill opacity=0.77] (11.92,10.69) circle (0.073);
\fill[fill={rgb,255:red,60;green,153;blue,214}, fill opacity=0.90] (4.81,7.13) circle (0.088);
\fill[fill={rgb,255:red,58;green,153;blue,216}, fill opacity=0.95] (7.98,9.05) circle (0.095);
\fill[fill={rgb,255:red,58;green,154;blue,216}, fill opacity=0.92] (7.54,9.75) circle (0.091);
\fill[fill={rgb,255:red,58;green,153;blue,216}, fill opacity=0.93] (6.41,8.29) circle (0.093);
\fill[fill={rgb,255:red,57;green,157;blue,209}, fill opacity=0.84] (10.76,11.60) circle (0.081);
\fill[fill={rgb,255:red,58;green,155;blue,213}, fill opacity=0.87] (9.43,11.38) circle (0.084);
\fill[fill={rgb,255:red,54;green,164;blue,195}, fill opacity=0.79] (11.76,11.52) circle (0.075);
\fill[fill={rgb,255:red,59;green,153;blue,215}, fill opacity=0.90] (8.83,8.14) circle (0.088);
\fill[fill={rgb,255:red,73;green,151;blue,203}, fill opacity=0.84] (2.43,7.03) circle (0.081);
\fill[fill={rgb,255:red,58;green,154;blue,216}, fill opacity=0.90] (8.98,10.69) circle (0.089);
\fill[fill={rgb,255:red,152;green,138;blue,135}, fill opacity=0.74] (8.45,6.36) circle (0.069);
\fill[fill={rgb,255:red,58;green,153;blue,216}, fill opacity=0.95] (7.42,8.65) circle (0.095);
\fill[fill={rgb,255:red,58;green,153;blue,216}, fill opacity=0.92] (8.44,8.16) circle (0.091);
\fill[fill={rgb,255:red,58;green,154;blue,216}, fill opacity=0.94] (8.95,9.56) circle (0.094);
\fill[fill={rgb,255:red,61;green,153;blue,214}, fill opacity=0.89] (4.55,7.16) circle (0.087);
\fill[fill={rgb,255:red,59;green,154;blue,214}, fill opacity=0.88] (7.27,10.42) circle (0.086);
\fill[fill={rgb,255:red,58;green,154;blue,216}, fill opacity=0.92] (9.87,9.94) circle (0.091);
\fill[fill={rgb,255:red,48;green,179;blue,164}, fill opacity=0.78] (12.37,13.94) circle (0.073);
\fill[fill={rgb,255:red,60;green,153;blue,214}, fill opacity=0.89] (4.86,8.06) circle (0.087);
\fill[fill={rgb,255:red,58;green,153;blue,216}, fill opacity=0.93] (6.14,7.98) circle (0.092);
\fill[fill={rgb,255:red,58;green,153;blue,216}, fill opacity=0.95] (7.62,8.57) circle (0.095);
\fill[fill={rgb,255:red,58;green,154;blue,216}, fill opacity=0.94] (8.64,9.62) circle (0.094);
\fill[fill={rgb,255:red,65;green,152;blue,210}, fill opacity=0.86] (3.51,6.71) circle (0.084);
\fill[fill={rgb,255:red,58;green,153;blue,216}, fill opacity=0.94] (8.57,8.83) circle (0.094);
\fill[fill={rgb,255:red,58;green,156;blue,210}, fill opacity=0.84] (9.01,11.69) circle (0.081);
\fill[fill={rgb,255:red,58;green,153;blue,216}, fill opacity=0.94] (6.62,8.26) circle (0.093);
\fill[fill={rgb,255:red,58;green,154;blue,215}, fill opacity=0.90] (9.67,8.96) circle (0.089);
\fill[fill={rgb,255:red,58;green,154;blue,216}, fill opacity=0.91] (7.80,10.01) circle (0.090);
\fill[fill={rgb,255:red,58;green,154;blue,216}, fill opacity=0.94] (9.32,9.76) circle (0.093);
\fill[fill={rgb,255:red,58;green,154;blue,216}, fill opacity=0.94] (8.65,9.70) circle (0.094);
\fill[fill={rgb,255:red,47;green,181;blue,161}, fill opacity=0.79] (12.72,12.07) circle (0.075);
\fill[fill={rgb,255:red,57;green,155;blue,212}, fill opacity=0.86] (10.02,11.54) circle (0.083);
\fill[fill={rgb,255:red,58;green,154;blue,214}, fill opacity=0.88] (9.95,11.13) circle (0.086);
\fill[fill={rgb,255:red,66;green,152;blue,209}, fill opacity=0.87] (3.95,5.99) circle (0.084);
\fill[fill={rgb,255:red,58;green,154;blue,216}, fill opacity=0.91] (9.69,10.50) circle (0.091);
\fill[fill={rgb,255:red,58;green,153;blue,216}, fill opacity=0.93] (8.64,8.39) circle (0.092);
\fill[fill={rgb,255:red,62;green,153;blue,211}, fill opacity=0.85] (5.69,9.91) circle (0.083);
\fill[fill={rgb,255:red,58;green,154;blue,216}, fill opacity=0.94] (8.93,9.52) circle (0.094);
\fill[fill={rgb,255:red,57;green,157;blue,210}, fill opacity=0.84] (10.87,11.45) circle (0.082);
\fill[fill={rgb,255:red,51;green,173;blue,177}, fill opacity=0.74] (12.22,10.70) circle (0.069);
\fill[fill={rgb,255:red,66;green,152;blue,210}, fill opacity=0.87] (4.97,5.99) circle (0.085);
\fill[fill={rgb,255:red,64;green,153;blue,210}, fill opacity=0.86] (4.20,8.32) circle (0.084);
\fill[fill={rgb,255:red,58;green,153;blue,216}, fill opacity=0.93] (7.05,9.13) circle (0.092);
\fill[fill={rgb,255:red,73;green,151;blue,203}, fill opacity=0.84] (2.68,5.26) circle (0.081);
\fill[fill={rgb,255:red,60;green,154;blue,212}, fill opacity=0.86] (10.50,9.24) circle (0.083);
\fill[fill={rgb,255:red,59;green,154;blue,214}, fill opacity=0.88] (7.16,10.29) circle (0.086);
\fill[fill={rgb,255:red,61;green,153;blue,213}, fill opacity=0.89] (5.89,6.58) circle (0.087);
\fill[fill={rgb,255:red,58;green,155;blue,214}, fill opacity=0.88] (10.78,10.43) circle (0.086);
\fill[fill={rgb,255:red,58;green,153;blue,216}, fill opacity=0.95] (7.38,8.45) circle (0.095);
\fill[fill={rgb,255:red,58;green,153;blue,216}, fill opacity=0.94] (7.32,9.01) circle (0.094);
\fill[fill={rgb,255:red,58;green,154;blue,216}, fill opacity=0.94] (8.96,9.73) circle (0.094);
\fill[fill={rgb,255:red,58;green,153;blue,216}, fill opacity=0.95] (7.55,8.72) circle (0.095);
\fill[fill={rgb,255:red,62;green,153;blue,212}, fill opacity=0.86] (5.65,9.71) circle (0.084);
\fill[fill={rgb,255:red,58;green,153;blue,216}, fill opacity=0.93] (6.39,8.09) circle (0.093);
\fill[fill={rgb,255:red,48;green,180;blue,164}, fill opacity=0.78] (12.64,12.75) circle (0.074);
\fill[fill={rgb,255:red,43;green,191;blue,142}, fill opacity=0.91] (14.36,11.81) circle (0.091);
\fill[fill={rgb,255:red,59;green,153;blue,215}, fill opacity=0.90] (5.29,8.10) circle (0.089);
\fill[fill={rgb,255:red,65;green,153;blue,209}, fill opacity=0.85] (9.45,8.10) circle (0.083);
\fill[fill={rgb,255:red,58;green,154;blue,216}, fill opacity=0.93] (9.51,9.36) circle (0.092);
\fill[fill={rgb,255:red,65;green,153;blue,208}, fill opacity=0.84] (5.06,9.76) circle (0.082);
\fill[fill={rgb,255:red,58;green,153;blue,216}, fill opacity=0.94] (7.90,8.14) circle (0.093);
\fill[fill={rgb,255:red,56;green,160;blue,204}, fill opacity=0.82] (11.43,11.40) circle (0.078);
\fill[fill={rgb,255:red,58;green,153;blue,216}, fill opacity=0.91] (9.22,8.60) circle (0.090);
\fill[fill={rgb,255:red,58;green,153;blue,216}, fill opacity=0.91] (8.39,8.04) circle (0.090);
\fill[fill={rgb,255:red,62;green,153;blue,212}, fill opacity=0.88] (4.37,7.89) circle (0.086);
\fill[fill={rgb,255:red,60;green,153;blue,214}, fill opacity=0.90] (5.28,6.71) circle (0.088);
\fill[fill={rgb,255:red,58;green,153;blue,216}, fill opacity=0.94] (6.71,8.08) circle (0.093);
\fill[fill={rgb,255:red,58;green,154;blue,216}, fill opacity=0.94] (9.37,9.66) circle (0.093);
\fill[fill={rgb,255:red,48;green,181;blue,161}, fill opacity=0.78] (12.55,10.54) circle (0.074);
\fill[fill={rgb,255:red,58;green,153;blue,216}, fill opacity=0.92] (5.95,7.50) circle (0.091);
\fill[fill={rgb,255:red,58;green,153;blue,216}, fill opacity=0.91] (6.93,9.57) circle (0.090);
\fill[fill={rgb,255:red,44;green,189;blue,146}, fill opacity=0.86] (13.42,11.90) circle (0.083);
\fill[fill={rgb,255:red,58;green,153;blue,216}, fill opacity=0.94] (7.38,9.16) circle (0.093);
\fill[fill={rgb,255:red,58;green,154;blue,215}, fill opacity=0.89] (8.59,10.73) circle (0.088);
\fill[fill={rgb,255:red,59;green,153;blue,215}, fill opacity=0.91] (8.16,7.81) circle (0.090);
\fill[fill={rgb,255:red,57;green,156;blue,211}, fill opacity=0.85] (10.28,11.62) circle (0.082);
\fill[fill={rgb,255:red,80;green,150;blue,197}, fill opacity=0.82] (7.82,6.77) circle (0.079);
\fill[fill={rgb,255:red,63;green,153;blue,212}, fill opacity=0.87] (3.95,7.06) circle (0.086);
\fill[fill={rgb,255:red,63;green,153;blue,212}, fill opacity=0.88] (4.50,6.28) circle (0.086);
\fill[fill={rgb,255:red,58;green,154;blue,216}, fill opacity=0.90] (9.54,8.85) circle (0.089);
\fill[fill={rgb,255:red,83;green,149;blue,194}, fill opacity=0.82] (1.02,4.15) circle (0.079);
\fill[fill={rgb,255:red,69;green,152;blue,207}, fill opacity=0.86] (3.96,5.60) circle (0.083);
\fill[fill={rgb,255:red,65;green,152;blue,209}, fill opacity=0.86] (8.54,7.53) circle (0.084);
\fill[fill={rgb,255:red,58;green,154;blue,215}, fill opacity=0.90] (8.63,10.71) circle (0.088);
\fill[fill={rgb,255:red,54;green,164;blue,195}, fill opacity=0.78] (11.04,12.46) circle (0.074);
\fill[fill={rgb,255:red,58;green,153;blue,216}, fill opacity=0.92] (7.71,7.70) circle (0.091);
\fill[fill={rgb,255:red,58;green,154;blue,216}, fill opacity=0.92] (9.76,9.78) circle (0.092);
\fill[fill={rgb,255:red,58;green,153;blue,216}, fill opacity=0.95] (8.52,8.85) circle (0.094);
\fill[fill={rgb,255:red,58;green,154;blue,215}, fill opacity=0.89] (10.59,10.54) circle (0.087);
\fill[fill={rgb,255:red,56;green,163;blue,195}, fill opacity=0.78] (11.70,10.13) circle (0.074);
\fill[fill={rgb,255:red,58;green,154;blue,216}, fill opacity=0.91] (7.74,9.98) circle (0.090);
\fill[fill={rgb,255:red,49;green,176;blue,171}, fill opacity=0.76] (12.41,12.14) circle (0.071);
\fill[fill={rgb,255:red,58;green,154;blue,214}, fill opacity=0.88] (10.46,9.66) circle (0.086);
\fill[fill={rgb,255:red,59;green,153;blue,216}, fill opacity=0.91] (5.96,7.12) circle (0.090);
\fill[fill={rgb,255:red,68;green,152;blue,207}, fill opacity=0.86] (4.28,5.67) circle (0.083);
\fill[fill={rgb,255:red,58;green,153;blue,216}, fill opacity=0.94] (7.12,8.26) circle (0.094);
\fill[fill={rgb,255:red,58;green,154;blue,214}, fill opacity=0.88] (9.90,11.11) circle (0.087);
\fill[fill={rgb,255:red,58;green,154;blue,216}, fill opacity=0.92] (8.20,9.96) circle (0.092);
\fill[fill={rgb,255:red,58;green,153;blue,216}, fill opacity=0.92] (8.82,8.41) circle (0.091);
\fill[fill={rgb,255:red,58;green,154;blue,216}, fill opacity=0.95] (8.35,9.42) circle (0.094);
\fill[fill={rgb,255:red,61;green,153;blue,213}, fill opacity=0.89] (5.88,6.61) circle (0.088);
\fill[fill={rgb,255:red,66;green,153;blue,209}, fill opacity=0.85] (9.14,7.88) circle (0.083);
\fill[fill={rgb,255:red,60;green,153;blue,215}, fill opacity=0.90] (6.84,7.06) circle (0.089);
\fill[fill={rgb,255:red,45;green,185;blue,153}, fill opacity=0.82] (13.23,12.62) circle (0.079);
\fill[fill={rgb,255:red,75;green,151;blue,201}, fill opacity=0.84] (5.80,5.83) circle (0.081);
\fill[fill={rgb,255:red,56;green,160;blue,203}, fill opacity=0.81] (11.57,10.69) circle (0.077);
\fill[fill={rgb,255:red,56;green,158;blue,206}, fill opacity=0.82] (11.08,11.67) circle (0.079);
\fill[fill={rgb,255:red,58;green,154;blue,214}, fill opacity=0.87] (10.77,10.16) circle (0.085);
\fill[fill={rgb,255:red,58;green,153;blue,216}, fill opacity=0.92] (6.85,7.45) circle (0.092);
\fill[fill={rgb,255:red,58;green,153;blue,216}, fill opacity=0.92] (7.40,7.57) circle (0.091);
\fill[fill={rgb,255:red,50;green,181;blue,159}, fill opacity=0.78] (12.52,10.05) circle (0.074);
\fill[fill={rgb,255:red,72;green,151;blue,204}, fill opacity=0.84] (3.00,5.34) circle (0.082);
\fill[fill={rgb,255:red,58;green,153;blue,216}, fill opacity=0.95] (7.42,8.69) circle (0.095);
\fill[fill={rgb,255:red,58;green,153;blue,216}, fill opacity=0.95] (7.61,8.91) circle (0.094);
\fill[fill={rgb,255:red,79;green,150;blue,198}, fill opacity=0.83] (3.26,4.67) circle (0.080);
\fill[fill={rgb,255:red,45;green,187;blue,150}, fill opacity=0.83] (13.04,11.38) circle (0.080);
\fill[fill={rgb,255:red,54;green,166;blue,192}, fill opacity=0.78] (11.45,12.36) circle (0.073);
\fill[fill={rgb,255:red,60;green,153;blue,214}, fill opacity=0.90] (6.30,6.83) circle (0.088);
\fill[fill={rgb,255:red,51;green,172;blue,179}, fill opacity=0.74] (12.21,11.58) circle (0.069);
\fill[fill={rgb,255:red,150;green,138;blue,138}, fill opacity=0.74] (5.70,4.24) circle (0.069);
\fill[fill={rgb,255:red,59;green,153;blue,215}, fill opacity=0.91] (7.19,7.31) circle (0.090);
\fill[fill={rgb,255:red,64;green,153;blue,211}, fill opacity=0.87] (5.74,6.31) circle (0.086);
\fill[fill={rgb,255:red,61;green,153;blue,213}, fill opacity=0.87] (9.72,8.53) circle (0.085);
\fill[fill={rgb,255:red,57;green,159;blue,204}, fill opacity=0.82] (10.10,12.18) circle (0.078);
\fill[fill={rgb,255:red,58;green,153;blue,216}, fill opacity=0.92] (5.99,7.28) circle (0.091);
\fill[fill={rgb,255:red,62;green,163;blue,188}, fill opacity=0.76] (11.58,9.35) circle (0.071);
\fill[fill={rgb,255:red,58;green,154;blue,216}, fill opacity=0.94] (9.22,9.92) circle (0.093);
\fill[fill={rgb,255:red,58;green,154;blue,216}, fill opacity=0.95] (8.38,9.27) circle (0.095);
\fill[fill={rgb,255:red,56;green,158;blue,207}, fill opacity=0.83] (11.30,11.29) circle (0.080);
\fill[fill={rgb,255:red,58;green,154;blue,215}, fill opacity=0.89] (7.49,10.29) circle (0.088);
\fill[fill={rgb,255:red,71;green,151;blue,204}, fill opacity=0.85] (2.89,5.47) circle (0.082);
\fill[fill={rgb,255:red,58;green,154;blue,215}, fill opacity=0.90] (10.29,10.61) circle (0.088);
\fill[fill={rgb,255:red,58;green,154;blue,216}, fill opacity=0.90] (10.21,9.76) circle (0.089);
\fill[fill={rgb,255:red,58;green,153;blue,216}, fill opacity=0.92] (8.83,8.46) circle (0.091);
\fill[fill={rgb,255:red,58;green,153;blue,216}, fill opacity=0.93] (6.43,7.93) circle (0.093);
\fill[fill={rgb,255:red,58;green,153;blue,216}, fill opacity=0.95] (7.48,8.55) circle (0.095);
\fill[fill={rgb,255:red,58;green,154;blue,214}, fill opacity=0.87] (10.66,9.83) circle (0.085);
\fill[fill={rgb,255:red,58;green,154;blue,216}, fill opacity=0.94] (9.06,9.20) circle (0.094);
\fill[fill={rgb,255:red,58;green,153;blue,216}, fill opacity=0.95] (7.43,8.73) circle (0.094);
\fill[fill={rgb,255:red,53;green,167;blue,188}, fill opacity=0.77] (11.90,11.86) circle (0.072);
\fill[fill={rgb,255:red,51;green,172;blue,179}, fill opacity=0.74] (12.22,11.28) circle (0.069);
\fill[fill={rgb,255:red,78;green,152;blue,195}, fill opacity=0.80] (10.05,8.09) circle (0.076);
\fill[fill={rgb,255:red,58;green,153;blue,216}, fill opacity=0.94] (7.75,9.13) circle (0.094);
\fill[fill={rgb,255:red,58;green,154;blue,216}, fill opacity=0.95] (8.77,9.15) circle (0.094);
\fill[fill={rgb,255:red,58;green,154;blue,216}, fill opacity=0.91] (9.68,9.12) circle (0.090);
\fill[fill={rgb,255:red,58;green,154;blue,216}, fill opacity=0.94] (9.15,9.26) circle (0.094);
\fill[fill={rgb,255:red,58;green,153;blue,216}, fill opacity=0.91] (8.49,8.04) circle (0.090);
\fill[fill={rgb,255:red,58;green,153;blue,216}, fill opacity=0.94] (7.35,8.84) circle (0.094);
\fill[fill={rgb,255:red,58;green,154;blue,216}, fill opacity=0.95] (8.73,9.10) circle (0.094);
\fill[fill={rgb,255:red,70;green,152;blue,206}, fill opacity=0.85] (7.44,6.79) circle (0.082);
\fill[fill={rgb,255:red,58;green,153;blue,216}, fill opacity=0.93] (6.66,8.49) circle (0.093);
\fill[fill={rgb,255:red,58;green,153;blue,216}, fill opacity=0.93] (7.08,8.94) circle (0.093);
\fill[fill={rgb,255:red,43;green,191;blue,141}, fill opacity=0.92] (15.92,9.16) circle (0.091);
\fill[fill={rgb,255:red,43;green,191;blue,142}, fill opacity=0.91] (15.11,12.06) circle (0.090);
\fill[fill={rgb,255:red,47;green,188;blue,144}, fill opacity=0.84] (13.23,9.48) circle (0.081);
\fill[fill={rgb,255:red,43;green,190;blue,143}, fill opacity=0.88] (16.17,12.57) circle (0.086);
\fill[fill={rgb,255:red,43;green,190;blue,144}, fill opacity=0.88] (14.38,12.52) circle (0.086);
\fill[fill={rgb,255:red,70;green,182;blue,131}, fill opacity=0.82] (18.61,8.70) circle (0.078);
\fill[fill={rgb,255:red,135;green,155;blue,114}, fill opacity=0.71] (12.40,7.91) circle (0.065);
\fill[fill={rgb,255:red,47;green,187;blue,146}, fill opacity=0.83] (13.03,9.80) circle (0.080);
\fill[fill={rgb,255:red,82;green,178;blue,126}, fill opacity=0.80] (16.28,7.68) circle (0.077);
\fill[fill={rgb,255:red,45;green,186;blue,151}, fill opacity=0.83] (13.28,12.42) circle (0.080);
\fill[fill={rgb,255:red,45;green,189;blue,145}, fill opacity=0.85] (17.25,12.65) circle (0.083);
\fill[fill={rgb,255:red,46;green,184;blue,156}, fill opacity=0.81] (13.20,13.23) circle (0.077);
\fill[fill={rgb,255:red,46;green,190;blue,141}, fill opacity=0.87] (14.13,8.94) circle (0.085);
\fill[fill={rgb,255:red,44;green,191;blue,141}, fill opacity=0.91] (16.84,9.73) circle (0.090);
\fill[fill={rgb,255:red,43;green,191;blue,142}, fill opacity=0.94] (15.02,11.33) circle (0.094);
\fill[fill={rgb,255:red,44;green,188;blue,149}, fill opacity=0.84] (13.13,11.55) circle (0.081);
\fill[fill={rgb,255:red,43;green,191;blue,142}, fill opacity=0.92] (16.39,11.13) circle (0.091);
\fill[fill={rgb,255:red,43;green,190;blue,143}, fill opacity=0.88] (15.99,12.56) circle (0.086);
\fill[fill={rgb,255:red,43;green,191;blue,142}, fill opacity=0.94] (14.49,10.94) circle (0.094);
\fill[fill={rgb,255:red,43;green,191;blue,142}, fill opacity=0.89] (13.67,10.23) circle (0.088);
\fill[fill={rgb,255:red,78;green,179;blue,128}, fill opacity=0.81] (18.26,8.22) circle (0.077);
\fill[fill={rgb,255:red,43;green,191;blue,142}, fill opacity=0.94] (15.86,10.73) circle (0.094);
\fill[fill={rgb,255:red,124;green,163;blue,109}, fill opacity=0.76] (17.15,7.23) circle (0.071);
\fill[fill={rgb,255:red,46;green,190;blue,140}, fill opacity=0.89] (15.09,8.68) circle (0.087);
\fill[fill={rgb,255:red,188;green,141;blue,84}, fill opacity=0.79] (17.91,6.29) circle (0.075);
\fill[fill={rgb,255:red,43;green,191;blue,143}, fill opacity=0.89] (14.71,12.46) circle (0.088);
\fill[fill={rgb,255:red,43;green,191;blue,141}, fill opacity=0.92] (16.16,9.48) circle (0.092);
\fill[fill={rgb,255:red,61;green,156;blue,206}, fill opacity=0.82] (10.94,9.28) circle (0.079);
\fill[fill={rgb,255:red,43;green,191;blue,142}, fill opacity=0.93] (16.24,10.83) circle (0.092);
\fill[fill={rgb,255:red,45;green,188;blue,148}, fill opacity=0.84] (16.65,13.32) circle (0.081);
\fill[fill={rgb,255:red,44;green,189;blue,145}, fill opacity=0.86] (14.32,12.81) circle (0.084);
\fill[fill={rgb,255:red,43;green,191;blue,142}, fill opacity=0.91] (16.00,11.80) circle (0.090);
\fill[fill={rgb,255:red,43;green,191;blue,141}, fill opacity=0.93] (16.05,9.56) circle (0.092);
\fill[fill={rgb,255:red,46;green,185;blue,153}, fill opacity=0.81] (16.68,13.96) circle (0.078);
\fill[fill={rgb,255:red,43;green,191;blue,142}, fill opacity=0.95] (15.55,10.22) circle (0.094);
\fill[fill={rgb,255:red,44;green,190;blue,143}, fill opacity=0.88] (16.83,12.14) circle (0.086);
\fill[fill={rgb,255:red,44;green,190;blue,143}, fill opacity=0.87] (17.29,11.86) circle (0.085);
\fill[fill={rgb,255:red,47;green,190;blue,140}, fill opacity=0.88] (16.42,8.68) circle (0.087);
\fill[fill={rgb,255:red,43;green,191;blue,142}, fill opacity=0.91] (16.65,10.93) circle (0.091);
\fill[fill={rgb,255:red,52;green,170;blue,182}, fill opacity=0.75] (12.14,10.95) circle (0.070);
\fill[fill={rgb,255:red,45;green,190;blue,141}, fill opacity=0.89] (17.01,9.26) circle (0.088);
\fill[fill={rgb,255:red,79;green,178;blue,128}, fill opacity=0.80] (14.34,7.97) circle (0.076);
\fill[fill={rgb,255:red,43;green,190;blue,145}, fill opacity=0.87] (14.67,12.82) circle (0.085);
\fill[fill={rgb,255:red,43;green,191;blue,142}, fill opacity=0.93] (14.16,10.96) circle (0.092);
\fill[fill={rgb,255:red,43;green,190;blue,144}, fill opacity=0.87] (13.76,12.11) circle (0.085);
\fill[fill={rgb,255:red,46;green,184;blue,156}, fill opacity=0.81] (12.83,11.60) circle (0.077);
\fill[fill={rgb,255:red,45;green,191;blue,141}, fill opacity=0.90] (16.77,9.24) circle (0.089);
\fill[fill={rgb,255:red,44;green,191;blue,141}, fill opacity=0.90] (17.16,9.95) circle (0.089);
\fill[fill={rgb,255:red,43;green,191;blue,142}, fill opacity=0.91] (13.90,10.13) circle (0.090);
\fill[fill={rgb,255:red,43;green,191;blue,142}, fill opacity=0.95] (14.86,10.34) circle (0.095);
\fill[fill={rgb,255:red,43;green,190;blue,145}, fill opacity=0.87] (13.41,11.55) circle (0.085);
\fill[fill={rgb,255:red,44;green,191;blue,142}, fill opacity=0.88] (17.49,10.91) circle (0.087);
\fill[fill={rgb,255:red,49;green,180;blue,163}, fill opacity=0.78] (15.01,16.32) circle (0.073);
\fill[fill={rgb,255:red,43;green,191;blue,142}, fill opacity=0.94] (15.59,10.86) circle (0.094);
\fill[fill={rgb,255:red,46;green,184;blue,155}, fill opacity=0.81] (13.25,13.10) circle (0.078);
\fill[fill={rgb,255:red,46;green,190;blue,140}, fill opacity=0.89] (16.81,8.95) circle (0.087);
\fill[fill={rgb,255:red,43;green,191;blue,142}, fill opacity=0.91] (13.77,10.55) circle (0.090);
\fill[fill={rgb,255:red,220;green,129;blue,71}, fill opacity=0.84] (15.73,6.32) circle (0.081);
\fill[fill={rgb,255:red,43;green,191;blue,143}, fill opacity=0.89] (13.58,10.99) circle (0.088);
\fill[fill={rgb,255:red,43;green,191;blue,142}, fill opacity=0.94] (15.85,10.49) circle (0.094);
\fill[fill={rgb,255:red,43;green,191;blue,142}, fill opacity=0.92] (14.17,11.31) circle (0.092);
\fill[fill={rgb,255:red,43;green,191;blue,142}, fill opacity=0.90] (13.63,10.58) circle (0.088);
\fill[fill={rgb,255:red,45;green,191;blue,141}, fill opacity=0.89] (14.58,8.95) circle (0.088);
\fill[fill={rgb,255:red,44;green,189;blue,145}, fill opacity=0.86] (13.89,12.49) circle (0.084);
\fill[fill={rgb,255:red,44;green,191;blue,141}, fill opacity=0.91] (15.92,8.92) circle (0.089);
\fill[fill={rgb,255:red,43;green,191;blue,142}, fill opacity=0.92] (16.57,10.81) circle (0.091);
\fill[fill={rgb,255:red,179;green,144;blue,87}, fill opacity=0.78] (16.17,6.78) circle (0.074);
\fill[fill={rgb,255:red,51;green,186;blue,144}, fill opacity=0.82] (19.08,12.25) circle (0.078);
\fill[fill={rgb,255:red,43;green,191;blue,143}, fill opacity=0.89] (14.02,12.10) circle (0.087);
\fill[fill={rgb,255:red,46;green,186;blue,151}, fill opacity=0.82] (17.06,13.71) circle (0.078);
\fill[fill={rgb,255:red,43;green,191;blue,141}, fill opacity=0.94] (15.24,9.79) circle (0.094);
\fill[fill={rgb,255:red,44;green,189;blue,147}, fill opacity=0.85] (14.45,13.12) circle (0.083);
\fill[fill={rgb,255:red,43;green,191;blue,142}, fill opacity=0.91] (13.89,11.24) circle (0.090);
\fill[fill={rgb,255:red,48;green,181;blue,161}, fill opacity=0.78] (15.05,15.50) circle (0.074);
\fill[fill={rgb,255:red,43;green,191;blue,142}, fill opacity=0.91] (15.52,12.17) circle (0.089);
\fill[fill={rgb,255:red,118;green,165;blue,112}, fill opacity=0.76] (18.84,7.60) circle (0.071);
\fill[fill={rgb,255:red,43;green,191;blue,142}, fill opacity=0.93] (16.20,10.61) circle (0.093);
\fill[fill={rgb,255:red,43;green,191;blue,142}, fill opacity=0.92] (14.08,11.41) circle (0.091);
\fill[fill={rgb,255:red,87;green,176;blue,124}, fill opacity=0.79] (18.06,7.93) circle (0.076);
\fill[fill={rgb,255:red,44;green,191;blue,141}, fill opacity=0.91] (15.95,8.94) circle (0.089);
\fill[fill={rgb,255:red,43;green,191;blue,141}, fill opacity=0.92] (15.75,9.17) circle (0.091);
\fill[fill={rgb,255:red,133;green,159;blue,106}, fill opacity=0.74] (19.93,7.65) circle (0.069);
\fill[fill={rgb,255:red,43;green,191;blue,142}, fill opacity=0.89] (16.47,12.02) circle (0.088);
\fill[fill={rgb,255:red,47;green,181;blue,161}, fill opacity=0.79] (12.86,14.02) circle (0.075);
\fill[fill={rgb,255:red,190;green,140;blue,83}, fill opacity=0.79] (16.62,6.56) circle (0.076);
\fill[fill={rgb,255:red,64;green,184;blue,133}, fill opacity=0.83] (16.34,7.96) circle (0.080);
\fill[fill={rgb,255:red,65;green,180;blue,143}, fill opacity=0.78] (12.78,8.81) circle (0.073);
\fill[fill={rgb,255:red,47;green,187;blue,146}, fill opacity=0.83] (13.05,9.72) circle (0.080);
\fill[fill={rgb,255:red,43;green,191;blue,142}, fill opacity=0.95] (14.87,10.96) circle (0.095);
\fill[fill={rgb,255:red,43;green,191;blue,141}, fill opacity=0.94] (15.71,9.92) circle (0.094);
\fill[fill={rgb,255:red,43;green,191;blue,141}, fill opacity=0.91] (15.33,9.02) circle (0.090);
\fill[fill={rgb,255:red,100;green,171;blue,119}, fill opacity=0.78] (17.55,7.57) circle (0.074);
\fill[fill={rgb,255:red,48;green,180;blue,163}, fill opacity=0.78] (15.49,15.62) circle (0.074);
\fill[fill={rgb,255:red,45;green,187;blue,149}, fill opacity=0.84] (14.57,13.51) circle (0.081);
\fill[fill={rgb,255:red,47;green,190;blue,141}, fill opacity=0.86] (18.09,10.71) circle (0.084);
\fill[fill={rgb,255:red,43;green,191;blue,142}, fill opacity=0.95] (15.39,10.38) circle (0.095);
\fill[fill={rgb,255:red,43;green,191;blue,141}, fill opacity=0.94] (15.89,10.12) circle (0.094);
\fill[fill={rgb,255:red,43;green,191;blue,141}, fill opacity=0.93] (15.86,9.70) circle (0.093);
\fill[fill={rgb,255:red,53;green,188;blue,138}, fill opacity=0.86] (16.34,8.30) circle (0.084);
\fill[fill={rgb,255:red,48;green,180;blue,162}, fill opacity=0.78] (12.35,14.46) circle (0.074);
\fill[fill={rgb,255:red,43;green,191;blue,142}, fill opacity=0.93] (16.38,10.62) circle (0.092);
\fill[fill={rgb,255:red,44;green,188;blue,147}, fill opacity=0.84] (13.08,10.73) circle (0.082);
\fill[fill={rgb,255:red,43;green,191;blue,142}, fill opacity=0.91] (16.82,11.04) circle (0.090);
\fill[fill={rgb,255:red,43;green,191;blue,142}, fill opacity=0.94] (15.20,11.13) circle (0.094);
\fill[fill={rgb,255:red,48;green,178;blue,167}, fill opacity=0.77] (12.53,12.62) circle (0.073);
\fill[fill={rgb,255:red,43;green,191;blue,142}, fill opacity=0.94] (14.64,10.20) circle (0.094);
\fill[fill={rgb,255:red,57;green,157;blue,210}, fill opacity=0.84] (11.01,11.32) circle (0.082);
\fill[fill={rgb,255:red,44;green,191;blue,141}, fill opacity=0.90] (15.60,8.87) circle (0.089);
\fill[fill={rgb,255:red,43;green,191;blue,142}, fill opacity=0.90] (15.30,12.27) circle (0.089);
\fill[fill={rgb,255:red,55;green,162;blue,199}, fill opacity=0.80] (11.61,11.50) circle (0.076);
\fill[fill={rgb,255:red,43;green,191;blue,142}, fill opacity=0.91] (14.06,11.58) circle (0.090);
\fill[fill={rgb,255:red,44;green,191;blue,141}, fill opacity=0.91] (16.94,9.88) circle (0.089);
\fill[fill={rgb,255:red,204;green,135;blue,77}, fill opacity=0.81] (13.98,6.99) circle (0.077);
\fill[fill={rgb,255:red,56;green,186;blue,137}, fill opacity=0.83] (14.00,8.49) circle (0.080);
\fill[fill={rgb,255:red,43;green,191;blue,141}, fill opacity=0.91] (16.71,9.76) circle (0.090);
\fill[fill={rgb,255:red,44;green,189;blue,144}, fill opacity=0.86] (16.99,12.57) circle (0.084);
\fill[fill={rgb,255:red,52;green,187;blue,139}, fill opacity=0.84] (13.73,8.76) circle (0.081);
\fill[fill={rgb,255:red,50;green,189;blue,139}, fill opacity=0.87] (15.22,8.36) circle (0.084);
\fill[fill={rgb,255:red,52;green,188;blue,138}, fill opacity=0.86] (17.53,8.82) circle (0.083);
\fill[fill={rgb,255:red,141;green,156;blue,104}, fill opacity=0.73] (13.75,7.56) circle (0.068);
\fill[fill={rgb,255:red,43;green,191;blue,142}, fill opacity=0.92] (14.35,11.65) circle (0.091);
\fill[fill={rgb,255:red,43;green,191;blue,142}, fill opacity=0.94] (15.07,11.26) circle (0.094);
\fill[fill={rgb,255:red,43;green,191;blue,142}, fill opacity=0.93] (14.63,11.68) circle (0.092);
\fill[fill={rgb,255:red,44;green,188;blue,147}, fill opacity=0.85] (14.33,13.10) circle (0.083);
\fill[fill={rgb,255:red,43;green,191;blue,142}, fill opacity=0.93] (15.41,11.61) circle (0.092);
\fill[fill={rgb,255:red,77;green,179;blue,128}, fill opacity=0.81] (17.08,7.84) circle (0.078);
\fill[fill={rgb,255:red,43;green,191;blue,142}, fill opacity=0.94] (16.00,10.28) circle (0.093);
\fill[fill={rgb,255:red,63;green,184;blue,135}, fill opacity=0.82] (19.22,9.70) circle (0.079);
\fill[fill={rgb,255:red,44;green,190;blue,142}, fill opacity=0.88] (13.74,9.59) circle (0.087);
\fill[fill={rgb,255:red,120;green,164;blue,111}, fill opacity=0.76] (17.27,7.29) circle (0.071);
\fill[fill={rgb,255:red,43;green,191;blue,142}, fill opacity=0.93] (14.61,9.71) circle (0.092);
\fill[fill={rgb,255:red,43;green,191;blue,142}, fill opacity=0.94] (14.47,10.10) circle (0.093);
\fill[fill={rgb,255:red,43;green,191;blue,142}, fill opacity=0.91] (13.77,10.83) circle (0.090);
\fill[fill={rgb,255:red,43;green,191;blue,142}, fill opacity=0.90] (16.68,11.39) circle (0.089);
\fill[fill={rgb,255:red,51;green,175;blue,172}, fill opacity=0.75] (10.79,14.09) circle (0.070);
\fill[fill={rgb,255:red,44;green,191;blue,141}, fill opacity=0.90] (14.95,8.94) circle (0.089);
\fill[fill={rgb,255:red,45;green,187;blue,150}, fill opacity=0.83] (12.94,10.84) circle (0.080);
\fill[fill={rgb,255:red,43;green,191;blue,141}, fill opacity=0.93] (15.99,9.67) circle (0.093);
\fill[fill={rgb,255:red,66;green,183;blue,133}, fill opacity=0.83] (15.47,7.94) circle (0.080);
\fill[fill={rgb,255:red,130;green,160;blue,109}, fill opacity=0.73] (13.37,7.74) circle (0.068);
\fill[fill={rgb,255:red,88;green,175;blue,124}, fill opacity=0.80] (16.76,7.63) circle (0.076);
\fill[fill={rgb,255:red,43;green,191;blue,142}, fill opacity=0.95] (14.97,10.63) circle (0.095);
\fill[fill={rgb,255:red,145;green,155;blue,101}, fill opacity=0.75] (17.12,7.01) circle (0.069);
\fill[fill={rgb,255:red,43;green,191;blue,143}, fill opacity=0.89] (15.06,12.45) circle (0.088);
\fill[fill={rgb,255:red,43;green,191;blue,142}, fill opacity=0.92] (14.38,9.57) circle (0.091);
\fill[fill={rgb,255:red,43;green,191;blue,142}, fill opacity=0.91] (13.85,11.10) circle (0.090);
\fill[fill={rgb,255:red,77;green,179;blue,128}, fill opacity=0.81] (15.91,7.75) circle (0.078);
\fill[fill={rgb,255:red,45;green,186;blue,151}, fill opacity=0.83] (13.79,13.28) circle (0.080);
\fill[fill={rgb,255:red,44;green,187;blue,149}, fill opacity=0.84] (13.29,12.12) circle (0.081);
\fill[fill={rgb,255:red,44;green,189;blue,145}, fill opacity=0.87] (13.42,11.62) circle (0.084);
\fill[fill={rgb,255:red,49;green,186;blue,146}, fill opacity=0.82] (18.73,12.59) circle (0.078);
\fill[fill={rgb,255:red,43;green,191;blue,143}, fill opacity=0.90] (13.77,11.41) circle (0.089);
\fill[fill={rgb,255:red,43;green,191;blue,142}, fill opacity=0.94] (15.92,10.70) circle (0.094);
\fill[fill={rgb,255:red,45;green,187;blue,151}, fill opacity=0.83] (13.21,12.20) circle (0.080);
\fill[fill={rgb,255:red,125;green,163;blue,109}, fill opacity=0.76] (17.31,7.24) circle (0.071);
\fill[fill={rgb,255:red,43;green,191;blue,142}, fill opacity=0.95] (15.52,10.67) circle (0.095);
\fill[fill={rgb,255:red,44;green,188;blue,148}, fill opacity=0.84] (14.09,13.14) circle (0.082);
\fill[fill={rgb,255:red,45;green,187;blue,150}, fill opacity=0.83] (12.96,10.93) circle (0.080);
\fill[fill={rgb,255:red,47;green,184;blue,153}, fill opacity=0.81] (16.91,14.02) circle (0.077);
\fill[fill={rgb,255:red,43;green,191;blue,142}, fill opacity=0.91] (14.71,12.13) circle (0.090);
\fill[fill={rgb,255:red,45;green,189;blue,144}, fill opacity=0.85] (17.45,12.48) circle (0.083);
\fill[fill={rgb,255:red,46;green,190;blue,140}, fill opacity=0.89] (16.46,8.82) circle (0.088);
\fill[fill={rgb,255:red,43;green,191;blue,142}, fill opacity=0.91] (15.74,12.04) circle (0.090);
\fill[fill={rgb,255:red,43;green,191;blue,141}, fill opacity=0.93] (16.35,10.06) circle (0.092);
\fill[fill={rgb,255:red,43;green,191;blue,141}, fill opacity=0.94] (15.91,10.17) circle (0.094);
\fill[fill={rgb,255:red,44;green,191;blue,141}, fill opacity=0.89] (14.13,9.24) circle (0.087);
\fill[fill={rgb,255:red,43;green,191;blue,142}, fill opacity=0.91] (16.53,11.32) circle (0.090);
\fill[fill={rgb,255:red,47;green,190;blue,140}, fill opacity=0.88] (16.35,8.63) circle (0.086);
\fill[fill={rgb,255:red,44;green,189;blue,146}, fill opacity=0.86] (13.19,10.70) circle (0.083);
\fill[fill={rgb,255:red,43;green,191;blue,142}, fill opacity=0.90] (13.83,9.87) circle (0.088);
\fill[fill={rgb,255:red,48;green,182;blue,159}, fill opacity=0.79] (16.84,14.69) circle (0.075);
\fill[fill={rgb,255:red,48;green,180;blue,163}, fill opacity=0.78] (12.50,13.87) circle (0.074);
\fill[fill={rgb,255:red,43;green,191;blue,142}, fill opacity=0.91] (14.59,11.95) circle (0.091);
\fill[fill={rgb,255:red,44;green,190;blue,143}, fill opacity=0.87] (17.21,12.17) circle (0.085);
\fill[fill={rgb,255:red,58;green,186;blue,136}, fill opacity=0.84] (18.44,9.23) circle (0.081);
\fill[fill={rgb,255:red,44;green,190;blue,142}, fill opacity=0.88] (17.24,11.77) circle (0.086);
\fill[fill={rgb,255:red,43;green,191;blue,142}, fill opacity=0.93] (14.31,10.24) circle (0.093);
\fill[fill={rgb,255:red,43;green,191;blue,142}, fill opacity=0.92] (15.97,11.71) circle (0.091);
\fill[fill={rgb,255:red,43;green,191;blue,142}, fill opacity=0.93] (15.97,11.06) circle (0.093);
\fill[fill={rgb,255:red,143;green,155;blue,104}, fill opacity=0.73] (13.39,7.63) circle (0.068);
\fill[fill={rgb,255:red,43;green,190;blue,144}, fill opacity=0.88] (15.11,12.77) circle (0.086);
\fill[fill={rgb,255:red,43;green,191;blue,142}, fill opacity=0.95] (15.29,10.86) circle (0.095);
\fill[fill={rgb,255:red,44;green,189;blue,146}, fill opacity=0.86] (13.48,11.92) circle (0.084);
\fill[fill={rgb,255:red,43;green,191;blue,141}, fill opacity=0.91] (16.08,9.13) circle (0.090);
\fill[fill={rgb,255:red,52;green,169;blue,184}, fill opacity=0.75] (12.11,11.32) circle (0.071);
\fill[fill={rgb,255:red,45;green,185;blue,153}, fill opacity=0.82] (13.05,12.02) circle (0.079);
\fill[fill={rgb,255:red,52;green,188;blue,139}, fill opacity=0.85] (18.41,9.80) circle (0.082);
\fill[fill={rgb,255:red,46;green,190;blue,141}, fill opacity=0.87] (13.89,9.15) circle (0.085);
\fill[fill={rgb,255:red,43;green,191;blue,142}, fill opacity=0.90] (17.04,10.79) circle (0.089);
\fill[fill={rgb,255:red,44;green,187;blue,149}, fill opacity=0.84] (13.31,12.22) circle (0.081);
\fill[fill={rgb,255:red,43;green,191;blue,142}, fill opacity=0.91] (15.77,11.92) circle (0.090);
\fill[fill={rgb,255:red,43;green,191;blue,141}, fill opacity=0.91] (16.68,9.75) circle (0.090);
\fill[fill={rgb,255:red,43;green,191;blue,142}, fill opacity=0.94] (14.77,9.98) circle (0.094);
\fill[fill={rgb,255:red,43;green,191;blue,141}, fill opacity=0.93] (16.12,9.74) circle (0.092);
\fill[fill={rgb,255:red,43;green,191;blue,142}, fill opacity=0.94] (14.65,10.06) circle (0.094);
\fill[fill={rgb,255:red,44;green,191;blue,142}, fill opacity=0.89] (17.35,10.84) circle (0.088);
\fill[fill={rgb,255:red,43;green,191;blue,142}, fill opacity=0.91] (16.23,11.74) circle (0.090);
\fill[fill={rgb,255:red,44;green,189;blue,147}, fill opacity=0.85] (14.58,13.19) circle (0.083);
\fill[fill={rgb,255:red,50;green,175;blue,173}, fill opacity=0.75] (11.41,13.73) circle (0.070);
\fill[fill={rgb,255:red,43;green,191;blue,142}, fill opacity=0.90] (13.95,11.65) circle (0.089);
\fill[fill={rgb,255:red,48;green,181;blue,161}, fill opacity=0.79] (15.13,15.37) circle (0.074);
\fill[fill={rgb,255:red,208;green,134;blue,75}, fill opacity=0.82] (16.72,6.22) circle (0.079);
\fill[fill={rgb,255:red,43;green,191;blue,143}, fill opacity=0.89] (15.53,12.44) circle (0.088);
\fill[fill={rgb,255:red,43;green,191;blue,142}, fill opacity=0.92] (14.35,11.76) circle (0.091);
\fill[fill={rgb,255:red,43;green,191;blue,142}, fill opacity=0.91] (15.63,11.95) circle (0.090);
\fill[fill={rgb,255:red,43;green,191;blue,142}, fill opacity=0.90] (16.79,11.37) circle (0.089);
\fill[fill={rgb,255:red,143;green,156;blue,102}, fill opacity=0.74] (19.13,7.19) circle (0.069);
\fill[fill={rgb,255:red,49;green,189;blue,139}, fill opacity=0.86] (17.95,9.63) circle (0.084);
\fill[fill={rgb,255:red,43;green,191;blue,141}, fill opacity=0.91] (14.58,9.31) circle (0.090);
\fill[fill={rgb,255:red,43;green,191;blue,142}, fill opacity=0.90] (13.77,10.34) circle (0.089);
\fill[fill={rgb,255:red,126;green,162;blue,109}, fill opacity=0.75] (14.61,7.48) circle (0.070);
\fill[fill={rgb,255:red,43;green,191;blue,142}, fill opacity=0.93] (16.06,10.96) circle (0.093);
\fill[fill={rgb,255:red,43;green,190;blue,144}, fill opacity=0.88] (13.89,12.19) circle (0.086);
\fill[fill={rgb,255:red,43;green,190;blue,144}, fill opacity=0.88] (14.06,12.39) circle (0.086);
\fill[fill={rgb,255:red,43;green,191;blue,142}, fill opacity=0.93] (16.24,10.91) circle (0.092);
\fill[fill={rgb,255:red,44;green,189;blue,145}, fill opacity=0.86] (13.59,12.07) circle (0.084);
\fill[fill={rgb,255:red,43;green,191;blue,141}, fill opacity=0.94] (15.51,9.96) circle (0.094);
\fill[fill={rgb,255:red,43;green,191;blue,142}, fill opacity=0.92] (14.16,9.84) circle (0.091);
\fill[fill={rgb,255:red,43;green,189;blue,145}, fill opacity=0.87] (13.56,11.96) circle (0.084);
\fill[fill={rgb,255:red,43;green,191;blue,142}, fill opacity=0.91] (13.76,10.58) circle (0.089);
\fill[fill={rgb,255:red,239;green,123;blue,63}, fill opacity=0.92] (11.62,5.85) circle (0.091);
\fill[fill={rgb,255:red,68;green,182;blue,133}, fill opacity=0.81] (13.97,8.23) circle (0.077);
\fill[fill={rgb,255:red,169;green,141;blue,108}, fill opacity=0.74] (11.35,7.57) circle (0.069);
\fill[fill={rgb,255:red,238;green,123;blue,64}, fill opacity=0.90] (9.78,4.27) circle (0.089);
\fill[fill={rgb,255:red,236;green,124;blue,64}, fill opacity=0.90] (14.11,6.24) circle (0.088);
\fill[fill={rgb,255:red,239;green,123;blue,63}, fill opacity=0.94] (12.68,5.56) circle (0.094);
\fill[fill={rgb,255:red,234;green,124;blue,65}, fill opacity=0.88] (14.16,3.41) circle (0.086);
\fill[fill={rgb,255:red,229;green,126;blue,67}, fill opacity=0.86] (14.02,2.88) circle (0.084);
\fill[fill={rgb,255:red,239;green,123;blue,63}, fill opacity=0.93] (11.76,4.46) circle (0.093);
\fill[fill={rgb,255:red,237;green,124;blue,64}, fill opacity=0.90] (15.34,5.58) circle (0.089);
\fill[fill={rgb,255:red,238;green,123;blue,63}, fill opacity=0.91] (12.21,3.63) circle (0.089);
\fill[fill={rgb,255:red,201;green,136;blue,78}, fill opacity=0.81] (18.95,4.99) circle (0.077);
\fill[fill={rgb,255:red,192;green,139;blue,82}, fill opacity=0.79] (19.92,5.25) circle (0.076);
\fill[fill={rgb,255:red,223;green,126;blue,75}, fill opacity=0.83] (7.72,3.50) circle (0.081);
\fill[fill={rgb,255:red,237;green,123;blue,64}, fill opacity=0.90] (10.93,3.53) circle (0.089);
\fill[fill={rgb,255:red,239;green,123;blue,63}, fill opacity=0.94] (11.94,5.00) circle (0.094);
\fill[fill={rgb,255:red,238;green,123;blue,63}, fill opacity=0.90] (13.00,6.26) circle (0.089);
\fill[fill={rgb,255:red,174;green,134;blue,117}, fill opacity=0.76] (6.95,4.98) circle (0.072);
\fill[fill={rgb,255:red,239;green,123;blue,63}, fill opacity=0.93] (12.64,4.26) circle (0.092);
\fill[fill={rgb,255:red,239;green,123;blue,63}, fill opacity=0.93] (10.97,4.96) circle (0.093);
\fill[fill={rgb,255:red,238;green,123;blue,63}, fill opacity=0.91] (10.10,5.10) circle (0.090);
\fill[fill={rgb,255:red,215;green,127;blue,81}, fill opacity=0.81] (7.37,2.28) circle (0.078);
\fill[fill={rgb,255:red,223;green,128;blue,69}, fill opacity=0.84] (13.74,6.76) circle (0.082);
\fill[fill={rgb,255:red,239;green,123;blue,63}, fill opacity=0.93] (14.20,4.70) circle (0.092);
\fill[fill={rgb,255:red,237;green,123;blue,63}, fill opacity=0.90] (14.02,3.93) circle (0.089);
\fill[fill={rgb,255:red,89;green,166;blue,145}, fill opacity=0.72] (12.09,8.50) circle (0.066);
\fill[fill={rgb,255:red,235;green,124;blue,66}, fill opacity=0.88] (9.23,3.64) circle (0.086);
\fill[fill={rgb,255:red,234;green,125;blue,65}, fill opacity=0.88] (15.90,5.42) circle (0.087);
\fill[fill={rgb,255:red,237;green,123;blue,64}, fill opacity=0.89] (9.65,4.06) circle (0.088);
\fill[fill={rgb,255:red,186;green,141;blue,84}, fill opacity=0.79] (15.07,6.94) circle (0.075);
\fill[fill={rgb,255:red,238;green,123;blue,63}, fill opacity=0.91] (10.05,4.67) circle (0.090);
\fill[fill={rgb,255:red,221;green,129;blue,70}, fill opacity=0.84] (16.91,4.11) circle (0.082);
\fill[fill={rgb,255:red,222;green,129;blue,70}, fill opacity=0.84] (15.61,3.07) circle (0.081);
\fill[fill={rgb,255:red,239;green,123;blue,63}, fill opacity=0.92] (10.48,4.41) circle (0.091);
\fill[fill={rgb,255:red,152;green,142;blue,124}, fill opacity=0.72] (10.96,7.63) circle (0.067);
\fill[fill={rgb,255:red,234;green,124;blue,66}, fill opacity=0.87] (9.14,3.57) circle (0.085);
\fill[fill={rgb,255:red,204;green,129;blue,90}, fill opacity=0.79] (9.78,6.57) circle (0.076);
\fill[fill={rgb,255:red,235;green,124;blue,64}, fill opacity=0.89] (15.56,4.68) circle (0.088);
\fill[fill={rgb,255:red,236;green,124;blue,64}, fill opacity=0.90] (13.53,3.59) circle (0.088);
\fill[fill={rgb,255:red,221;green,127;blue,76}, fill opacity=0.82] (8.35,1.88) circle (0.079);
\fill[fill={rgb,255:red,223;green,128;blue,69}, fill opacity=0.85] (16.66,5.69) circle (0.082);
\fill[fill={rgb,255:red,239;green,123;blue,63}, fill opacity=0.93] (13.63,5.84) circle (0.093);
\fill[fill={rgb,255:red,239;green,123;blue,63}, fill opacity=0.95] (12.02,5.14) circle (0.094);
\fill[fill={rgb,255:red,229;green,125;blue,69}, fill opacity=0.85] (9.43,2.36) circle (0.082);
\fill[fill={rgb,255:red,239;green,123;blue,63}, fill opacity=0.94] (13.01,4.64) circle (0.094);
\fill[fill={rgb,255:red,239;green,123;blue,63}, fill opacity=0.94] (12.41,5.74) circle (0.093);
\fill[fill={rgb,255:red,233;green,125;blue,65}, fill opacity=0.88] (16.07,5.14) circle (0.086);
\fill[fill={rgb,255:red,237;green,123;blue,64}, fill opacity=0.89] (10.17,5.60) circle (0.088);
\fill[fill={rgb,255:red,52;green,188;blue,138}, fill opacity=0.86] (14.84,8.37) circle (0.083);
\fill[fill={rgb,255:red,176;green,134;blue,115}, fill opacity=0.76] (5.37,2.67) circle (0.072);
\fill[fill={rgb,255:red,88;green,159;blue,165}, fill opacity=0.72] (11.46,8.51) circle (0.066);
\fill[fill={rgb,255:red,238;green,123;blue,64}, fill opacity=0.90] (10.22,5.45) circle (0.089);
\fill[fill={rgb,255:red,239;green,123;blue,63}, fill opacity=0.94] (13.68,4.95) circle (0.094);
\fill[fill={rgb,255:red,239;green,123;blue,63}, fill opacity=0.94] (13.70,4.94) circle (0.094);
\fill[fill={rgb,255:red,239;green,123;blue,63}, fill opacity=0.92] (10.45,4.70) circle (0.091);
\fill[fill={rgb,255:red,237;green,123;blue,63}, fill opacity=0.90] (14.58,4.23) circle (0.089);
\fill[fill={rgb,255:red,239;green,123;blue,63}, fill opacity=0.93] (13.21,4.44) circle (0.093);
\fill[fill={rgb,255:red,219;green,130;blue,72}, fill opacity=0.83] (13.33,6.90) circle (0.080);
\fill[fill={rgb,255:red,238;green,123;blue,63}, fill opacity=0.91] (11.58,5.99) circle (0.090);
\fill[fill={rgb,255:red,239;green,123;blue,63}, fill opacity=0.94] (11.69,5.44) circle (0.093);
\fill[fill={rgb,255:red,239;green,123;blue,63}, fill opacity=0.95] (13.37,5.20) circle (0.095);
\fill[fill={rgb,255:red,238;green,123;blue,63}, fill opacity=0.90] (10.52,3.84) circle (0.089);
\fill[fill={rgb,255:red,225;green,125;blue,74}, fill opacity=0.84] (8.38,4.97) circle (0.082);
\fill[fill={rgb,255:red,237;green,124;blue,64}, fill opacity=0.90] (15.49,5.26) circle (0.089);
\fill[fill={rgb,255:red,234;green,124;blue,65}, fill opacity=0.89] (15.14,6.03) circle (0.087);
\fill[fill={rgb,255:red,142;green,157;blue,102}, fill opacity=0.74] (15.12,7.27) circle (0.069);
\fill[fill={rgb,255:red,219;green,129;blue,71}, fill opacity=0.84] (15.90,3.06) circle (0.081);
\fill[fill={rgb,255:red,225;green,128;blue,69}, fill opacity=0.85] (15.69,3.36) circle (0.082);
\fill[fill={rgb,255:red,223;green,128;blue,70}, fill opacity=0.84] (12.85,6.86) circle (0.081);
\fill[fill={rgb,255:red,238;green,123;blue,64}, fill opacity=0.90] (10.94,3.58) circle (0.089);
\fill[fill={rgb,255:red,60;green,185;blue,135}, fill opacity=0.84] (15.16,8.08) circle (0.081);
\fill[fill={rgb,255:red,239;green,123;blue,63}, fill opacity=0.94] (11.82,5.39) circle (0.094);
\fill[fill={rgb,255:red,60;green,153;blue,214}, fill opacity=0.89] (8.93,8.10) circle (0.087);
\fill[fill={rgb,255:red,226;green,125;blue,74}, fill opacity=0.84] (8.51,5.12) circle (0.082);
\fill[fill={rgb,255:red,238;green,123;blue,64}, fill opacity=0.90] (11.60,6.13) circle (0.089);
\fill[fill={rgb,255:red,237;green,123;blue,64}, fill opacity=0.90] (15.13,4.49) circle (0.089);
\fill[fill={rgb,255:red,239;green,123;blue,63}, fill opacity=0.93] (12.87,4.33) circle (0.093);
\fill[fill={rgb,255:red,239;green,123;blue,63}, fill opacity=0.94] (13.08,5.68) circle (0.094);
\fill[fill={rgb,255:red,156;green,137;blue,132}, fill opacity=0.74] (5.16,3.42) circle (0.069);
\fill[fill={rgb,255:red,239;green,123;blue,63}, fill opacity=0.95] (13.33,5.38) circle (0.095);
\fill[fill={rgb,255:red,228;green,125;blue,69}, fill opacity=0.85] (11.25,6.63) circle (0.082);
\fill[fill={rgb,255:red,238;green,123;blue,63}, fill opacity=0.92] (14.68,4.73) circle (0.091);
\fill[fill={rgb,255:red,204;green,135;blue,77}, fill opacity=0.81] (18.57,5.12) circle (0.078);
\fill[fill={rgb,255:red,239;green,123;blue,63}, fill opacity=0.93] (11.08,5.34) circle (0.092);
\fill[fill={rgb,255:red,229;green,126;blue,67}, fill opacity=0.86] (16.37,5.52) circle (0.084);
\fill[fill={rgb,255:red,238;green,123;blue,63}, fill opacity=0.91] (12.70,6.14) circle (0.090);
\fill[fill={rgb,255:red,238;green,123;blue,63}, fill opacity=0.92] (14.93,5.48) circle (0.091);
\fill[fill={rgb,255:red,199;green,137;blue,79}, fill opacity=0.80] (19.20,4.92) circle (0.077);
\fill[fill={rgb,255:red,219;green,127;blue,78}, fill opacity=0.82] (7.82,2.08) circle (0.079);
\fill[fill={rgb,255:red,239;green,123;blue,63}, fill opacity=0.92] (13.99,5.95) circle (0.091);
\fill[fill={rgb,255:red,239;green,123;blue,63}, fill opacity=0.93] (14.29,4.69) circle (0.092);
\fill[fill={rgb,255:red,239;green,123;blue,63}, fill opacity=0.95] (13.10,5.55) circle (0.094);
\fill[fill={rgb,255:red,234;green,124;blue,65}, fill opacity=0.88] (13.48,3.23) circle (0.086);
\fill[fill={rgb,255:red,239;green,123;blue,63}, fill opacity=0.93] (13.53,4.53) circle (0.093);
\fill[fill={rgb,255:red,236;green,124;blue,64}, fill opacity=0.90] (14.70,6.06) circle (0.088);
\fill[fill={rgb,255:red,230;green,125;blue,69}, fill opacity=0.85] (10.59,6.35) circle (0.083);
\fill[fill={rgb,255:red,205;green,135;blue,77}, fill opacity=0.81] (14.41,6.90) circle (0.078);
\fill[fill={rgb,255:red,167;green,136;blue,122}, fill opacity=0.75] (8.76,6.41) circle (0.070);
\fill[fill={rgb,255:red,226;green,125;blue,73}, fill opacity=0.84] (8.24,4.72) circle (0.082);
\fill[fill={rgb,255:red,230;green,126;blue,66}, fill opacity=0.86] (13.34,6.63) circle (0.084);
\fill[fill={rgb,255:red,187;green,141;blue,84}, fill opacity=0.79] (17.45,6.40) circle (0.075);
\fill[fill={rgb,255:red,239;green,123;blue,63}, fill opacity=0.92] (14.74,5.31) circle (0.092);
\fill[fill={rgb,255:red,200;green,136;blue,78}, fill opacity=0.81] (18.90,5.21) circle (0.077);
\fill[fill={rgb,255:red,231;green,125;blue,66}, fill opacity=0.86] (12.81,6.64) circle (0.084);
\fill[fill={rgb,255:red,237;green,123;blue,64}, fill opacity=0.89] (9.58,4.98) circle (0.088);
\fill[fill={rgb,255:red,239;green,123;blue,63}, fill opacity=0.93] (13.67,4.58) circle (0.093);
\fill[fill={rgb,255:red,239;green,123;blue,63}, fill opacity=0.94] (11.58,5.04) circle (0.094);
\fill[fill={rgb,255:red,158;green,140;blue,123}, fill opacity=0.73] (10.56,7.43) circle (0.068);
\fill[fill={rgb,255:red,237;green,123;blue,63}, fill opacity=0.90] (13.25,3.68) circle (0.089);
\fill[fill={rgb,255:red,229;green,126;blue,67}, fill opacity=0.86] (16.03,4.11) circle (0.084);
\fill[fill={rgb,255:red,83;green,177;blue,126}, fill opacity=0.80] (16.70,7.69) circle (0.077);
\fill[fill={rgb,255:red,237;green,123;blue,63}, fill opacity=0.90] (12.38,3.54) circle (0.089);
\fill[fill={rgb,255:red,211;green,129;blue,83}, fill opacity=0.80] (10.49,6.80) circle (0.077);
\fill[fill={rgb,255:red,239;green,123;blue,63}, fill opacity=0.93] (14.22,4.85) circle (0.093);
\fill[fill={rgb,255:red,238;green,123;blue,63}, fill opacity=0.91] (13.98,6.03) circle (0.091);
\fill[fill={rgb,255:red,217;green,128;blue,78}, fill opacity=0.81] (10.54,6.71) circle (0.078);
\fill[fill={rgb,255:red,153;green,138;blue,134}, fill opacity=0.74] (8.62,6.45) circle (0.069);
\fill[fill={rgb,255:red,193;green,131;blue,101}, fill opacity=0.79] (8.20,5.74) circle (0.074);
\fill[fill={rgb,255:red,195;green,138;blue,81}, fill opacity=0.80] (18.06,6.06) circle (0.076);
\fill[fill={rgb,255:red,239;green,123;blue,63}, fill opacity=0.94] (11.90,4.91) circle (0.094);
\fill[fill={rgb,255:red,238;green,123;blue,64}, fill opacity=0.91] (9.89,4.67) circle (0.089);
\fill[fill={rgb,255:red,78;green,179;blue,128}, fill opacity=0.81] (15.77,7.74) circle (0.077);
\fill[fill={rgb,255:red,238;green,123;blue,63}, fill opacity=0.91] (10.62,5.50) circle (0.090);
\fill[fill={rgb,255:red,208;green,129;blue,87}, fill opacity=0.80] (9.60,6.43) circle (0.076);
\fill[fill={rgb,255:red,229;green,126;blue,67}, fill opacity=0.85] (12.15,2.32) circle (0.083);
\fill[fill={rgb,255:red,238;green,123;blue,63}, fill opacity=0.91] (15.24,5.27) circle (0.090);
\fill[fill={rgb,255:red,217;green,130;blue,72}, fill opacity=0.84] (16.17,6.22) circle (0.081);
\fill[fill={rgb,255:red,235;green,124;blue,66}, fill opacity=0.88] (8.96,4.39) circle (0.086);
\fill[fill={rgb,255:red,48;green,189;blue,139}, fill opacity=0.88] (16.37,8.56) circle (0.086);
\fill[fill={rgb,255:red,237;green,123;blue,64}, fill opacity=0.90] (14.62,4.20) circle (0.089);
\fill[fill={rgb,255:red,239;green,123;blue,63}, fill opacity=0.92] (12.02,5.95) circle (0.091);
\fill[fill={rgb,255:red,197;green,137;blue,80}, fill opacity=0.80] (15.13,6.83) circle (0.076);
\fill[fill={rgb,255:red,238;green,123;blue,63}, fill opacity=0.91] (14.69,4.60) circle (0.091);
\fill[fill={rgb,255:red,236;green,124;blue,64}, fill opacity=0.89] (12.16,6.40) circle (0.087);
\fill[fill={rgb,255:red,232;green,125;blue,66}, fill opacity=0.88] (14.97,6.22) circle (0.086);
\fill[fill={rgb,255:red,58;green,182;blue,147}, fill opacity=0.78] (12.74,9.08) circle (0.074);
\fill[fill={rgb,255:red,104;green,170;blue,118}, fill opacity=0.77] (15.05,7.57) circle (0.073);
\fill[fill={rgb,255:red,239;green,123;blue,63}, fill opacity=0.95] (13.09,5.39) circle (0.095);
\fill[fill={rgb,255:red,228;green,125;blue,70}, fill opacity=0.85] (10.32,6.31) circle (0.082);
\fill[fill={rgb,255:red,222;green,126;blue,76}, fill opacity=0.83] (8.58,5.37) circle (0.080);
\fill[fill={rgb,255:red,238;green,123;blue,63}, fill opacity=0.92] (15.00,5.13) circle (0.091);
\fill[fill={rgb,255:red,237;green,123;blue,64}, fill opacity=0.90] (11.08,6.02) circle (0.088);
\fill[fill={rgb,255:red,207;green,128;blue,89}, fill opacity=0.80] (8.12,5.42) circle (0.077);
\fill[fill={rgb,255:red,129;green,142;blue,154}, fill opacity=0.75] (9.08,6.93) circle (0.070);
\fill[fill={rgb,255:red,238;green,123;blue,63}, fill opacity=0.90] (12.59,3.61) circle (0.089);
\fill[fill={rgb,255:red,220;green,129;blue,70}, fill opacity=0.84] (17.10,4.27) circle (0.081);
\fill[fill={rgb,255:red,223;green,126;blue,76}, fill opacity=0.83] (7.68,3.76) circle (0.081);
\fill[fill={rgb,255:red,213;green,127;blue,83}, fill opacity=0.82] (8.65,5.70) circle (0.078);
\fill[fill={rgb,255:red,98;green,147;blue,182}, fill opacity=0.80] (6.96,6.00) circle (0.076);
\fill[fill={rgb,255:red,239;green,123;blue,63}, fill opacity=0.95] (13.27,5.33) circle (0.095);
\fill[fill={rgb,255:red,239;green,123;blue,63}, fill opacity=0.94] (12.58,5.71) circle (0.094);
\fill[fill={rgb,255:red,238;green,123;blue,63}, fill opacity=0.92] (14.68,5.73) circle (0.091);
\fill[fill={rgb,255:red,222;green,126;blue,76}, fill opacity=0.83] (7.76,2.69) circle (0.080);
\fill[fill={rgb,255:red,239;green,123;blue,63}, fill opacity=0.93] (12.17,4.29) circle (0.093);
\fill[fill={rgb,255:red,234;green,124;blue,65}, fill opacity=0.88] (12.12,6.49) circle (0.086);
\fill[fill={rgb,255:red,234;green,124;blue,65}, fill opacity=0.88] (12.96,6.51) circle (0.086);
\fill[fill={rgb,255:red,105;green,170;blue,117}, fill opacity=0.78] (18.06,7.62) circle (0.073);
\fill[fill={rgb,255:red,221;green,129;blue,70}, fill opacity=0.84] (14.82,6.57) circle (0.082);
\fill[fill={rgb,255:red,227;green,127;blue,68}, fill opacity=0.85] (12.71,2.24) circle (0.082);
\fill[fill={rgb,255:red,230;green,126;blue,66}, fill opacity=0.87] (15.40,6.14) circle (0.085);
\fill[fill={rgb,255:red,196;green,130;blue,98}, fill opacity=0.79] (7.62,5.19) circle (0.075);
\fill[fill={rgb,255:red,238;green,123;blue,63}, fill opacity=0.91] (12.26,3.66) circle (0.090);
\fill[fill={rgb,255:red,65;green,152;blue,210}, fill opacity=0.87] (7.78,7.13) circle (0.084);
\fill[fill={rgb,255:red,239;green,123;blue,63}, fill opacity=0.95] (12.83,5.24) circle (0.095);
\fill[fill={rgb,255:red,123;green,163;blue,109}, fill opacity=0.76] (18.19,7.35) circle (0.071);
\fill[fill={rgb,255:red,73;green,151;blue,203}, fill opacity=0.84] (3.69,5.20) circle (0.082);
\fill[fill={rgb,255:red,213;green,127;blue,84}, fill opacity=0.81] (8.55,5.63) circle (0.078);
\fill[fill={rgb,255:red,86;green,176;blue,124}, fill opacity=0.80] (15.98,7.63) circle (0.076);
\fill[fill={rgb,255:red,215;green,131;blue,73}, fill opacity=0.83] (17.08,5.82) circle (0.080);
\fill[fill={rgb,255:red,234;green,124;blue,66}, fill opacity=0.87] (11.81,6.50) circle (0.085);
\fill[fill={rgb,255:red,239;green,123;blue,63}, fill opacity=0.92] (11.43,5.82) circle (0.091);
\fill[fill={rgb,255:red,233;green,125;blue,65}, fill opacity=0.88] (15.67,4.27) circle (0.086);
\fill[fill={rgb,255:red,239;green,123;blue,63}, fill opacity=0.94] (12.80,4.62) circle (0.094);
\fill[fill={rgb,255:red,217;green,130;blue,72}, fill opacity=0.83] (17.50,4.55) circle (0.081);
\fill[fill={rgb,255:red,231;green,125;blue,66}, fill opacity=0.87] (14.85,3.45) circle (0.085);
\fill[fill={rgb,255:red,213;green,132;blue,74}, fill opacity=0.83] (17.70,4.11) circle (0.079);
\fill[fill={rgb,255:red,239;green,123;blue,63}, fill opacity=0.92] (12.21,5.95) circle (0.092);
\fill[fill={rgb,255:red,201;green,130;blue,93}, fill opacity=0.79] (6.39,2.77) circle (0.075);
\fill[fill={rgb,255:red,184;green,142;blue,85}, fill opacity=0.78] (15.03,6.97) circle (0.074);
\fill[fill={rgb,255:red,239;green,123;blue,63}, fill opacity=0.95] (13.68,5.22) circle (0.094);
\fill[fill={rgb,255:red,207;green,128;blue,89}, fill opacity=0.81] (7.20,4.34) circle (0.077);
\fill[fill={rgb,255:red,222;green,128;blue,70}, fill opacity=0.85] (17.00,4.41) circle (0.082);
\fill[fill={rgb,255:red,239;green,123;blue,63}, fill opacity=0.93] (11.71,5.76) circle (0.092);
\fill[fill={rgb,255:red,216;green,131;blue,72}, fill opacity=0.83] (17.67,4.92) circle (0.080);
\fill[fill={rgb,255:red,237;green,123;blue,64}, fill opacity=0.90] (11.00,6.00) circle (0.088);
\fill[fill={rgb,255:red,218;green,130;blue,71}, fill opacity=0.84] (14.95,6.59) circle (0.081);
\fill[fill={rgb,255:red,235;green,124;blue,65}, fill opacity=0.88] (11.27,3.06) circle (0.087);
\fill[fill={rgb,255:red,237;green,123;blue,63}, fill opacity=0.90] (13.87,6.18) circle (0.089);
\fill[fill={rgb,255:red,226;green,127;blue,68}, fill opacity=0.85] (15.18,3.13) circle (0.083);
\fill[fill={rgb,255:red,230;green,126;blue,66}, fill opacity=0.86] (13.72,6.59) circle (0.084);
\fill[fill={rgb,255:red,220;green,127;blue,77}, fill opacity=0.82] (7.94,2.10) circle (0.079);
\fill[fill={rgb,255:red,139;green,140;blue,146}, fill opacity=0.74] (4.29,2.95) circle (0.069);
\fill[fill={rgb,255:red,187;green,132;blue,105}, fill opacity=0.77] (9.47,6.63) circle (0.073);
\fill[fill={rgb,255:red,236;green,123;blue,65}, fill opacity=0.88] (9.56,5.27) circle (0.087);
\fill[fill={rgb,255:red,166;green,148;blue,92}, fill opacity=0.77] (16.70,6.83) circle (0.072);
\fill[fill={rgb,255:red,216;green,130;blue,72}, fill opacity=0.83] (15.72,2.72) circle (0.080);
\fill[fill={rgb,255:red,73;green,151;blue,203}, fill opacity=0.84] (3.26,5.13) circle (0.081);
\fill[fill={rgb,255:red,209;green,128;blue,87}, fill opacity=0.81] (6.93,3.58) circle (0.077);
\fill[fill={rgb,255:red,73;green,151;blue,203}, fill opacity=0.84] (7.69,6.84) circle (0.081);
\fill[fill={rgb,255:red,212;green,128;blue,85}, fill opacity=0.81] (7.67,4.80) circle (0.078);
\fill[fill={rgb,255:red,216;green,131;blue,72}, fill opacity=0.83] (16.97,5.85) circle (0.080);
\fill[fill={rgb,255:red,189;green,140;blue,83}, fill opacity=0.79] (20.58,4.41) circle (0.075);
\fill[fill={rgb,255:red,239;green,123;blue,63}, fill opacity=0.93] (12.75,4.26) circle (0.092);
\fill[fill={rgb,255:red,239;green,123;blue,63}, fill opacity=0.95] (13.46,5.47) circle (0.094);
\fill[fill={rgb,255:red,67;green,152;blue,208}, fill opacity=0.86] (6.18,6.27) circle (0.084);
\fill[fill={rgb,255:red,213;green,131;blue,75}, fill opacity=0.81] (12.84,7.05) circle (0.078);
\fill[fill={rgb,255:red,236;green,124;blue,64}, fill opacity=0.89] (15.57,5.46) circle (0.088);
\fill[fill={rgb,255:red,237;green,123;blue,64}, fill opacity=0.89] (10.04,3.58) circle (0.087);
\fill[fill={rgb,255:red,238;green,123;blue,63}, fill opacity=0.91] (14.05,4.01) circle (0.090);
\fill[fill={rgb,255:red,180;green,141;blue,93}, fill opacity=0.76] (12.23,7.50) circle (0.071);
\fill[fill={rgb,255:red,230;green,126;blue,66}, fill opacity=0.87] (15.38,6.15) circle (0.085);
\fill[fill={rgb,255:red,231;green,125;blue,66}, fill opacity=0.87] (16.10,4.51) circle (0.085);
\fill[fill={rgb,255:red,71;green,181;blue,130}, fill opacity=0.82] (16.81,7.89) circle (0.079);
\fill[fill={rgb,255:red,86;green,149;blue,192}, fill opacity=0.82] (4.87,5.06) circle (0.078);
\fill[fill={rgb,255:red,205;green,135;blue,77}, fill opacity=0.81] (15.18,6.73) circle (0.078);
\fill[fill={rgb,255:red,226;green,127;blue,68}, fill opacity=0.85] (13.88,6.67) circle (0.083);
\fill[fill={rgb,255:red,239;green,123;blue,63}, fill opacity=0.93] (11.87,5.74) circle (0.092);
\fill[fill={rgb,255:red,239;green,123;blue,63}, fill opacity=0.93] (10.95,4.54) circle (0.092);
\fill[fill={rgb,255:red,234;green,125;blue,65}, fill opacity=0.88] (14.54,3.57) circle (0.086);
\fill[fill={rgb,255:red,236;green,124;blue,64}, fill opacity=0.89] (14.62,3.96) circle (0.088);
\fill[fill={rgb,255:red,239;green,123;blue,63}, fill opacity=0.93] (13.54,4.40) circle (0.092);
\fill[fill={rgb,255:red,239;green,123;blue,63}, fill opacity=0.92] (12.48,6.05) circle (0.091);
\fill[fill={rgb,255:red,213;green,132;blue,73}, fill opacity=0.83] (17.80,5.16) circle (0.080);
\fill[fill={rgb,255:red,173;green,146;blue,89}, fill opacity=0.77] (15.94,6.88) circle (0.073);
\fill[fill={rgb,255:red,115;green,144;blue,167}, fill opacity=0.77] (7.01,5.82) circle (0.073);
\fill[fill={rgb,255:red,147;green,154;blue,101}, fill opacity=0.74] (13.82,7.50) circle (0.069);
\fill[fill={rgb,255:red,75;green,180;blue,129}, fill opacity=0.81] (17.74,8.07) circle (0.078);
\fill[fill={rgb,255:red,212;green,132;blue,74}, fill opacity=0.82] (15.56,2.37) circle (0.079);
\fill[fill={rgb,255:red,65;green,152;blue,210}, fill opacity=0.87] (3.72,6.29) circle (0.085);
\fill[fill={rgb,255:red,215;green,131;blue,72}, fill opacity=0.83] (16.34,6.20) circle (0.080);
\fill[fill={rgb,255:red,239;green,123;blue,63}, fill opacity=0.94] (13.60,5.75) circle (0.093);
\fill[fill={rgb,255:red,209;green,133;blue,75}, fill opacity=0.82] (18.00,5.35) circle (0.079);
\fill[fill={rgb,255:red,239;green,123;blue,63}, fill opacity=0.92] (11.28,5.59) circle (0.092);
\fill[fill={rgb,255:red,231;green,124;blue,69}, fill opacity=0.86] (8.85,5.09) circle (0.084);
\fill[fill={rgb,255:red,150;green,138;blue,137}, fill opacity=0.74] (5.11,3.58) circle (0.069);
\fill[fill={rgb,255:red,235;green,124;blue,65}, fill opacity=0.88] (12.34,3.03) circle (0.086);
\fill[fill={rgb,255:red,81;green,178;blue,126}, fill opacity=0.80] (16.89,7.74) circle (0.077);
\fill[fill={rgb,255:red,198;green,137;blue,79}, fill opacity=0.80] (15.98,6.61) circle (0.077);
\fill[fill={rgb,255:red,239;green,123;blue,63}, fill opacity=0.93] (13.27,4.26) circle (0.092);
\fill[fill={rgb,255:red,239;green,123;blue,63}, fill opacity=0.94] (14.15,5.05) circle (0.093);
\draw[white, line width=0.8pt] (8.07,8.87) circle (0.16);
\fill[white] (8.07,8.87) circle (0.045);
\draw[white, line width=0.8pt] (12.86,5.19) circle (0.16);
\fill[white] (12.86,5.19) circle (0.045);
\draw[white, line width=0.8pt] (15.10,10.57) circle (0.16);
\fill[white] (15.10,10.57) circle (0.045);

    \end{scope}
    % fade the lower edge of the artwork into the background
    \fill[fcnavy, path fading=artfade, fading angle=180] (0,9.6) rectangle (21,14.2);
    % title block
    \fill[fcverm] (2.0,10.35) rectangle (3.6,10.47);
    \node[anchor=north west, inner sep=0, text=white, align=left, font=\fcblack\fontsize{40}{44}\selectfont]
      at (1.95,9.85) {Fuzzy Clustering\\and Data Analysis\\in Python};
    \node[anchor=north west, inner sep=0, text=fclight, font=\fcmedium\fontsize{17}{21}\selectfont]
      at (2.0,4.95) {The {\plexmono fuzzyclust} toolbox};
    % author and affiliation
    \node[anchor=south west, inner sep=0, text=white, font=\fcsemi\fontsize{16}{19}\selectfont] at (2.0,2.0) {J\'anos Abonyi};
    \node[anchor=south west, inner sep=0, text=fclight, font=\fclightface\fontsize{10.5}{13}\selectfont] at (2.0,1.35) {University of Pannonia, Veszpr\'em};
    \node[anchor=south east, inner sep=0] at (19.0,1.3) {\fclogo[0.42]};
  \end{scope}
\end{tikzpicture}
\newpage
% ----------------------------------------------------------------- back cover
\noindent\begin{tikzpicture}[remember picture, overlay, x=1cm, y=1cm]
  \begin{scope}[shift={(current page.south west)}]
    \fill[fcnavy] (0,0) rectangle (21,29.7);
    \begin{scope}[shift={(-0.3,23.2)}, opacity=0.55]
      \clip (0.3,0) rectangle (21.3,6.5);
      \begin{scope}[yshift=-7.5cm]% generated by book/cover/make_cover_art.py -- do not edit
\draw[line width=0.45pt, draw={rgb,255:red,58;green,154;blue,217}, draw opacity=0.30] (4.71,-2.49) -- (4.74,-2.28) -- (4.78,-2.00) -- (4.81,-1.78) -- (4.85,-1.50) -- (4.88,-1.28) -- (4.91,-1.01) -- (4.95,-0.76) -- (4.98,-0.52) -- (5.02,-0.26) -- (5.06,-0.02) -- (5.09,0.21) -- (5.14,0.47) -- (5.17,0.68) -- (5.23,0.96) -- (5.27,1.17) -- (5.32,1.39) -- (5.38,1.63) -- (5.45,1.88) -- (5.51,2.09) -- (5.58,2.30) -- (5.65,2.51) -- (5.74,2.72) -- (5.81,2.90) -- (5.89,3.08) -- (6.00,3.29) -- (6.10,3.47) -- (6.19,3.64) -- (6.31,3.83) -- (6.42,3.99) -- (6.53,4.14) -- (6.67,4.33) -- (6.80,4.49) -- (6.92,4.63) -- (7.05,4.77) -- (7.18,4.91) -- (7.32,5.05) -- (7.47,5.19) -- (7.63,5.33) -- (7.79,5.47) -- (7.96,5.61) -- (8.11,5.73) -- (8.25,5.84) -- (8.41,5.97) -- (8.60,6.11) -- (8.75,6.21) -- (8.91,6.32) -- (9.11,6.46) -- (9.26,6.55) -- (9.44,6.67) -- (9.62,6.78) -- (9.77,6.88) -- (9.97,7.01) -- (10.12,7.10) -- (10.32,7.23) -- (10.48,7.34) -- (10.64,7.44) -- (10.84,7.58) -- (10.98,7.68) -- (11.13,7.80) -- (11.30,7.94) -- (11.46,8.08) -- (11.60,8.22) -- (11.73,8.36) -- (11.84,8.50) -- (11.98,8.71) -- (12.07,8.85) -- (12.17,9.06) -- (12.26,9.28) -- (12.33,9.49) -- (12.40,9.70) -- (12.45,9.91) -- (12.49,10.12) -- (12.55,10.40) -- (12.58,10.61) -- (12.62,10.90) -- (12.65,11.11) -- (12.69,11.39) -- (12.72,11.60) -- (12.76,11.88) -- (12.78,12.09) -- (12.82,12.38) -- (12.85,12.63) -- (12.87,12.87) -- (12.91,13.15) -- (12.93,13.36) -- (12.95,13.64) -- (12.97,13.93) -- (12.98,14.21) -- (12.97,14.49) -- (12.92,14.73) -- (12.84,14.91) -- (12.70,15.07) -- (12.53,15.19) -- (12.34,15.28) -- (12.13,15.36) -- (11.91,15.43) -- (11.70,15.48) -- (11.41,15.54) -- (11.19,15.58) -- (10.97,15.62) -- (10.69,15.66) -- (10.48,15.69) -- (10.19,15.72) -- (9.90,15.76) -- (9.69,15.78) -- (9.40,15.80) -- (9.11,15.82) -- (8.90,15.84) -- (8.61,15.86) -- (8.32,15.87) -- (8.04,15.88) -- (7.75,15.89) -- (7.46,15.89) -- (7.17,15.89) -- (6.89,15.89) -- (6.60,15.88) -- (6.31,15.87) -- (6.03,15.86) -- (5.74,15.85) -- (5.45,15.83) -- (5.24,15.81) -- (4.95,15.79) -- (4.66,15.76) -- (4.45,15.74) -- (4.16,15.70) -- (3.94,15.68) -- (3.66,15.64) -- (3.44,15.61) -- (3.15,15.56) -- (2.94,15.52) -- (2.67,15.48) -- (2.44,15.43) -- (2.22,15.39) -- (1.97,15.33) -- (1.72,15.28) -- (1.50,15.23) -- (1.29,15.17) -- (1.07,15.12) -- (0.82,15.05) -- (0.57,14.98) -- (0.35,14.92) -- (0.14,14.86) -- (-0.08,14.79) -- (-0.29,14.72) -- (-0.51,14.65) -- (-0.72,14.58) -- (-0.94,14.50) -- (-1.15,14.42) -- (-1.35,14.35) -- (-1.53,14.28) -- (-1.73,14.20) -- (-1.94,14.11);
\draw[line width=0.45pt, draw={rgb,255:red,58;green,154;blue,217}, draw opacity=0.30] (14.63,20.69) -- (14.62,20.62) -- (14.60,20.55) -- (14.58,20.48) -- (14.57,20.41) -- (14.57,20.41) -- (14.55,20.34) -- (14.54,20.27) -- (14.52,20.20) -- (14.50,20.12) -- (14.50,20.10) -- (14.49,20.05) -- (14.47,19.98) -- (14.46,19.91) -- (14.44,19.84) -- (14.43,19.77) -- (14.43,19.77) -- (14.41,19.70) -- (14.39,19.63) -- (14.38,19.56) -- (14.36,19.49) -- (14.35,19.45) -- (14.35,19.42) -- (14.33,19.35) -- (14.32,19.28) -- (14.30,19.21) -- (14.29,19.14) -- (14.28,19.11) -- (14.27,19.07) -- (14.26,19.00) -- (14.24,18.93) -- (14.23,18.86) -- (14.21,18.79) -- (14.21,18.77) -- (14.20,18.72) -- (14.19,18.65) -- (14.17,18.57) -- (14.16,18.50) -- (14.14,18.43) -- (14.14,18.41) -- (14.13,18.36) -- (14.12,18.29) -- (14.10,18.22) -- (14.09,18.15) -- (14.08,18.08) -- (14.07,18.03) -- (14.06,18.01) -- (14.05,17.94) -- (14.04,17.87) -- (14.03,17.80) -- (14.01,17.73) -- (14.00,17.66) -- (14.00,17.62) -- (13.99,17.59) -- (13.98,17.52) -- (13.97,17.45) -- (13.96,17.38) -- (13.95,17.31) -- (13.94,17.24) -- (13.93,17.17) -- (13.92,17.14) -- (13.92,17.10) -- (13.91,17.03) -- (13.90,16.95) -- (13.89,16.88) -- (13.89,16.81) -- (13.88,16.74) -- (13.88,16.67) -- (13.87,16.60) -- (13.87,16.53) -- (13.87,16.46) -- (13.87,16.39) -- (13.87,16.32) -- (13.88,16.25) -- (13.89,16.18) -- (13.90,16.11) -- (13.91,16.04) -- (13.92,16.01) -- (13.94,15.97) -- (13.96,15.90) -- (14.00,15.84) -- (14.00,15.83) -- (14.05,15.76) -- (14.07,15.74) -- (14.11,15.69) -- (14.14,15.66) -- (14.19,15.62) -- (14.21,15.61) -- (14.28,15.56) -- (14.30,15.55) -- (14.35,15.51) -- (14.43,15.48) -- (14.43,15.48) -- (14.50,15.44) -- (14.57,15.41) -- (14.59,15.40) -- (14.64,15.39) -- (14.71,15.36) -- (14.78,15.34) -- (14.79,15.33) -- (14.86,15.31) -- (14.93,15.29) -- (15.00,15.27) -- (15.01,15.26) -- (15.07,15.25) -- (15.14,15.23) -- (15.22,15.21) -- (15.27,15.19) -- (15.29,15.19) -- (15.36,15.17) -- (15.43,15.15) -- (15.50,15.14) -- (15.56,15.12) -- (15.57,15.12) -- (15.65,15.10) -- (15.72,15.09) -- (15.79,15.07) -- (15.86,15.05) -- (15.87,15.05) -- (15.93,15.04) -- (16.01,15.02) -- (16.08,15.01) -- (16.15,14.99) -- (16.21,14.98) -- (16.22,14.98) -- (16.29,14.97) -- (16.36,14.95) -- (16.44,14.94) -- (16.51,14.92) -- (16.58,14.91) -- (16.58,14.91) -- (16.65,14.90) -- (16.72,14.89) -- (16.80,14.87) -- (16.87,14.86) -- (16.94,14.85) -- (16.98,14.84) -- (17.01,14.84) -- (17.08,14.82) -- (17.15,14.81) -- (17.23,14.80) -- (17.30,14.79) -- (17.37,14.78) -- (17.43,14.77) -- (17.44,14.77) -- (17.51,14.76) -- (17.59,14.75) -- (17.66,14.74) -- (17.73,14.73) -- (17.80,14.72) -- (17.87,14.71) -- (17.94,14.70) -- (17.96,14.70) -- (18.02,14.69) -- (18.09,14.69) -- (18.16,14.68) -- (18.23,14.67) -- (18.30,14.66) -- (18.38,14.66) -- (18.45,14.65) -- (18.52,14.64) -- (18.59,14.64) -- (18.65,14.63) -- (18.66,14.63) -- (18.73,14.62) -- (18.81,14.62) -- (18.88,14.61) -- (18.95,14.61) -- (19.02,14.60) -- (19.09,14.60) -- (19.16,14.59) -- (19.24,14.59) -- (19.31,14.59) -- (19.38,14.58) -- (19.45,14.58) -- (19.52,14.58) -- (19.60,14.58) -- (19.67,14.57) -- (19.74,14.57) -- (19.81,14.57) -- (19.88,14.57) -- (19.95,14.57) -- (20.03,14.57) -- (20.10,14.57) -- (20.17,14.57) -- (20.24,14.57) -- (20.31,14.57) -- (20.39,14.57) -- (20.46,14.57) -- (20.53,14.57) -- (20.60,14.57) -- (20.67,14.58) -- (20.74,14.58) -- (20.82,14.58) -- (20.89,14.58) -- (20.96,14.59) -- (21.03,14.59) -- (21.10,14.59) -- (21.18,14.60) -- (21.25,14.60) -- (21.32,14.61) -- (21.39,14.61) -- (21.46,14.61) -- (21.53,14.62) -- (21.61,14.63) -- (21.67,14.63) -- (21.68,14.63) -- (21.75,14.64) -- (21.82,14.64) -- (21.89,14.65) -- (21.97,14.66) -- (22.04,14.66) -- (22.11,14.67) -- (22.18,14.68) -- (22.25,14.68) -- (22.32,14.69) -- (22.40,14.70) -- (22.42,14.70) -- (22.47,14.71) -- (22.54,14.71) -- (22.61,14.72) -- (22.68,14.73) -- (22.75,14.74) -- (22.83,14.75) -- (22.90,14.76) -- (22.97,14.77) -- (23.00,14.77) -- (23.04,14.78) -- (23.11,14.79) -- (23.19,14.80) -- (23.26,14.81) -- (23.33,14.82) -- (23.40,14.83) -- (23.47,14.84) -- (23.51,14.84) -- (23.54,14.85) -- (23.62,14.86) -- (23.69,14.87);
\draw[line width=0.45pt, draw={rgb,255:red,58;green,154;blue,217}, draw opacity=0.44] (0.38,-2.49) -- (0.50,-2.30) -- (0.59,-2.14) -- (0.71,-1.93) -- (0.80,-1.78) -- (0.92,-1.57) -- (1.01,-1.43) -- (1.13,-1.22) -- (1.22,-1.08) -- (1.35,-0.87) -- (1.43,-0.73) -- (1.56,-0.52) -- (1.65,-0.37) -- (1.78,-0.16) -- (1.87,-0.02) -- (2.00,0.18) -- (2.10,0.33) -- (2.22,0.51) -- (2.33,0.68) -- (2.44,0.84) -- (2.57,1.03) -- (2.67,1.17) -- (2.79,1.35) -- (2.92,1.53) -- (3.02,1.67) -- (3.15,1.84) -- (3.29,2.02) -- (3.40,2.16) -- (3.51,2.30) -- (3.66,2.48) -- (3.80,2.65) -- (3.92,2.80) -- (4.05,2.94) -- (4.17,3.08) -- (4.30,3.22) -- (4.45,3.37) -- (4.59,3.52) -- (4.73,3.66) -- (4.88,3.80) -- (5.02,3.94) -- (5.16,4.07) -- (5.31,4.21) -- (5.46,4.34) -- (5.62,4.49) -- (5.79,4.63) -- (5.95,4.76) -- (6.10,4.88) -- (6.24,4.99) -- (6.40,5.12) -- (6.59,5.26) -- (6.74,5.38) -- (6.89,5.48) -- (7.06,5.61) -- (7.25,5.74) -- (7.39,5.84) -- (7.56,5.97) -- (7.75,6.09) -- (7.89,6.19) -- (8.09,6.32) -- (8.25,6.43) -- (8.41,6.53) -- (8.61,6.66) -- (8.75,6.76) -- (8.95,6.88) -- (9.11,6.99) -- (9.27,7.09) -- (9.47,7.23) -- (9.62,7.32) -- (9.79,7.44) -- (9.97,7.57) -- (10.12,7.68) -- (10.28,7.80) -- (10.46,7.94) -- (10.62,8.07) -- (10.76,8.20) -- (10.91,8.34) -- (11.05,8.48) -- (11.19,8.64) -- (11.31,8.78) -- (11.41,8.93) -- (11.55,9.14) -- (11.63,9.28) -- (11.73,9.49) -- (11.82,9.70) -- (11.90,9.91) -- (11.97,10.12) -- (12.02,10.33) -- (12.07,10.54) -- (12.11,10.83) -- (12.14,11.04) -- (12.15,11.32) -- (12.14,11.60) -- (12.12,11.81) -- (12.06,12.08) -- (11.98,12.28) -- (11.90,12.45) -- (11.77,12.65) -- (11.63,12.80) -- (11.48,12.93) -- (11.34,13.03) -- (11.12,13.14) -- (10.93,13.22) -- (10.72,13.29) -- (10.48,13.35) -- (10.26,13.40) -- (9.99,13.43) -- (9.76,13.45) -- (9.47,13.47) -- (9.18,13.47) -- (8.90,13.47) -- (8.61,13.45) -- (8.36,13.43) -- (8.11,13.41) -- (7.82,13.37) -- (7.60,13.34) -- (7.33,13.29) -- (7.10,13.25) -- (6.89,13.20) -- (6.65,13.15) -- (6.38,13.08) -- (6.17,13.03) -- (5.95,12.97) -- (5.74,12.90) -- (5.52,12.83) -- (5.31,12.76) -- (5.09,12.69) -- (4.88,12.61) -- (4.66,12.52) -- (4.47,12.45) -- (4.30,12.37) -- (4.09,12.28) -- (3.87,12.18) -- (3.69,12.09) -- (3.51,12.01) -- (3.30,11.90) -- (3.14,11.81) -- (2.94,11.70) -- (2.75,11.60) -- (2.58,11.50) -- (2.40,11.39) -- (2.22,11.28) -- (2.06,11.18) -- (1.86,11.05) -- (1.72,10.95) -- (1.54,10.83) -- (1.36,10.70) -- (1.21,10.59) -- (1.05,10.47) -- (0.87,10.33) -- (0.71,10.21) -- (0.57,10.10) -- (0.42,9.98) -- (0.26,9.84) -- (0.09,9.70) -- (-0.08,9.56) -- (-0.22,9.43) -- (-0.37,9.30) -- (-0.51,9.17) -- (-0.65,9.04) -- (-0.80,8.91) -- (-0.94,8.78) -- (-1.09,8.64) -- (-1.24,8.50) -- (-1.38,8.36) -- (-1.53,8.22) -- (-1.68,8.08) -- (-1.83,7.94) -- (-1.97,7.80);
\draw[line width=0.45pt, draw={rgb,255:red,58;green,154;blue,217}, draw opacity=0.44] (20.59,20.69) -- (20.57,20.62) -- (20.56,20.55) -- (20.55,20.48) -- (20.53,20.41) -- (20.53,20.37) -- (20.52,20.34) -- (20.51,20.27) -- (20.51,20.20) -- (20.50,20.12) -- (20.50,20.05) -- (20.49,19.98) -- (20.49,19.91) -- (20.49,19.84) -- (20.50,19.77) -- (20.50,19.70) -- (20.51,19.63) -- (20.52,19.56) -- (20.53,19.53) -- (20.54,19.49) -- (20.56,19.42) -- (20.58,19.35) -- (20.60,19.30) -- (20.61,19.28) -- (20.65,19.21) -- (20.67,19.16) -- (20.69,19.14) -- (20.74,19.07) -- (20.74,19.06) -- (20.80,19.00) -- (20.82,18.98) -- (20.88,18.93) -- (20.89,18.92) -- (20.96,18.86) -- (20.97,18.86) -- (21.03,18.82) -- (21.09,18.79) -- (21.10,18.78) -- (21.18,18.74) -- (21.24,18.72) -- (21.25,18.72) -- (21.32,18.69) -- (21.39,18.67) -- (21.46,18.65) -- (21.48,18.65) -- (21.53,18.63) -- (21.61,18.62) -- (21.68,18.61) -- (21.75,18.60) -- (21.82,18.59) -- (21.89,18.59) -- (21.97,18.58) -- (22.04,18.58) -- (22.11,18.58) -- (22.18,18.58) -- (22.25,18.58) -- (22.32,18.58) -- (22.40,18.58) -- (22.47,18.59) -- (22.54,18.59) -- (22.61,18.60) -- (22.68,18.60) -- (22.75,18.61) -- (22.83,18.62) -- (22.90,18.63) -- (22.97,18.64) -- (23.04,18.65) -- (23.04,18.65) -- (23.11,18.66) -- (23.19,18.67) -- (23.26,18.68) -- (23.33,18.69) -- (23.40,18.70) -- (23.47,18.72) -- (23.48,18.72) -- (23.54,18.73) -- (23.62,18.74) -- (23.69,18.76);
\draw[line width=0.45pt, draw={rgb,255:red,58;green,154;blue,217}, draw opacity=0.58] (0.50,3.35) -- (0.78,3.34) -- (1.00,3.36) -- (1.29,3.41) -- (1.50,3.47) -- (1.72,3.53) -- (1.93,3.60) -- (2.15,3.68) -- (2.36,3.76) -- (2.58,3.85) -- (2.75,3.92) -- (2.94,4.00) -- (3.15,4.10) -- (3.37,4.20) -- (3.52,4.27) -- (3.73,4.38) -- (3.94,4.48) -- (4.09,4.56) -- (4.30,4.67) -- (4.49,4.77) -- (4.66,4.86) -- (4.88,4.97) -- (5.02,5.05) -- (5.24,5.17) -- (5.41,5.26) -- (5.59,5.36) -- (5.79,5.47) -- (5.95,5.57) -- (6.16,5.68) -- (6.31,5.77) -- (6.53,5.89) -- (6.67,5.98) -- (6.89,6.10) -- (7.03,6.19) -- (7.25,6.31) -- (7.39,6.40) -- (7.60,6.53) -- (7.75,6.62) -- (7.95,6.74) -- (8.11,6.84) -- (8.29,6.95) -- (8.47,7.06) -- (8.62,7.16) -- (8.83,7.30) -- (8.97,7.39) -- (9.15,7.51) -- (9.33,7.64) -- (9.47,7.75) -- (9.63,7.87) -- (9.82,8.01) -- (9.97,8.14) -- (10.12,8.26) -- (10.26,8.39) -- (10.41,8.53) -- (10.55,8.68) -- (10.69,8.84) -- (10.82,8.99) -- (10.93,9.14) -- (11.05,9.32) -- (11.15,9.49) -- (11.27,9.70) -- (11.34,9.86) -- (11.41,10.05) -- (11.48,10.27) -- (11.54,10.54) -- (11.57,10.76) -- (11.57,11.04) -- (11.54,11.25) -- (11.47,11.46) -- (11.36,11.67) -- (11.24,11.81) -- (11.07,11.95) -- (10.91,12.05) -- (10.69,12.14) -- (10.48,12.20) -- (10.26,12.24) -- (9.97,12.27) -- (9.69,12.28) -- (9.40,12.26) -- (9.11,12.24) -- (8.90,12.21) -- (8.66,12.16) -- (8.39,12.11) -- (8.18,12.06) -- (7.96,12.00) -- (7.75,11.94) -- (7.53,11.87) -- (7.32,11.79) -- (7.10,11.72) -- (6.89,11.63) -- (6.67,11.54) -- (6.49,11.46) -- (6.31,11.38) -- (6.10,11.27) -- (5.91,11.18) -- (5.74,11.09) -- (5.52,10.97) -- (5.38,10.88) -- (5.17,10.76) -- (5.02,10.66) -- (4.84,10.54) -- (4.66,10.42) -- (4.52,10.32) -- (4.33,10.19) -- (4.16,10.06) -- (4.01,9.95) -- (3.87,9.83) -- (3.71,9.70) -- (3.54,9.56) -- (3.39,9.42) -- (3.23,9.28) -- (3.08,9.14) -- (2.94,8.99) -- (2.80,8.85) -- (2.66,8.71) -- (2.53,8.57) -- (2.40,8.43) -- (2.28,8.29) -- (2.15,8.14) -- (2.00,7.97) -- (1.86,7.80) -- (1.75,7.66) -- (1.65,7.51) -- (1.50,7.32) -- (1.39,7.16) -- (1.29,7.02) -- (1.15,6.81) -- (1.05,6.67) -- (0.93,6.47) -- (0.84,6.32) -- (0.71,6.11) -- (0.63,5.97) -- (0.52,5.75) -- (0.42,5.57) -- (0.35,5.40) -- (0.25,5.19) -- (0.17,4.98) -- (0.09,4.77) -- (0.03,4.56) -- (-0.02,4.34) -- (-0.05,4.06) -- (-0.02,3.78) -- (0.07,3.60) -- (0.21,3.46) -- (0.42,3.37);
\draw[line width=0.45pt, draw={rgb,255:red,58;green,154;blue,217}, draw opacity=0.72] (4.45,6.03) -- (4.59,6.02) -- (4.73,6.03) -- (4.83,6.04) -- (4.95,6.05) -- (5.09,6.07) -- (5.24,6.10) -- (5.31,6.12) -- (5.45,6.15) -- (5.53,6.18) -- (5.67,6.22) -- (5.76,6.25) -- (5.88,6.29) -- (5.97,6.32) -- (6.10,6.36) -- (6.17,6.39) -- (6.31,6.45) -- (6.38,6.48) -- (6.50,6.53) -- (6.60,6.57) -- (6.67,6.60) -- (6.82,6.67) -- (6.89,6.70) -- (6.97,6.74) -- (7.10,6.81) -- (7.17,6.84) -- (7.25,6.88) -- (7.38,6.95) -- (7.46,6.99) -- (7.53,7.03) -- (7.64,7.09) -- (7.75,7.15) -- (7.82,7.19) -- (7.89,7.23) -- (8.01,7.30) -- (8.11,7.36) -- (8.18,7.41) -- (8.25,7.45) -- (8.35,7.51) -- (8.46,7.59) -- (8.54,7.64) -- (8.61,7.69) -- (8.68,7.73) -- (8.77,7.80) -- (8.87,7.87) -- (8.97,7.94) -- (9.04,7.99) -- (9.11,8.05) -- (9.18,8.10) -- (9.26,8.16) -- (9.33,8.22) -- (9.42,8.29) -- (9.50,8.36) -- (9.58,8.43) -- (9.66,8.50) -- (9.74,8.57) -- (9.81,8.64) -- (9.88,8.71) -- (9.95,8.78) -- (10.02,8.85) -- (10.08,8.92) -- (10.14,8.99) -- (10.20,9.06) -- (10.26,9.14) -- (10.33,9.23) -- (10.41,9.33) -- (10.47,9.42) -- (10.51,9.49) -- (10.55,9.56) -- (10.62,9.68) -- (10.67,9.77) -- (10.70,9.84) -- (10.76,9.98) -- (10.78,10.05) -- (10.82,10.19) -- (10.84,10.26) -- (10.86,10.40) -- (10.86,10.54) -- (10.84,10.69) -- (10.83,10.76) -- (10.77,10.90) -- (10.72,10.97) -- (10.67,11.04) -- (10.59,11.11) -- (10.48,11.18) -- (10.41,11.22) -- (10.32,11.25) -- (10.19,11.29) -- (10.05,11.31) -- (9.97,11.32) -- (9.83,11.33) -- (9.69,11.33) -- (9.54,11.32) -- (9.47,11.31) -- (9.33,11.30) -- (9.18,11.27) -- (9.06,11.25) -- (8.97,11.23) -- (8.83,11.19) -- (8.75,11.17) -- (8.61,11.13) -- (8.53,11.11) -- (8.39,11.06) -- (8.32,11.04) -- (8.18,10.98) -- (8.11,10.96) -- (7.96,10.90) -- (7.89,10.87) -- (7.80,10.83) -- (7.68,10.77) -- (7.60,10.74) -- (7.50,10.69) -- (7.39,10.63) -- (7.32,10.59) -- (7.22,10.54) -- (7.10,10.48) -- (7.03,10.44) -- (6.96,10.40) -- (6.85,10.33) -- (6.74,10.27) -- (6.67,10.22) -- (6.60,10.18) -- (6.51,10.12) -- (6.41,10.05) -- (6.31,9.98) -- (6.24,9.93) -- (6.17,9.88) -- (6.10,9.83) -- (6.01,9.77) -- (5.92,9.70) -- (5.83,9.63) -- (5.75,9.56) -- (5.67,9.49) -- (5.59,9.43) -- (5.52,9.37) -- (5.45,9.31) -- (5.38,9.24) -- (5.31,9.18) -- (5.24,9.11) -- (5.16,9.04) -- (5.09,8.97) -- (5.02,8.90) -- (4.95,8.82) -- (4.88,8.75) -- (4.80,8.67) -- (4.73,8.59) -- (4.66,8.51) -- (4.60,8.43) -- (4.54,8.36) -- (4.48,8.29) -- (4.43,8.22) -- (4.37,8.15) -- (4.30,8.05) -- (4.23,7.94) -- (4.18,7.87) -- (4.14,7.80) -- (4.09,7.71) -- (4.02,7.59) -- (3.98,7.51) -- (3.94,7.44) -- (3.88,7.30) -- (3.86,7.23) -- (3.81,7.09) -- (3.79,7.02) -- (3.76,6.88) -- (3.74,6.74) -- (3.75,6.60) -- (3.78,6.46) -- (3.81,6.39) -- (3.87,6.29) -- (3.94,6.21) -- (4.01,6.16) -- (4.12,6.11) -- (4.23,6.07) -- (4.37,6.04) -- (4.45,6.03);
\draw[line width=0.45pt, draw={rgb,255:red,58;green,154;blue,217}, draw opacity=0.86] (5.95,7.29) -- (6.03,7.27) -- (6.10,7.25) -- (6.17,7.24) -- (6.24,7.24) -- (6.31,7.24) -- (6.38,7.24) -- (6.46,7.25) -- (6.53,7.26) -- (6.60,7.27) -- (6.67,7.29) -- (6.73,7.30) -- (6.74,7.31) -- (6.82,7.33) -- (6.89,7.35) -- (6.96,7.37) -- (6.97,7.37) -- (7.03,7.39) -- (7.10,7.42) -- (7.17,7.44) -- (7.17,7.45) -- (7.25,7.48) -- (7.32,7.51) -- (7.34,7.51) -- (7.39,7.54) -- (7.46,7.57) -- (7.49,7.59) -- (7.53,7.60) -- (7.60,7.64) -- (7.64,7.66) -- (7.68,7.67) -- (7.75,7.71) -- (7.78,7.73) -- (7.82,7.75) -- (7.89,7.79) -- (7.90,7.80) -- (7.96,7.83) -- (8.03,7.87) -- (8.04,7.87) -- (8.11,7.92) -- (8.14,7.94) -- (8.18,7.96) -- (8.25,8.01) -- (8.25,8.01) -- (8.32,8.05) -- (8.36,8.08) -- (8.39,8.10) -- (8.46,8.15) -- (8.47,8.15) -- (8.54,8.20) -- (8.56,8.22) -- (8.61,8.25) -- (8.66,8.29) -- (8.68,8.31) -- (8.75,8.36) -- (8.75,8.36) -- (8.83,8.42) -- (8.84,8.43) -- (8.90,8.48) -- (8.92,8.50) -- (8.97,8.54) -- (9.01,8.57) -- (9.04,8.60) -- (9.09,8.64) -- (9.11,8.67) -- (9.16,8.71) -- (9.18,8.73) -- (9.24,8.78) -- (9.26,8.80) -- (9.31,8.85) -- (9.33,8.88) -- (9.37,8.92) -- (9.40,8.95) -- (9.44,8.99) -- (9.47,9.03) -- (9.50,9.06) -- (9.54,9.12) -- (9.56,9.14) -- (9.61,9.21) -- (9.62,9.21) -- (9.66,9.28) -- (9.69,9.31) -- (9.71,9.35) -- (9.76,9.42) -- (9.76,9.42) -- (9.80,9.49) -- (9.83,9.55) -- (9.84,9.56) -- (9.87,9.63) -- (9.90,9.70) -- (9.90,9.71) -- (9.92,9.77) -- (9.94,9.84) -- (9.95,9.91) -- (9.96,9.98) -- (9.96,10.05) -- (9.94,10.12) -- (9.92,10.19) -- (9.90,10.22) -- (9.87,10.26) -- (9.83,10.30) -- (9.79,10.33) -- (9.76,10.35) -- (9.69,10.39) -- (9.64,10.40) -- (9.62,10.41) -- (9.54,10.42) -- (9.47,10.43) -- (9.40,10.43) -- (9.33,10.43) -- (9.26,10.43) -- (9.18,10.42) -- (9.11,10.41) -- (9.08,10.40) -- (9.04,10.40) -- (8.97,10.38) -- (8.90,10.36) -- (8.83,10.34) -- (8.79,10.33) -- (8.75,10.32) -- (8.68,10.30) -- (8.61,10.28) -- (8.57,10.26) -- (8.54,10.25) -- (8.47,10.22) -- (8.39,10.19) -- (8.39,10.19) -- (8.32,10.16) -- (8.25,10.13) -- (8.23,10.12) -- (8.18,10.10) -- (8.11,10.06) -- (8.08,10.05) -- (8.04,10.03) -- (7.96,9.99) -- (7.95,9.98) -- (7.89,9.95) -- (7.82,9.91) -- (7.82,9.91) -- (7.75,9.87) -- (7.70,9.84) -- (7.68,9.83) -- (7.60,9.78) -- (7.58,9.77) -- (7.53,9.74) -- (7.47,9.70) -- (7.46,9.69) -- (7.39,9.65) -- (7.36,9.63) -- (7.32,9.60) -- (7.26,9.56) -- (7.25,9.55) -- (7.17,9.49) -- (7.16,9.49) -- (7.10,9.44) -- (7.07,9.42) -- (7.03,9.39) -- (6.98,9.35) -- (6.96,9.33) -- (6.89,9.28) -- (6.89,9.27) -- (6.82,9.21) -- (6.81,9.21) -- (6.74,9.15) -- (6.73,9.14) -- (6.67,9.08) -- (6.65,9.06) -- (6.60,9.02) -- (6.57,8.99) -- (6.53,8.95) -- (6.50,8.92) -- (6.46,8.88) -- (6.43,8.85) -- (6.38,8.80) -- (6.36,8.78) -- (6.31,8.73) -- (6.30,8.71) -- (6.24,8.64) -- (6.24,8.64) -- (6.18,8.57) -- (6.17,8.56) -- (6.12,8.50) -- (6.10,8.47) -- (6.07,8.43) -- (6.03,8.37) -- (6.02,8.36) -- (5.97,8.29) -- (5.95,8.26) -- (5.93,8.22) -- (5.89,8.15) -- (5.88,8.14) -- (5.85,8.08) -- (5.82,8.01) -- (5.81,7.99) -- (5.79,7.94) -- (5.77,7.87) -- (5.75,7.80) -- (5.74,7.74) -- (5.74,7.73) -- (5.73,7.66) -- (5.73,7.59) -- (5.74,7.56) -- (5.75,7.51) -- (5.78,7.44) -- (5.81,7.40) -- (5.83,7.37) -- (5.88,7.33) -- (5.93,7.30) -- (5.95,7.29);
\draw[line width=0.45pt, draw={rgb,255:red,240;green,123;blue,63}, draw opacity=0.30] (-2.09,5.99) -- (-2.06,6.04) -- (-2.02,6.11) -- (-2.02,6.12) -- (-1.98,6.18) -- (-1.95,6.25) -- (-1.94,6.25) -- (-1.91,6.32) -- (-1.87,6.39) -- (-1.87,6.39) -- (-1.84,6.46) -- (-1.80,6.53) -- (-1.80,6.54) -- (-1.77,6.60) -- (-1.74,6.67) -- (-1.73,6.69) -- (-1.71,6.74) -- (-1.68,6.81) -- (-1.66,6.86) -- (-1.65,6.88) -- (-1.62,6.95) -- (-1.59,7.02) -- (-1.59,7.03) -- (-1.56,7.09) -- (-1.54,7.16) -- (-1.51,7.23) -- (-1.51,7.23) -- (-1.49,7.30) -- (-1.47,7.37) -- (-1.44,7.44) -- (-1.44,7.45) -- (-1.42,7.51) -- (-1.40,7.59) -- (-1.39,7.66) -- (-1.37,7.71) -- (-1.37,7.73) -- (-1.35,7.80) -- (-1.33,7.87) -- (-1.32,7.94) -- (-1.31,8.01) -- (-1.30,8.05) -- (-1.29,8.08) -- (-1.28,8.15) -- (-1.27,8.22) -- (-1.26,8.29) -- (-1.26,8.36) -- (-1.25,8.43) -- (-1.24,8.50) -- (-1.24,8.57) -- (-1.24,8.64) -- (-1.23,8.71) -- (-1.23,8.78) -- (-1.23,8.85) -- (-1.23,8.92) -- (-1.24,8.99) -- (-1.24,9.06) -- (-1.25,9.14) -- (-1.25,9.21) -- (-1.26,9.28) -- (-1.27,9.35) -- (-1.28,9.42) -- (-1.29,9.49) -- (-1.30,9.51) -- (-1.31,9.56) -- (-1.32,9.63) -- (-1.34,9.70) -- (-1.36,9.77) -- (-1.37,9.82) -- (-1.38,9.84) -- (-1.40,9.91) -- (-1.42,9.98) -- (-1.44,10.05) -- (-1.44,10.06) -- (-1.46,10.12) -- (-1.49,10.19) -- (-1.51,10.25) -- (-1.52,10.26) -- (-1.55,10.33) -- (-1.58,10.40) -- (-1.59,10.42) -- (-1.61,10.47) -- (-1.64,10.54) -- (-1.66,10.58) -- (-1.67,10.61) -- (-1.71,10.69) -- (-1.73,10.72) -- (-1.75,10.76) -- (-1.78,10.83) -- (-1.80,10.86) -- (-1.82,10.90) -- (-1.86,10.97) -- (-1.87,10.98) -- (-1.91,11.04) -- (-1.94,11.10) -- (-1.95,11.11) -- (-2.00,11.18) -- (-2.02,11.21) -- (-2.04,11.25) -- (-2.09,11.32) -- (-2.09,11.32);
\draw[line width=0.45pt, draw={rgb,255:red,240;green,123;blue,63}, draw opacity=0.30] (23.69,11.34) -- (23.55,11.25) -- (23.45,11.18) -- (23.33,11.09) -- (23.19,10.99) -- (23.05,10.90) -- (22.95,10.83) -- (22.83,10.74) -- (22.68,10.63) -- (22.56,10.54) -- (22.46,10.47) -- (22.32,10.37) -- (22.18,10.27) -- (22.08,10.19) -- (21.97,10.11) -- (21.82,10.01) -- (21.69,9.91) -- (21.59,9.84) -- (21.46,9.74) -- (21.32,9.64) -- (21.21,9.56) -- (21.10,9.48) -- (20.96,9.38) -- (20.82,9.28) -- (20.71,9.21) -- (20.60,9.13) -- (20.46,9.03) -- (20.31,8.94) -- (20.18,8.85) -- (20.07,8.78) -- (19.95,8.71) -- (19.81,8.62) -- (19.67,8.54) -- (19.52,8.46) -- (19.38,8.38) -- (19.24,8.30) -- (19.09,8.23) -- (18.95,8.16) -- (18.81,8.10) -- (18.66,8.04) -- (18.52,7.98) -- (18.38,7.93) -- (18.20,7.87) -- (18.02,7.81) -- (17.87,7.77) -- (17.71,7.73) -- (17.51,7.68) -- (17.37,7.66) -- (17.15,7.62) -- (16.94,7.59) -- (16.80,7.58) -- (16.58,7.57) -- (16.36,7.56) -- (16.15,7.56) -- (15.93,7.57) -- (15.72,7.58) -- (15.57,7.59) -- (15.36,7.61) -- (15.14,7.64) -- (15.00,7.66) -- (14.78,7.70) -- (14.64,7.73) -- (14.43,7.77) -- (14.28,7.80) -- (14.07,7.84) -- (13.92,7.88) -- (13.71,7.92) -- (13.56,7.96) -- (13.35,8.00) -- (13.21,8.03) -- (12.99,8.07) -- (12.85,8.09) -- (12.63,8.12) -- (12.42,8.14) -- (12.20,8.14) -- (11.98,8.14) -- (11.77,8.12) -- (11.55,8.09) -- (11.41,8.07) -- (11.19,8.02) -- (11.05,7.98) -- (10.91,7.94) -- (10.70,7.87) -- (10.55,7.81) -- (10.41,7.76) -- (10.26,7.70) -- (10.12,7.63) -- (9.97,7.57) -- (9.83,7.50) -- (9.69,7.43) -- (9.54,7.36) -- (9.40,7.29) -- (9.26,7.22) -- (9.11,7.14) -- (8.97,7.07) -- (8.83,6.99) -- (8.68,6.91) -- (8.54,6.83) -- (8.39,6.75) -- (8.25,6.67) -- (8.13,6.60) -- (8.00,6.53) -- (7.88,6.46) -- (7.75,6.38) -- (7.60,6.30) -- (7.46,6.21) -- (7.32,6.13) -- (7.17,6.04) -- (7.05,5.97) -- (6.93,5.89) -- (6.82,5.82) -- (6.67,5.74) -- (6.53,5.65) -- (6.38,5.56) -- (6.24,5.47) -- (6.13,5.40) -- (6.02,5.33) -- (5.88,5.24) -- (5.74,5.15) -- (5.59,5.06) -- (5.47,4.98) -- (5.37,4.91) -- (5.24,4.82) -- (5.09,4.73) -- (4.95,4.63) -- (4.84,4.56) -- (4.73,4.48) -- (4.59,4.38) -- (4.45,4.28) -- (4.34,4.20) -- (4.23,4.13) -- (4.09,4.02) -- (3.95,3.92) -- (3.86,3.85) -- (3.73,3.75) -- (3.58,3.64) -- (3.49,3.57) -- (3.37,3.48) -- (3.22,3.36) -- (3.14,3.29) -- (3.01,3.18) -- (2.88,3.08) -- (2.79,3.00) -- (2.65,2.87) -- (2.56,2.80) -- (2.44,2.68) -- (2.34,2.58) -- (2.22,2.47) -- (2.12,2.37) -- (2.00,2.26) -- (1.91,2.16) -- (1.79,2.04) -- (1.71,1.95) -- (1.58,1.81) -- (1.50,1.72) -- (1.40,1.60) -- (1.29,1.46) -- (1.21,1.38) -- (1.11,1.25) -- (1.00,1.10) -- (0.93,1.00) -- (0.85,0.89) -- (0.75,0.75) -- (0.66,0.61) -- (0.57,0.48) -- (0.50,0.36) -- (0.42,0.24) -- (0.35,0.12) -- (0.27,-0.02) -- (0.19,-0.16) -- (0.12,-0.30) -- (0.04,-0.45) -- (-0.03,-0.59) -- (-0.09,-0.73) -- (-0.16,-0.87) -- (-0.22,-1.01) -- (-0.29,-1.18) -- (-0.37,-1.36) -- (-0.42,-1.50) -- (-0.47,-1.64) -- (-0.52,-1.78) -- (-0.58,-1.95) -- (-0.64,-2.14) -- (-0.68,-2.28) -- (-0.73,-2.42);
\draw[line width=0.45pt, draw={rgb,255:red,240;green,123;blue,63}, draw opacity=0.44] (11.14,-2.49) -- (11.19,-2.48) -- (11.27,-2.47) -- (11.34,-2.45) -- (11.41,-2.44) -- (11.48,-2.42) -- (11.51,-2.42) -- (11.55,-2.41) -- (11.63,-2.39) -- (11.70,-2.38) -- (11.77,-2.36) -- (11.83,-2.35) -- (11.84,-2.35) -- (11.91,-2.33) -- (11.98,-2.31) -- (12.06,-2.29) -- (12.13,-2.28) -- (12.13,-2.28) -- (12.20,-2.26) -- (12.27,-2.24) -- (12.34,-2.22) -- (12.40,-2.21) -- (12.42,-2.20) -- (12.49,-2.18) -- (12.56,-2.16) -- (12.63,-2.14) -- (12.65,-2.14) -- (12.70,-2.12) -- (12.77,-2.10) -- (12.85,-2.08) -- (12.89,-2.07) -- (12.92,-2.06) -- (12.99,-2.04) -- (13.06,-2.01) -- (13.12,-2.00) -- (13.13,-1.99) -- (13.21,-1.97) -- (13.28,-1.94) -- (13.34,-1.92) -- (13.35,-1.92) -- (13.42,-1.90) -- (13.49,-1.87) -- (13.55,-1.85) -- (13.56,-1.85) -- (13.64,-1.82) -- (13.71,-1.80) -- (13.75,-1.78) -- (13.78,-1.77) -- (13.85,-1.75) -- (13.92,-1.72) -- (13.95,-1.71) -- (14.00,-1.70) -- (14.07,-1.67) -- (14.14,-1.65) -- (14.15,-1.64) -- (14.21,-1.62) -- (14.28,-1.59) -- (14.34,-1.57) -- (14.35,-1.57) -- (14.43,-1.54) -- (14.50,-1.51) -- (14.52,-1.50) -- (14.57,-1.48) -- (14.64,-1.46) -- (14.70,-1.43) -- (14.71,-1.43) -- (14.78,-1.40) -- (14.86,-1.37) -- (14.88,-1.36) -- (14.93,-1.34) -- (15.00,-1.31) -- (15.06,-1.29) -- (15.07,-1.29) -- (15.14,-1.26) -- (15.22,-1.23) -- (15.23,-1.22) -- (15.29,-1.20) -- (15.36,-1.17) -- (15.41,-1.15) -- (15.43,-1.14) -- (15.50,-1.11) -- (15.57,-1.08) -- (15.58,-1.08) -- (15.65,-1.05) -- (15.72,-1.02) -- (15.75,-1.01) -- (15.79,-0.99) -- (15.86,-0.96) -- (15.92,-0.94) -- (15.93,-0.93) -- (16.01,-0.90) -- (16.08,-0.87) -- (16.08,-0.87) -- (16.15,-0.84) -- (16.22,-0.81) -- (16.25,-0.80) -- (16.29,-0.78) -- (16.36,-0.75) -- (16.42,-0.73) -- (16.44,-0.72) -- (16.51,-0.69) -- (16.58,-0.66) -- (16.59,-0.66) -- (16.65,-0.63) -- (16.72,-0.60) -- (16.75,-0.59) -- (16.80,-0.57) -- (16.87,-0.54) -- (16.92,-0.52) -- (16.94,-0.51) -- (17.01,-0.48) -- (17.08,-0.45) -- (17.09,-0.45) -- (17.15,-0.42) -- (17.23,-0.39) -- (17.26,-0.37) -- (17.30,-0.36) -- (17.37,-0.33) -- (17.43,-0.30) -- (17.44,-0.30) -- (17.51,-0.27) -- (17.59,-0.24) -- (17.60,-0.23) -- (17.66,-0.21) -- (17.73,-0.18) -- (17.77,-0.16) -- (17.80,-0.15) -- (17.87,-0.12) -- (17.94,-0.10) -- (17.95,-0.09) -- (18.02,-0.07) -- (18.09,-0.04) -- (18.13,-0.02) -- (18.16,-0.01) -- (18.23,0.02) -- (18.30,0.04) -- (18.31,0.05) -- (18.38,0.07) -- (18.45,0.10) -- (18.50,0.12) -- (18.52,0.12) -- (18.59,0.15) -- (18.66,0.18) -- (18.70,0.19) -- (18.73,0.20) -- (18.81,0.23) -- (18.88,0.25) -- (18.90,0.26) -- (18.95,0.28) -- (19.02,0.30) -- (19.09,0.33) -- (19.10,0.33) -- (19.16,0.35) -- (19.24,0.37) -- (19.31,0.40) -- (19.32,0.40) -- (19.38,0.42) -- (19.45,0.44) -- (19.52,0.46) -- (19.54,0.47) -- (19.60,0.49) -- (19.67,0.51) -- (19.74,0.53) -- (19.79,0.54) -- (19.81,0.55) -- (19.88,0.57) -- (19.95,0.59) -- (20.03,0.61) -- (20.04,0.61) -- (20.10,0.63) -- (20.17,0.64) -- (20.24,0.66) -- (20.31,0.68) -- (20.32,0.68) -- (20.39,0.70) -- (20.46,0.71) -- (20.53,0.73) -- (20.60,0.75) -- (20.63,0.75) -- (20.67,0.76) -- (20.74,0.78) -- (20.82,0.79) -- (20.89,0.80) -- (20.96,0.82) -- (20.99,0.82) -- (21.03,0.83) -- (21.10,0.84) -- (21.18,0.86) -- (21.25,0.87) -- (21.32,0.88) -- (21.39,0.89) -- (21.41,0.89) -- (21.46,0.90) -- (21.53,0.91) -- (21.61,0.92) -- (21.68,0.93) -- (21.75,0.94) -- (21.82,0.95) -- (21.89,0.95) -- (21.97,0.96) -- (21.99,0.96) -- (22.04,0.97) -- (22.11,0.97) -- (22.18,0.98) -- (22.25,0.99) -- (22.32,0.99) -- (22.40,1.00) -- (22.47,1.00) -- (22.54,1.01) -- (22.61,1.01) -- (22.68,1.01) -- (22.75,1.02) -- (22.83,1.02) -- (22.90,1.02) -- (22.97,1.02) -- (23.04,1.03) -- (23.11,1.03) -- (23.19,1.03) -- (23.26,1.03) -- (23.33,1.03) -- (23.40,1.03) -- (23.47,1.03) -- (23.54,1.03) -- (23.62,1.03) -- (23.69,1.03);
\draw[line width=0.45pt, draw={rgb,255:red,240;green,123;blue,63}, draw opacity=0.44] (23.69,8.26) -- (23.54,8.21) -- (23.37,8.15) -- (23.19,8.08) -- (23.04,8.03) -- (22.90,7.98) -- (22.75,7.93) -- (22.58,7.87) -- (22.40,7.80) -- (22.25,7.75) -- (22.11,7.71) -- (21.95,7.66) -- (21.75,7.59) -- (21.61,7.55) -- (21.46,7.50) -- (21.25,7.44) -- (21.10,7.40) -- (20.96,7.36) -- (20.74,7.31) -- (20.60,7.27) -- (20.42,7.23) -- (20.24,7.19) -- (20.09,7.16) -- (19.88,7.12) -- (19.70,7.09) -- (19.52,7.06) -- (19.31,7.03) -- (19.16,7.02) -- (18.95,6.99) -- (18.73,6.98) -- (18.52,6.96) -- (18.33,6.95) -- (18.16,6.95) -- (17.94,6.94) -- (17.73,6.94) -- (17.51,6.95) -- (17.35,6.95) -- (17.15,6.96) -- (16.94,6.97) -- (16.72,6.99) -- (16.51,7.01) -- (16.36,7.03) -- (16.15,7.05) -- (15.93,7.08) -- (15.79,7.10) -- (15.57,7.13) -- (15.36,7.16) -- (15.22,7.18) -- (15.00,7.22) -- (14.86,7.25) -- (14.64,7.28) -- (14.50,7.31) -- (14.28,7.35) -- (14.12,7.37) -- (13.92,7.41) -- (13.71,7.44) -- (13.56,7.47) -- (13.35,7.50) -- (13.21,7.52) -- (12.99,7.54) -- (12.77,7.56) -- (12.56,7.57) -- (12.34,7.58) -- (12.13,7.57) -- (11.91,7.56) -- (11.70,7.54) -- (11.51,7.51) -- (11.34,7.49) -- (11.15,7.44) -- (10.98,7.40) -- (10.84,7.36) -- (10.64,7.30) -- (10.48,7.25) -- (10.33,7.19) -- (10.19,7.14) -- (10.05,7.08) -- (9.90,7.02) -- (9.75,6.95) -- (9.60,6.88) -- (9.46,6.81) -- (9.32,6.74) -- (9.18,6.67) -- (9.04,6.60) -- (8.90,6.52) -- (8.75,6.44) -- (8.61,6.35) -- (8.47,6.27) -- (8.32,6.18) -- (8.21,6.11) -- (8.10,6.04) -- (7.96,5.95) -- (7.82,5.85) -- (7.68,5.75) -- (7.58,5.68) -- (7.46,5.60) -- (7.32,5.49) -- (7.20,5.40) -- (7.10,5.32) -- (6.96,5.21) -- (6.85,5.12) -- (6.74,5.02) -- (6.61,4.91) -- (6.53,4.83) -- (6.38,4.70) -- (6.31,4.63) -- (6.17,4.49) -- (6.10,4.41) -- (5.97,4.27) -- (5.88,4.17) -- (5.79,4.06) -- (5.68,3.92) -- (5.59,3.82) -- (5.52,3.71) -- (5.42,3.57) -- (5.32,3.43) -- (5.24,3.29) -- (5.16,3.16) -- (5.09,3.03) -- (5.02,2.88) -- (4.95,2.73) -- (4.89,2.58) -- (4.83,2.44) -- (4.79,2.30) -- (4.73,2.12) -- (4.69,1.95) -- (4.66,1.81) -- (4.63,1.60) -- (4.60,1.39) -- (4.59,1.17) -- (4.60,0.96) -- (4.61,0.75) -- (4.64,0.54) -- (4.67,0.40) -- (4.72,0.19) -- (4.76,0.05) -- (4.81,-0.09) -- (4.88,-0.27) -- (4.95,-0.42) -- (5.02,-0.56) -- (5.09,-0.68) -- (5.17,-0.80) -- (5.27,-0.94) -- (5.38,-1.08) -- (5.45,-1.17) -- (5.57,-1.29) -- (5.67,-1.39) -- (5.79,-1.50) -- (5.88,-1.58) -- (6.03,-1.69) -- (6.16,-1.78) -- (6.26,-1.85) -- (6.38,-1.93) -- (6.53,-2.01) -- (6.67,-2.08) -- (6.82,-2.15) -- (6.96,-2.21) -- (7.12,-2.28) -- (7.32,-2.35) -- (7.46,-2.39) -- (7.60,-2.43) -- (7.82,-2.49);
\draw[line width=0.45pt, draw={rgb,255:red,240;green,123;blue,63}, draw opacity=0.58] (10.12,0.82) -- (10.33,0.81) -- (10.55,0.80) -- (10.76,0.80) -- (10.98,0.81) -- (11.19,0.82) -- (11.34,0.83) -- (11.55,0.85) -- (11.77,0.88) -- (11.91,0.90) -- (12.13,0.93) -- (12.32,0.96) -- (12.49,1.00) -- (12.67,1.03) -- (12.85,1.07) -- (12.99,1.11) -- (13.21,1.16) -- (13.35,1.20) -- (13.51,1.25) -- (13.71,1.30) -- (13.85,1.35) -- (14.00,1.39) -- (14.18,1.46) -- (14.35,1.52) -- (14.50,1.57) -- (14.64,1.62) -- (14.78,1.68) -- (14.93,1.74) -- (15.10,1.81) -- (15.27,1.88) -- (15.43,1.95) -- (15.57,2.02) -- (15.72,2.08) -- (15.86,2.15) -- (16.01,2.22) -- (16.15,2.30) -- (16.29,2.37) -- (16.43,2.44) -- (16.56,2.51) -- (16.68,2.58) -- (16.80,2.65) -- (16.94,2.74) -- (17.08,2.83) -- (17.23,2.92) -- (17.36,3.01) -- (17.46,3.08) -- (17.59,3.17) -- (17.73,3.28) -- (17.83,3.36) -- (17.94,3.45) -- (18.08,3.57) -- (18.16,3.64) -- (18.30,3.78) -- (18.38,3.87) -- (18.48,3.99) -- (18.58,4.13) -- (18.66,4.27) -- (18.73,4.42) -- (18.77,4.56) -- (18.80,4.77) -- (18.77,4.98) -- (18.72,5.12) -- (18.64,5.26) -- (18.53,5.40) -- (18.45,5.48) -- (18.30,5.61) -- (18.20,5.68) -- (18.09,5.76) -- (17.94,5.84) -- (17.80,5.92) -- (17.66,5.99) -- (17.51,6.05) -- (17.37,6.11) -- (17.20,6.18) -- (17.01,6.25) -- (16.87,6.30) -- (16.72,6.34) -- (16.58,6.39) -- (16.36,6.45) -- (16.22,6.49) -- (16.08,6.53) -- (15.86,6.59) -- (15.72,6.63) -- (15.54,6.67) -- (15.36,6.71) -- (15.22,6.74) -- (15.00,6.79) -- (14.86,6.82) -- (14.64,6.86) -- (14.50,6.89) -- (14.28,6.93) -- (14.12,6.95) -- (13.92,6.98) -- (13.71,7.01) -- (13.56,7.03) -- (13.35,7.05) -- (13.13,7.07) -- (12.92,7.08) -- (12.70,7.09) -- (12.49,7.09) -- (12.27,7.08) -- (12.06,7.07) -- (11.84,7.05) -- (11.63,7.02) -- (11.48,7.00) -- (11.27,6.95) -- (11.12,6.92) -- (10.96,6.88) -- (10.76,6.82) -- (10.62,6.78) -- (10.48,6.73) -- (10.32,6.67) -- (10.14,6.60) -- (9.98,6.53) -- (9.83,6.46) -- (9.69,6.39) -- (9.54,6.32) -- (9.41,6.25) -- (9.29,6.18) -- (9.17,6.11) -- (9.04,6.03) -- (8.90,5.94) -- (8.75,5.84) -- (8.63,5.75) -- (8.54,5.68) -- (8.39,5.57) -- (8.27,5.47) -- (8.18,5.39) -- (8.04,5.26) -- (7.96,5.19) -- (7.82,5.05) -- (7.75,4.97) -- (7.63,4.84) -- (7.53,4.71) -- (7.46,4.61) -- (7.37,4.49) -- (7.29,4.34) -- (7.21,4.20) -- (7.14,4.06) -- (7.09,3.92) -- (7.03,3.75) -- (6.99,3.57) -- (6.96,3.38) -- (6.95,3.22) -- (6.96,3.01) -- (6.97,2.87) -- (7.02,2.65) -- (7.07,2.51) -- (7.14,2.37) -- (7.21,2.23) -- (7.30,2.09) -- (7.39,1.98) -- (7.48,1.88) -- (7.60,1.75) -- (7.70,1.67) -- (7.82,1.57) -- (7.96,1.47) -- (8.10,1.39) -- (8.23,1.32) -- (8.37,1.25) -- (8.53,1.17) -- (8.68,1.11) -- (8.83,1.06) -- (8.97,1.02) -- (9.18,0.96) -- (9.33,0.93) -- (9.53,0.89) -- (9.69,0.87) -- (9.90,0.84) -- (10.10,0.82);
\draw[line width=0.45pt, draw={rgb,255:red,240;green,123;blue,63}, draw opacity=0.72] (11.34,2.72) -- (11.48,2.72) -- (11.63,2.72) -- (11.70,2.72) -- (11.84,2.73) -- (11.98,2.74) -- (12.13,2.75) -- (12.27,2.76) -- (12.42,2.78) -- (12.53,2.80) -- (12.63,2.81) -- (12.77,2.83) -- (12.92,2.86) -- (12.99,2.87) -- (13.13,2.90) -- (13.28,2.94) -- (13.35,2.95) -- (13.49,2.99) -- (13.56,3.01) -- (13.71,3.05) -- (13.80,3.08) -- (13.92,3.11) -- (14.02,3.15) -- (14.14,3.19) -- (14.22,3.22) -- (14.35,3.27) -- (14.43,3.30) -- (14.57,3.35) -- (14.64,3.39) -- (14.74,3.43) -- (14.86,3.49) -- (14.93,3.52) -- (15.03,3.57) -- (15.14,3.63) -- (15.22,3.67) -- (15.29,3.72) -- (15.39,3.78) -- (15.50,3.85) -- (15.57,3.90) -- (15.65,3.96) -- (15.72,4.01) -- (15.79,4.07) -- (15.86,4.13) -- (15.94,4.20) -- (16.01,4.27) -- (16.08,4.35) -- (16.15,4.45) -- (16.22,4.55) -- (16.26,4.63) -- (16.30,4.70) -- (16.35,4.84) -- (16.36,4.91) -- (16.38,5.05) -- (16.36,5.19) -- (16.35,5.26) -- (16.29,5.40) -- (16.25,5.47) -- (16.20,5.54) -- (16.15,5.61) -- (16.08,5.69) -- (16.00,5.75) -- (15.91,5.82) -- (15.81,5.89) -- (15.72,5.95) -- (15.65,6.00) -- (15.57,6.04) -- (15.43,6.10) -- (15.36,6.14) -- (15.26,6.18) -- (15.14,6.22) -- (15.07,6.25) -- (14.93,6.29) -- (14.85,6.32) -- (14.71,6.36) -- (14.60,6.39) -- (14.50,6.41) -- (14.35,6.44) -- (14.28,6.46) -- (14.14,6.48) -- (14.00,6.51) -- (13.86,6.53) -- (13.78,6.54) -- (13.64,6.55) -- (13.49,6.57) -- (13.35,6.58) -- (13.21,6.59) -- (13.06,6.59) -- (12.92,6.60) -- (12.77,6.60) -- (12.63,6.59) -- (12.49,6.59) -- (12.34,6.58) -- (12.20,6.56) -- (12.06,6.55) -- (11.92,6.53) -- (11.84,6.52) -- (11.70,6.49) -- (11.55,6.46) -- (11.48,6.45) -- (11.34,6.41) -- (11.24,6.39) -- (11.12,6.35) -- (11.00,6.32) -- (10.91,6.29) -- (10.80,6.25) -- (10.69,6.21) -- (10.62,6.18) -- (10.48,6.12) -- (10.41,6.08) -- (10.30,6.04) -- (10.19,5.98) -- (10.12,5.94) -- (10.03,5.89) -- (9.92,5.82) -- (9.83,5.77) -- (9.76,5.72) -- (9.69,5.67) -- (9.61,5.61) -- (9.52,5.54) -- (9.43,5.47) -- (9.35,5.40) -- (9.28,5.33) -- (9.21,5.26) -- (9.15,5.19) -- (9.09,5.12) -- (9.03,5.05) -- (8.97,4.95) -- (8.90,4.84) -- (8.86,4.77) -- (8.83,4.68) -- (8.78,4.56) -- (8.75,4.43) -- (8.74,4.34) -- (8.74,4.20) -- (8.75,4.06) -- (8.77,3.99) -- (8.81,3.85) -- (8.85,3.78) -- (8.90,3.69) -- (8.97,3.59) -- (9.04,3.50) -- (9.11,3.43) -- (9.18,3.36) -- (9.26,3.31) -- (9.33,3.26) -- (9.40,3.21) -- (9.50,3.15) -- (9.62,3.09) -- (9.69,3.06) -- (9.80,3.01) -- (9.90,2.97) -- (10.00,2.94) -- (10.12,2.90) -- (10.24,2.87) -- (10.33,2.84) -- (10.48,2.81) -- (10.58,2.80) -- (10.69,2.78) -- (10.84,2.76) -- (10.98,2.74) -- (11.12,2.73) -- (11.27,2.73) -- (11.34,2.72);
\draw[line width=0.45pt, draw={rgb,255:red,240;green,123;blue,63}, draw opacity=0.86] (11.55,3.98) -- (11.63,3.97) -- (11.70,3.96) -- (11.77,3.95) -- (11.84,3.95) -- (11.91,3.94) -- (11.98,3.94) -- (12.06,3.94) -- (12.13,3.93) -- (12.20,3.93) -- (12.27,3.93) -- (12.34,3.94) -- (12.42,3.94) -- (12.49,3.94) -- (12.56,3.95) -- (12.63,3.96) -- (12.70,3.96) -- (12.77,3.97) -- (12.85,3.98) -- (12.92,3.99) -- (12.93,3.99) -- (12.99,4.00) -- (13.06,4.01) -- (13.13,4.03) -- (13.21,4.04) -- (13.28,4.06) -- (13.29,4.06) -- (13.35,4.08) -- (13.42,4.09) -- (13.49,4.11) -- (13.56,4.13) -- (13.56,4.13) -- (13.64,4.16) -- (13.71,4.18) -- (13.78,4.20) -- (13.78,4.20) -- (13.85,4.23) -- (13.92,4.26) -- (13.96,4.27) -- (14.00,4.29) -- (14.07,4.32) -- (14.12,4.34) -- (14.14,4.35) -- (14.21,4.39) -- (14.26,4.42) -- (14.28,4.43) -- (14.35,4.47) -- (14.38,4.49) -- (14.43,4.52) -- (14.49,4.56) -- (14.50,4.56) -- (14.57,4.62) -- (14.58,4.63) -- (14.64,4.68) -- (14.66,4.70) -- (14.71,4.75) -- (14.73,4.77) -- (14.78,4.83) -- (14.80,4.84) -- (14.85,4.91) -- (14.86,4.93) -- (14.89,4.98) -- (14.92,5.05) -- (14.93,5.10) -- (14.94,5.12) -- (14.95,5.19) -- (14.94,5.26) -- (14.93,5.33) -- (14.93,5.35) -- (14.91,5.40) -- (14.87,5.47) -- (14.86,5.50) -- (14.83,5.54) -- (14.78,5.59) -- (14.76,5.61) -- (14.71,5.66) -- (14.68,5.68) -- (14.64,5.71) -- (14.58,5.75) -- (14.57,5.76) -- (14.50,5.80) -- (14.45,5.82) -- (14.43,5.84) -- (14.35,5.87) -- (14.29,5.89) -- (14.28,5.90) -- (14.21,5.92) -- (14.14,5.95) -- (14.07,5.97) -- (14.07,5.97) -- (14.00,5.99) -- (13.92,6.00) -- (13.85,6.02) -- (13.78,6.03) -- (13.76,6.04) -- (13.71,6.04) -- (13.64,6.05) -- (13.56,6.06) -- (13.49,6.07) -- (13.42,6.08) -- (13.35,6.08) -- (13.28,6.09) -- (13.21,6.09) -- (13.13,6.09) -- (13.06,6.10) -- (12.99,6.10) -- (12.92,6.09) -- (12.85,6.09) -- (12.77,6.09) -- (12.70,6.08) -- (12.63,6.08) -- (12.56,6.07) -- (12.49,6.06) -- (12.42,6.05) -- (12.34,6.05) -- (12.28,6.04) -- (12.27,6.03) -- (12.20,6.02) -- (12.13,6.01) -- (12.06,5.99) -- (11.98,5.98) -- (11.93,5.97) -- (11.91,5.96) -- (11.84,5.94) -- (11.77,5.92) -- (11.70,5.90) -- (11.68,5.89) -- (11.63,5.88) -- (11.55,5.85) -- (11.48,5.83) -- (11.47,5.82) -- (11.41,5.80) -- (11.34,5.77) -- (11.30,5.75) -- (11.27,5.74) -- (11.19,5.70) -- (11.15,5.68) -- (11.12,5.67) -- (11.05,5.63) -- (11.03,5.61) -- (10.98,5.58) -- (10.91,5.54) -- (10.91,5.54) -- (10.84,5.49) -- (10.81,5.47) -- (10.76,5.43) -- (10.73,5.40) -- (10.69,5.37) -- (10.65,5.33) -- (10.62,5.30) -- (10.59,5.26) -- (10.55,5.21) -- (10.53,5.19) -- (10.49,5.12) -- (10.48,5.10) -- (10.45,5.05) -- (10.42,4.98) -- (10.41,4.92) -- (10.40,4.91) -- (10.39,4.84) -- (10.39,4.77) -- (10.40,4.70) -- (10.41,4.69) -- (10.42,4.63) -- (10.46,4.56) -- (10.48,4.52) -- (10.50,4.49) -- (10.55,4.42) -- (10.56,4.42) -- (10.62,4.35) -- (10.63,4.34) -- (10.69,4.30) -- (10.72,4.27) -- (10.76,4.25) -- (10.84,4.21) -- (10.84,4.20) -- (10.91,4.17) -- (10.98,4.14) -- (10.99,4.13) -- (11.05,4.11) -- (11.12,4.08) -- (11.19,4.06) -- (11.19,4.06) -- (11.27,4.04) -- (11.34,4.02) -- (11.41,4.01) -- (11.48,3.99) -- (11.48,3.99) -- (11.55,3.98);
\draw[line width=0.45pt, draw={rgb,255:red,43;green,192;blue,142}, draw opacity=0.30] (-2.09,14.55) -- (-1.84,14.49) -- (-1.59,14.43) -- (-1.37,14.38) -- (-1.15,14.34) -- (-0.88,14.28) -- (-0.65,14.23) -- (-0.44,14.18) -- (-0.22,14.14) -- (0.07,14.08) -- (0.28,14.04) -- (0.50,13.99) -- (0.78,13.94) -- (1.00,13.90) -- (1.22,13.86) -- (1.50,13.80) -- (1.72,13.77) -- (2.00,13.72) -- (2.22,13.68) -- (2.44,13.64) -- (2.72,13.60) -- (2.94,13.56) -- (3.22,13.52) -- (3.44,13.49) -- (3.73,13.45) -- (3.94,13.42) -- (4.23,13.38) -- (4.45,13.36) -- (4.73,13.32) -- (5.02,13.29) -- (5.24,13.27) -- (5.52,13.24) -- (5.74,13.22) -- (6.03,13.19) -- (6.31,13.17) -- (6.53,13.15) -- (6.82,13.12) -- (7.10,13.10) -- (7.38,13.08) -- (7.60,13.06) -- (7.89,13.04) -- (8.18,13.02) -- (8.39,13.00) -- (8.68,12.97) -- (8.94,12.94) -- (9.18,12.91) -- (9.46,12.87) -- (9.69,12.83) -- (9.90,12.79) -- (10.15,12.73) -- (10.40,12.66) -- (10.59,12.59) -- (10.76,12.51) -- (10.98,12.40) -- (11.13,12.31) -- (11.30,12.16) -- (11.43,12.02) -- (11.55,11.84) -- (11.64,11.67) -- (11.71,11.46) -- (11.77,11.18) -- (11.78,10.97) -- (11.78,10.69) -- (11.77,10.44) -- (11.74,10.19) -- (11.70,9.91) -- (11.67,9.70) -- (11.64,9.42) -- (11.62,9.21) -- (11.60,8.92) -- (11.63,8.64) -- (11.68,8.43) -- (11.77,8.23) -- (11.88,8.08) -- (12.02,7.94) -- (12.20,7.81) -- (12.34,7.72) -- (12.56,7.62) -- (12.77,7.53) -- (12.99,7.45) -- (13.21,7.38) -- (13.42,7.32) -- (13.64,7.26) -- (13.85,7.20) -- (14.07,7.15) -- (14.29,7.09) -- (14.57,7.03) -- (14.78,6.98) -- (15.00,6.93) -- (15.22,6.88) -- (15.50,6.81) -- (15.72,6.76) -- (15.93,6.71) -- (16.15,6.66) -- (16.40,6.60) -- (16.65,6.54) -- (16.87,6.49) -- (17.08,6.43) -- (17.30,6.38) -- (17.55,6.32) -- (17.80,6.25) -- (18.02,6.20) -- (18.23,6.14) -- (18.45,6.08) -- (18.66,6.02) -- (18.88,5.96) -- (19.09,5.89) -- (19.32,5.82) -- (19.54,5.75) -- (19.75,5.68) -- (19.95,5.61) -- (20.17,5.52) -- (20.39,5.43) -- (20.58,5.33) -- (20.74,5.22) -- (20.89,5.10) -- (21.01,4.91) -- (20.98,4.63) -- (20.89,4.46) -- (20.75,4.27) -- (20.61,4.13) -- (20.47,3.99) -- (20.31,3.86) -- (20.17,3.74) -- (20.03,3.62) -- (19.87,3.50) -- (19.69,3.36) -- (19.52,3.23) -- (19.38,3.13) -- (19.21,3.01) -- (19.02,2.87) -- (18.88,2.76) -- (18.72,2.65) -- (18.52,2.51) -- (18.38,2.41) -- (18.23,2.30) -- (18.03,2.16) -- (17.87,2.05) -- (17.73,1.95) -- (17.53,1.81) -- (17.37,1.69) -- (17.23,1.59) -- (17.04,1.46) -- (16.87,1.33) -- (16.72,1.23) -- (16.56,1.10) -- (16.37,0.96) -- (16.22,0.85) -- (16.08,0.74) -- (15.90,0.61) -- (15.72,0.47) -- (15.57,0.35) -- (15.43,0.24) -- (15.28,0.12) -- (15.11,-0.02) -- (14.94,-0.16) -- (14.78,-0.30) -- (14.64,-0.42) -- (14.50,-0.55) -- (14.35,-0.67) -- (14.21,-0.80) -- (14.06,-0.94) -- (13.91,-1.08) -- (13.76,-1.22) -- (13.62,-1.36) -- (13.47,-1.50) -- (13.33,-1.64) -- (13.19,-1.78) -- (13.06,-1.92) -- (12.92,-2.07) -- (12.77,-2.23) -- (12.63,-2.38);
\draw[line width=0.45pt, draw={rgb,255:red,43;green,192;blue,142}, draw opacity=0.30] (23.69,18.69) -- (23.65,18.72) -- (23.62,18.74) -- (23.55,18.79) -- (23.54,18.79) -- (23.47,18.84) -- (23.45,18.86) -- (23.40,18.90) -- (23.37,18.93) -- (23.33,18.96) -- (23.28,19.00) -- (23.26,19.02) -- (23.21,19.07) -- (23.19,19.09) -- (23.13,19.14) -- (23.11,19.16) -- (23.07,19.21) -- (23.04,19.24) -- (23.00,19.28) -- (22.97,19.32) -- (22.95,19.35) -- (22.90,19.41) -- (22.89,19.42) -- (22.84,19.49) -- (22.83,19.51) -- (22.80,19.56) -- (22.75,19.63) -- (22.75,19.63) -- (22.72,19.70) -- (22.68,19.77) -- (22.68,19.77) -- (22.65,19.84) -- (22.62,19.91) -- (22.61,19.93) -- (22.59,19.98) -- (22.57,20.05) -- (22.54,20.12) -- (22.54,20.14) -- (22.52,20.20) -- (22.51,20.27) -- (22.49,20.34) -- (22.48,20.41) -- (22.47,20.46) -- (22.46,20.48) -- (22.45,20.55) -- (22.45,20.62) -- (22.44,20.69);
\draw[line width=0.45pt, draw={rgb,255:red,43;green,192;blue,142}, draw opacity=0.44] (-2.09,20.50) -- (-1.87,20.37) -- (-1.66,20.23) -- (-1.44,20.10) -- (-1.23,19.97) -- (-1.01,19.84) -- (-0.79,19.70) -- (-0.55,19.56) -- (-0.32,19.42) -- (-0.08,19.28) -- (0.14,19.15) -- (0.35,19.02) -- (0.57,18.90) -- (0.78,18.77) -- (1.00,18.65) -- (1.24,18.50) -- (1.49,18.36) -- (1.72,18.24) -- (1.93,18.12) -- (2.15,18.00) -- (2.38,17.87) -- (2.63,17.73) -- (2.87,17.60) -- (3.08,17.49) -- (3.30,17.37) -- (3.56,17.24) -- (3.80,17.11) -- (4.01,17.00) -- (4.25,16.88) -- (4.52,16.75) -- (4.73,16.64) -- (4.97,16.53) -- (5.24,16.40) -- (5.45,16.30) -- (5.72,16.18) -- (5.95,16.07) -- (6.19,15.97) -- (6.46,15.85) -- (6.67,15.76) -- (6.96,15.64) -- (7.18,15.55) -- (7.46,15.43) -- (7.70,15.33) -- (7.96,15.23) -- (8.23,15.12) -- (8.47,15.03) -- (8.75,14.92) -- (8.97,14.84) -- (9.26,14.73) -- (9.50,14.63) -- (9.76,14.53) -- (10.01,14.42) -- (10.26,14.31) -- (10.48,14.21) -- (10.76,14.07) -- (10.98,13.94) -- (11.19,13.80) -- (11.40,13.64) -- (11.57,13.50) -- (11.76,13.29) -- (11.91,13.09) -- (12.04,12.87) -- (12.13,12.66) -- (12.22,12.38) -- (12.27,12.09) -- (12.31,11.74) -- (12.32,11.39) -- (12.32,11.04) -- (12.30,10.69) -- (12.28,10.33) -- (12.26,10.05) -- (12.26,9.70) -- (12.27,9.42) -- (12.33,9.06) -- (12.42,8.81) -- (12.55,8.57) -- (12.70,8.38) -- (12.88,8.22) -- (13.08,8.08) -- (13.33,7.94) -- (13.56,7.83) -- (13.81,7.73) -- (14.07,7.64) -- (14.35,7.55) -- (14.64,7.47) -- (14.93,7.40) -- (15.22,7.34) -- (15.50,7.29) -- (15.86,7.23) -- (16.15,7.20) -- (16.51,7.17) -- (16.80,7.16) -- (17.15,7.16) -- (17.44,7.17) -- (17.80,7.20) -- (18.09,7.24) -- (18.45,7.30) -- (18.73,7.37) -- (18.99,7.44) -- (19.24,7.53) -- (19.52,7.65) -- (19.74,7.74) -- (19.98,7.87) -- (20.23,8.01) -- (20.45,8.15) -- (20.65,8.29) -- (20.83,8.43) -- (21.03,8.60) -- (21.23,8.78) -- (21.39,8.94) -- (21.57,9.14) -- (21.75,9.35) -- (21.89,9.54) -- (22.04,9.76) -- (22.17,9.98) -- (22.28,10.19) -- (22.40,10.44) -- (22.49,10.69) -- (22.57,10.97) -- (22.63,11.25) -- (22.66,11.60) -- (22.65,11.95) -- (22.60,12.23) -- (22.53,12.52) -- (22.42,12.80) -- (22.31,13.01) -- (22.18,13.22) -- (22.02,13.43) -- (21.84,13.64) -- (21.68,13.80) -- (21.46,13.99) -- (21.28,14.14) -- (21.09,14.28) -- (20.88,14.42) -- (20.65,14.56) -- (20.40,14.70) -- (20.17,14.83) -- (19.95,14.94) -- (19.72,15.05) -- (19.45,15.18) -- (19.24,15.28) -- (18.98,15.40) -- (18.73,15.53) -- (18.52,15.63) -- (18.29,15.76) -- (18.04,15.90) -- (17.82,16.04) -- (17.62,16.18) -- (17.44,16.33) -- (17.23,16.53) -- (17.08,16.71) -- (16.94,16.91) -- (16.80,17.17) -- (16.70,17.38) -- (16.61,17.66) -- (16.54,17.94) -- (16.49,18.22) -- (16.45,18.57) -- (16.43,18.86) -- (16.41,19.21) -- (16.41,19.56) -- (16.42,19.91) -- (16.43,20.27) -- (16.44,20.55);
\draw[line width=0.45pt, draw={rgb,255:red,43;green,192;blue,142}, draw opacity=0.58] (15.36,7.80) -- (15.50,7.78) -- (15.65,7.76) -- (15.79,7.75) -- (15.93,7.74) -- (16.08,7.74) -- (16.22,7.74) -- (16.36,7.74) -- (16.51,7.75) -- (16.65,7.77) -- (16.80,7.79) -- (16.87,7.80) -- (17.01,7.82) -- (17.15,7.86) -- (17.23,7.87) -- (17.37,7.92) -- (17.44,7.94) -- (17.59,7.99) -- (17.66,8.02) -- (17.80,8.07) -- (17.87,8.11) -- (17.96,8.15) -- (18.09,8.22) -- (18.16,8.26) -- (18.23,8.30) -- (18.33,8.36) -- (18.43,8.43) -- (18.52,8.49) -- (18.59,8.55) -- (18.66,8.61) -- (18.73,8.67) -- (18.81,8.73) -- (18.88,8.80) -- (18.95,8.88) -- (19.02,8.95) -- (19.09,9.04) -- (19.16,9.13) -- (19.23,9.21) -- (19.28,9.28) -- (19.32,9.35) -- (19.38,9.44) -- (19.45,9.56) -- (19.49,9.63) -- (19.52,9.70) -- (19.59,9.84) -- (19.62,9.91) -- (19.67,10.05) -- (19.69,10.12) -- (19.73,10.26) -- (19.74,10.33) -- (19.77,10.47) -- (19.79,10.61) -- (19.80,10.76) -- (19.80,10.90) -- (19.79,11.04) -- (19.78,11.18) -- (19.75,11.32) -- (19.74,11.39) -- (19.70,11.53) -- (19.67,11.64) -- (19.63,11.74) -- (19.60,11.82) -- (19.53,11.95) -- (19.50,12.02) -- (19.45,12.10) -- (19.38,12.22) -- (19.32,12.31) -- (19.27,12.38) -- (19.21,12.45) -- (19.15,12.52) -- (19.09,12.59) -- (19.02,12.66) -- (18.95,12.73) -- (18.88,12.80) -- (18.81,12.87) -- (18.72,12.94) -- (18.64,13.01) -- (18.55,13.08) -- (18.45,13.15) -- (18.38,13.20) -- (18.30,13.25) -- (18.23,13.30) -- (18.13,13.36) -- (18.02,13.43) -- (17.94,13.47) -- (17.87,13.51) -- (17.75,13.57) -- (17.66,13.62) -- (17.59,13.65) -- (17.45,13.71) -- (17.37,13.75) -- (17.29,13.78) -- (17.15,13.84) -- (17.08,13.87) -- (16.94,13.92) -- (16.87,13.94) -- (16.72,13.99) -- (16.65,14.02) -- (16.51,14.06) -- (16.44,14.08) -- (16.29,14.12) -- (16.22,14.14) -- (16.08,14.17) -- (15.93,14.20) -- (15.86,14.21) -- (15.72,14.24) -- (15.57,14.26) -- (15.43,14.28) -- (15.36,14.28) -- (15.22,14.29) -- (15.07,14.30) -- (14.93,14.29) -- (14.78,14.29) -- (14.71,14.28) -- (14.57,14.26) -- (14.43,14.23) -- (14.35,14.21) -- (14.21,14.16) -- (14.14,14.13) -- (14.01,14.07) -- (13.92,14.01) -- (13.85,13.96) -- (13.78,13.91) -- (13.71,13.84) -- (13.64,13.77) -- (13.56,13.69) -- (13.49,13.59) -- (13.43,13.50) -- (13.39,13.43) -- (13.35,13.36) -- (13.28,13.22) -- (13.25,13.15) -- (13.21,13.04) -- (13.17,12.94) -- (13.13,12.84) -- (13.10,12.73) -- (13.06,12.59) -- (13.04,12.52) -- (13.01,12.38) -- (12.99,12.30) -- (12.96,12.16) -- (12.94,12.02) -- (12.92,11.92) -- (12.90,11.81) -- (12.88,11.67) -- (12.86,11.53) -- (12.85,11.42) -- (12.83,11.32) -- (12.82,11.18) -- (12.81,11.04) -- (12.80,10.90) -- (12.79,10.76) -- (12.79,10.61) -- (12.78,10.47) -- (12.78,10.33) -- (12.79,10.19) -- (12.80,10.05) -- (12.81,9.91) -- (12.83,9.77) -- (12.85,9.68) -- (12.87,9.56) -- (12.91,9.42) -- (12.94,9.35) -- (12.99,9.22) -- (13.03,9.14) -- (13.07,9.06) -- (13.13,8.96) -- (13.21,8.86) -- (13.26,8.78) -- (13.33,8.71) -- (13.40,8.64) -- (13.47,8.57) -- (13.56,8.50) -- (13.64,8.44) -- (13.71,8.39) -- (13.78,8.35) -- (13.88,8.29) -- (14.00,8.23) -- (14.07,8.19) -- (14.15,8.15) -- (14.28,8.09) -- (14.35,8.06) -- (14.50,8.01) -- (14.57,7.99) -- (14.71,7.94) -- (14.78,7.92) -- (14.93,7.88) -- (15.00,7.87) -- (15.14,7.84) -- (15.29,7.81) -- (15.36,7.80);
\draw[line width=0.45pt, draw={rgb,255:red,43;green,192;blue,142}, draw opacity=0.72] (15.22,8.42) -- (15.29,8.41) -- (15.36,8.40) -- (15.43,8.39) -- (15.50,8.38) -- (15.57,8.37) -- (15.65,8.37) -- (15.72,8.37) -- (15.79,8.37) -- (15.86,8.37) -- (15.93,8.38) -- (16.01,8.38) -- (16.08,8.39) -- (16.15,8.40) -- (16.22,8.41) -- (16.29,8.43) -- (16.30,8.43) -- (16.36,8.45) -- (16.44,8.47) -- (16.51,8.49) -- (16.55,8.50) -- (16.58,8.51) -- (16.65,8.54) -- (16.72,8.57) -- (16.73,8.57) -- (16.80,8.60) -- (16.87,8.64) -- (16.87,8.64) -- (16.94,8.68) -- (17.00,8.71) -- (17.01,8.72) -- (17.08,8.77) -- (17.10,8.78) -- (17.15,8.82) -- (17.20,8.85) -- (17.23,8.88) -- (17.28,8.92) -- (17.30,8.94) -- (17.36,8.99) -- (17.37,9.01) -- (17.43,9.06) -- (17.44,9.08) -- (17.49,9.14) -- (17.51,9.16) -- (17.55,9.21) -- (17.59,9.26) -- (17.60,9.28) -- (17.65,9.35) -- (17.66,9.36) -- (17.69,9.42) -- (17.73,9.48) -- (17.73,9.49) -- (17.77,9.56) -- (17.80,9.63) -- (17.80,9.63) -- (17.83,9.70) -- (17.86,9.77) -- (17.87,9.81) -- (17.88,9.84) -- (17.90,9.91) -- (17.92,9.98) -- (17.93,10.05) -- (17.94,10.10) -- (17.95,10.12) -- (17.96,10.19) -- (17.96,10.26) -- (17.97,10.33) -- (17.97,10.40) -- (17.97,10.47) -- (17.97,10.54) -- (17.96,10.61) -- (17.95,10.69) -- (17.94,10.75) -- (17.94,10.76) -- (17.93,10.83) -- (17.92,10.90) -- (17.90,10.97) -- (17.88,11.04) -- (17.87,11.05) -- (17.85,11.11) -- (17.83,11.18) -- (17.80,11.24) -- (17.80,11.25) -- (17.77,11.32) -- (17.73,11.39) -- (17.73,11.40) -- (17.69,11.46) -- (17.66,11.52) -- (17.65,11.53) -- (17.61,11.60) -- (17.59,11.64) -- (17.56,11.67) -- (17.51,11.74) -- (17.51,11.74) -- (17.45,11.81) -- (17.44,11.83) -- (17.40,11.88) -- (17.37,11.91) -- (17.33,11.95) -- (17.30,11.99) -- (17.26,12.02) -- (17.23,12.06) -- (17.19,12.09) -- (17.15,12.13) -- (17.11,12.16) -- (17.08,12.19) -- (17.03,12.23) -- (17.01,12.25) -- (16.94,12.30) -- (16.94,12.31) -- (16.87,12.36) -- (16.84,12.38) -- (16.80,12.41) -- (16.73,12.45) -- (16.72,12.45) -- (16.65,12.50) -- (16.62,12.52) -- (16.58,12.54) -- (16.51,12.58) -- (16.49,12.59) -- (16.44,12.61) -- (16.36,12.65) -- (16.34,12.66) -- (16.29,12.68) -- (16.22,12.71) -- (16.17,12.73) -- (16.15,12.74) -- (16.08,12.76) -- (16.01,12.79) -- (15.96,12.80) -- (15.93,12.81) -- (15.86,12.83) -- (15.79,12.84) -- (15.72,12.86) -- (15.67,12.87) -- (15.65,12.87) -- (15.57,12.89) -- (15.50,12.90) -- (15.43,12.90) -- (15.36,12.91) -- (15.29,12.91) -- (15.22,12.91) -- (15.14,12.91) -- (15.07,12.91) -- (15.00,12.90) -- (14.93,12.90) -- (14.86,12.88) -- (14.78,12.87) -- (14.78,12.87) -- (14.71,12.85) -- (14.64,12.83) -- (14.57,12.81) -- (14.54,12.80) -- (14.50,12.78) -- (14.43,12.75) -- (14.38,12.73) -- (14.35,12.72) -- (14.28,12.68) -- (14.25,12.66) -- (14.21,12.63) -- (14.15,12.59) -- (14.14,12.58) -- (14.07,12.53) -- (14.06,12.52) -- (14.00,12.46) -- (13.98,12.45) -- (13.92,12.39) -- (13.91,12.38) -- (13.85,12.31) -- (13.85,12.31) -- (13.79,12.23) -- (13.78,12.22) -- (13.74,12.16) -- (13.71,12.11) -- (13.70,12.09) -- (13.65,12.02) -- (13.64,11.99) -- (13.62,11.95) -- (13.58,11.88) -- (13.56,11.84) -- (13.55,11.81) -- (13.52,11.74) -- (13.50,11.67) -- (13.49,11.66) -- (13.47,11.60) -- (13.45,11.53) -- (13.43,11.46) -- (13.42,11.42) -- (13.41,11.39) -- (13.40,11.32) -- (13.38,11.25) -- (13.37,11.18) -- (13.36,11.11) -- (13.35,11.04) -- (13.35,11.01) -- (13.34,10.97) -- (13.34,10.90) -- (13.34,10.83) -- (13.33,10.76) -- (13.33,10.69) -- (13.33,10.61) -- (13.34,10.54) -- (13.34,10.47) -- (13.35,10.40) -- (13.35,10.38) -- (13.35,10.33) -- (13.36,10.26) -- (13.38,10.19) -- (13.39,10.12) -- (13.41,10.05) -- (13.42,9.99) -- (13.42,9.98) -- (13.44,9.91) -- (13.47,9.84) -- (13.49,9.77) -- (13.49,9.77) -- (13.52,9.70) -- (13.55,9.63) -- (13.56,9.60) -- (13.59,9.56) -- (13.62,9.49) -- (13.64,9.46) -- (13.66,9.42) -- (13.71,9.35) -- (13.71,9.35) -- (13.76,9.28) -- (13.78,9.25) -- (13.82,9.21) -- (13.85,9.17) -- (13.88,9.14) -- (13.92,9.09) -- (13.95,9.06) -- (14.00,9.02) -- (14.02,8.99) -- (14.07,8.95) -- (14.10,8.92) -- (14.14,8.90) -- (14.19,8.85) -- (14.21,8.84) -- (14.28,8.79) -- (14.30,8.78) -- (14.35,8.75) -- (14.41,8.71) -- (14.43,8.71) -- (14.50,8.67) -- (14.55,8.64) -- (14.57,8.63) -- (14.64,8.60) -- (14.70,8.57) -- (14.71,8.57) -- (14.78,8.54) -- (14.86,8.51) -- (14.90,8.50) -- (14.93,8.49) -- (15.00,8.47) -- (15.07,8.45) -- (15.14,8.43) -- (15.16,8.43) -- (15.22,8.42);
\draw[line width=0.45pt, draw={rgb,255:red,43;green,192;blue,142}, draw opacity=0.86] (15.36,9.14) -- (15.43,9.13) -- (15.50,9.13) -- (15.57,9.13) -- (15.60,9.14) -- (15.65,9.14) -- (15.72,9.15) -- (15.79,9.17) -- (15.86,9.19) -- (15.91,9.21) -- (15.93,9.21) -- (16.01,9.24) -- (16.07,9.28) -- (16.08,9.28) -- (16.15,9.32) -- (16.18,9.35) -- (16.22,9.38) -- (16.27,9.42) -- (16.29,9.44) -- (16.34,9.49) -- (16.36,9.51) -- (16.40,9.56) -- (16.44,9.60) -- (16.46,9.63) -- (16.50,9.70) -- (16.51,9.71) -- (16.54,9.77) -- (16.57,9.84) -- (16.58,9.87) -- (16.60,9.91) -- (16.62,9.98) -- (16.63,10.05) -- (16.65,10.12) -- (16.65,10.16) -- (16.66,10.19) -- (16.66,10.26) -- (16.66,10.33) -- (16.66,10.40) -- (16.65,10.47) -- (16.65,10.48) -- (16.64,10.54) -- (16.63,10.61) -- (16.61,10.69) -- (16.59,10.76) -- (16.58,10.80) -- (16.57,10.83) -- (16.54,10.90) -- (16.51,10.97) -- (16.51,10.98) -- (16.48,11.04) -- (16.44,11.11) -- (16.44,11.11) -- (16.39,11.18) -- (16.36,11.22) -- (16.35,11.25) -- (16.29,11.32) -- (16.29,11.32) -- (16.23,11.39) -- (16.22,11.40) -- (16.17,11.46) -- (16.15,11.48) -- (16.09,11.53) -- (16.08,11.54) -- (16.01,11.60) -- (16.01,11.60) -- (15.93,11.66) -- (15.92,11.67) -- (15.86,11.71) -- (15.81,11.74) -- (15.79,11.75) -- (15.72,11.79) -- (15.67,11.81) -- (15.65,11.83) -- (15.57,11.86) -- (15.50,11.88) -- (15.50,11.88) -- (15.43,11.91) -- (15.36,11.92) -- (15.29,11.94) -- (15.22,11.95) -- (15.20,11.95) -- (15.14,11.96) -- (15.07,11.96) -- (15.00,11.96) -- (14.93,11.96) -- (14.90,11.95) -- (14.86,11.95) -- (14.78,11.93) -- (14.71,11.91) -- (14.64,11.89) -- (14.62,11.88) -- (14.57,11.86) -- (14.50,11.82) -- (14.48,11.81) -- (14.43,11.78) -- (14.38,11.74) -- (14.35,11.72) -- (14.30,11.67) -- (14.28,11.66) -- (14.23,11.60) -- (14.21,11.58) -- (14.17,11.53) -- (14.14,11.48) -- (14.13,11.46) -- (14.09,11.39) -- (14.07,11.35) -- (14.05,11.32) -- (14.02,11.25) -- (14.00,11.18) -- (14.00,11.17) -- (13.98,11.11) -- (13.96,11.04) -- (13.95,10.97) -- (13.94,10.90) -- (13.94,10.83) -- (13.94,10.76) -- (13.94,10.69) -- (13.94,10.61) -- (13.95,10.54) -- (13.96,10.47) -- (13.97,10.40) -- (13.99,10.33) -- (14.00,10.31) -- (14.01,10.26) -- (14.03,10.19) -- (14.06,10.12) -- (14.07,10.11) -- (14.09,10.05) -- (14.13,9.98) -- (14.14,9.96) -- (14.17,9.91) -- (14.21,9.84) -- (14.21,9.84) -- (14.26,9.77) -- (14.28,9.74) -- (14.31,9.70) -- (14.35,9.65) -- (14.38,9.63) -- (14.43,9.58) -- (14.45,9.56) -- (14.50,9.51) -- (14.53,9.49) -- (14.57,9.45) -- (14.62,9.42) -- (14.64,9.40) -- (14.71,9.35) -- (14.72,9.35) -- (14.78,9.31) -- (14.85,9.28) -- (14.86,9.27) -- (14.93,9.24) -- (15.00,9.21) -- (15.03,9.21) -- (15.07,9.19) -- (15.14,9.17) -- (15.22,9.16) -- (15.29,9.14) -- (15.36,9.14) -- (15.36,9.14);
\fill[fill={rgb,255:red,58;green,154;blue,216}, fill opacity=0.94] (7.92,9.42) circle (0.094);
\fill[fill={rgb,255:red,58;green,154;blue,216}, fill opacity=0.92] (9.35,9.00) circle (0.092);
\fill[fill={rgb,255:red,58;green,154;blue,216}, fill opacity=0.90] (9.92,9.27) circle (0.089);
\fill[fill={rgb,255:red,59;green,154;blue,215}, fill opacity=0.89] (7.19,10.14) circle (0.087);
\fill[fill={rgb,255:red,58;green,154;blue,216}, fill opacity=0.92] (9.82,9.64) circle (0.091);
\fill[fill={rgb,255:red,58;green,153;blue,216}, fill opacity=0.93] (6.51,8.52) circle (0.092);
\fill[fill={rgb,255:red,58;green,153;blue,216}, fill opacity=0.93] (8.30,8.24) circle (0.092);
\fill[fill={rgb,255:red,58;green,154;blue,216}, fill opacity=0.92] (7.79,9.80) circle (0.092);
\fill[fill={rgb,255:red,52;green,170;blue,183}, fill opacity=0.75] (12.14,11.41) circle (0.070);
\fill[fill={rgb,255:red,43;green,191;blue,142}, fill opacity=0.91] (14.17,11.82) circle (0.090);
\fill[fill={rgb,255:red,45;green,187;blue,150}, fill opacity=0.84] (13.42,12.55) circle (0.081);
\fill[fill={rgb,255:red,55;green,163;blue,198}, fill opacity=0.79] (11.52,11.84) circle (0.075);
\fill[fill={rgb,255:red,58;green,153;blue,216}, fill opacity=0.95] (7.75,8.73) circle (0.095);
\fill[fill={rgb,255:red,45;green,187;blue,149}, fill opacity=0.84] (15.50,13.61) circle (0.081);
\fill[fill={rgb,255:red,58;green,154;blue,216}, fill opacity=0.94] (8.17,9.37) circle (0.094);
\fill[fill={rgb,255:red,47;green,181;blue,162}, fill opacity=0.79] (12.68,11.75) circle (0.075);
\fill[fill={rgb,255:red,57;green,157;blue,208}, fill opacity=0.83] (11.30,10.42) circle (0.080);
\fill[fill={rgb,255:red,61;green,153;blue,213}, fill opacity=0.89] (5.53,6.49) circle (0.087);
\fill[fill={rgb,255:red,58;green,154;blue,216}, fill opacity=0.91] (7.70,9.95) circle (0.090);
\fill[fill={rgb,255:red,58;green,154;blue,216}, fill opacity=0.92] (9.80,10.22) circle (0.091);
\fill[fill={rgb,255:red,58;green,153;blue,216}, fill opacity=0.95] (7.74,9.02) circle (0.095);
\fill[fill={rgb,255:red,55;green,161;blue,200}, fill opacity=0.80] (11.51,11.61) circle (0.077);
\fill[fill={rgb,255:red,81;green,150;blue,196}, fill opacity=0.83] (5.09,5.32) circle (0.079);
\fill[fill={rgb,255:red,58;green,153;blue,216}, fill opacity=0.91] (5.60,7.62) circle (0.091);
\fill[fill={rgb,255:red,58;green,157;blue,206}, fill opacity=0.82] (8.81,12.02) circle (0.079);
\fill[fill={rgb,255:red,62;green,153;blue,213}, fill opacity=0.89] (6.59,6.75) circle (0.087);
\fill[fill={rgb,255:red,58;green,154;blue,216}, fill opacity=0.93] (7.56,9.50) circle (0.092);
\fill[fill={rgb,255:red,58;green,154;blue,216}, fill opacity=0.92] (8.25,10.09) circle (0.091);
\fill[fill={rgb,255:red,58;green,154;blue,216}, fill opacity=0.92] (7.91,9.85) circle (0.092);
\fill[fill={rgb,255:red,66;green,152;blue,209}, fill opacity=0.86] (4.39,5.87) circle (0.084);
\fill[fill={rgb,255:red,58;green,153;blue,216}, fill opacity=0.95] (7.77,8.42) circle (0.095);
\fill[fill={rgb,255:red,58;green,153;blue,216}, fill opacity=0.91] (8.38,7.99) circle (0.090);
\fill[fill={rgb,255:red,58;green,154;blue,216}, fill opacity=0.92] (9.83,10.14) circle (0.091);
\fill[fill={rgb,255:red,61;green,153;blue,213}, fill opacity=0.88] (4.92,8.37) circle (0.087);
\fill[fill={rgb,255:red,51;green,173;blue,177}, fill opacity=0.74] (12.26,11.10) circle (0.069);
\fill[fill={rgb,255:red,90;green,148;blue,189}, fill opacity=0.81] (7.35,6.34) circle (0.077);
\fill[fill={rgb,255:red,58;green,154;blue,216}, fill opacity=0.93] (9.45,10.00) circle (0.093);
\fill[fill={rgb,255:red,62;green,153;blue,212}, fill opacity=0.88] (4.16,7.31) circle (0.086);
\fill[fill={rgb,255:red,58;green,154;blue,216}, fill opacity=0.94] (9.22,9.52) circle (0.094);
\fill[fill={rgb,255:red,59;green,154;blue,213}, fill opacity=0.87] (7.97,10.88) circle (0.085);
\fill[fill={rgb,255:red,58;green,154;blue,216}, fill opacity=0.93] (9.47,9.77) circle (0.093);
\fill[fill={rgb,255:red,58;green,153;blue,216}, fill opacity=0.94] (7.16,8.41) circle (0.094);
\fill[fill={rgb,255:red,58;green,153;blue,216}, fill opacity=0.91] (9.26,8.67) circle (0.090);
\fill[fill={rgb,255:red,58;green,153;blue,216}, fill opacity=0.95] (8.36,8.85) circle (0.095);
\fill[fill={rgb,255:red,61;green,153;blue,213}, fill opacity=0.88] (4.44,7.46) circle (0.087);
\fill[fill={rgb,255:red,58;green,154;blue,216}, fill opacity=0.92] (7.79,9.76) circle (0.092);
\fill[fill={rgb,255:red,58;green,154;blue,215}, fill opacity=0.89] (8.48,10.68) circle (0.088);
\fill[fill={rgb,255:red,58;green,154;blue,216}, fill opacity=0.93] (7.79,9.66) circle (0.092);
\fill[fill={rgb,255:red,54;green,164;blue,196}, fill opacity=0.79] (11.55,11.95) circle (0.075);
\fill[fill={rgb,255:red,43;green,191;blue,142}, fill opacity=0.92] (13.97,10.81) circle (0.091);
\fill[fill={rgb,255:red,58;green,154;blue,216}, fill opacity=0.92] (9.44,9.05) circle (0.091);
\fill[fill={rgb,255:red,62;green,154;blue,210}, fill opacity=0.84] (6.40,10.66) circle (0.082);
\fill[fill={rgb,255:red,58;green,155;blue,214}, fill opacity=0.87] (10.78,10.29) circle (0.086);
\fill[fill={rgb,255:red,58;green,153;blue,216}, fill opacity=0.93] (7.24,9.18) circle (0.093);
\fill[fill={rgb,255:red,58;green,154;blue,216}, fill opacity=0.95] (8.70,9.36) circle (0.095);
\fill[fill={rgb,255:red,58;green,153;blue,216}, fill opacity=0.91] (5.82,8.24) circle (0.091);
\fill[fill={rgb,255:red,57;green,156;blue,210}, fill opacity=0.85] (11.05,11.11) circle (0.082);
\fill[fill={rgb,255:red,58;green,153;blue,216}, fill opacity=0.92] (8.57,8.27) circle (0.091);
\fill[fill={rgb,255:red,58;green,153;blue,216}, fill opacity=0.93] (7.82,8.00) circle (0.093);
\fill[fill={rgb,255:red,59;green,153;blue,215}, fill opacity=0.91] (5.24,7.46) circle (0.090);
\fill[fill={rgb,255:red,58;green,153;blue,216}, fill opacity=0.95] (8.16,8.54) circle (0.094);
\fill[fill={rgb,255:red,76;green,151;blue,199}, fill opacity=0.81] (9.52,7.81) circle (0.078);
\fill[fill={rgb,255:red,56;green,158;blue,206}, fill opacity=0.82] (11.09,11.66) circle (0.079);
\fill[fill={rgb,255:red,46;green,184;blue,156}, fill opacity=0.81] (13.64,13.97) circle (0.077);
\fill[fill={rgb,255:red,48;green,180;blue,164}, fill opacity=0.78] (12.61,13.16) circle (0.074);
\fill[fill={rgb,255:red,57;green,155;blue,213}, fill opacity=0.86] (10.05,11.44) circle (0.084);
\fill[fill={rgb,255:red,62;green,153;blue,212}, fill opacity=0.88] (8.56,7.70) circle (0.086);
\fill[fill={rgb,255:red,66;green,152;blue,209}, fill opacity=0.86] (3.77,7.93) circle (0.083);
\fill[fill={rgb,255:red,58;green,153;blue,216}, fill opacity=0.95] (8.43,9.09) circle (0.095);
\fill[fill={rgb,255:red,58;green,153;blue,216}, fill opacity=0.92] (9.07,8.71) circle (0.092);
\fill[fill={rgb,255:red,61;green,153;blue,214}, fill opacity=0.89] (5.28,6.63) circle (0.088);
\fill[fill={rgb,255:red,58;green,153;blue,216}, fill opacity=0.94] (8.67,8.65) circle (0.093);
\fill[fill={rgb,255:red,59;green,154;blue,212}, fill opacity=0.86] (10.56,9.37) circle (0.084);
\fill[fill={rgb,255:red,61;green,153;blue,213}, fill opacity=0.89] (4.62,6.63) circle (0.087);
\fill[fill={rgb,255:red,59;green,153;blue,215}, fill opacity=0.91] (5.35,7.33) circle (0.090);
\fill[fill={rgb,255:red,58;green,154;blue,216}, fill opacity=0.91] (10.23,10.26) circle (0.090);
\fill[fill={rgb,255:red,58;green,154;blue,216}, fill opacity=0.92] (9.67,10.27) circle (0.091);
\fill[fill={rgb,255:red,58;green,154;blue,216}, fill opacity=0.94] (9.31,9.67) circle (0.093);
\fill[fill={rgb,255:red,49;green,177;blue,170}, fill opacity=0.76] (12.43,11.20) circle (0.072);
\fill[fill={rgb,255:red,70;green,151;blue,205}, fill opacity=0.85] (3.78,5.46) circle (0.083);
\fill[fill={rgb,255:red,57;green,157;blue,209}, fill opacity=0.84] (10.88,11.49) circle (0.081);
\fill[fill={rgb,255:red,79;green,150;blue,198}, fill opacity=0.83] (4.91,5.31) circle (0.080);
\fill[fill={rgb,255:red,58;green,154;blue,216}, fill opacity=0.94] (9.02,9.08) circle (0.094);
\fill[fill={rgb,255:red,48;green,180;blue,163}, fill opacity=0.78] (12.67,13.17) circle (0.074);
\fill[fill={rgb,255:red,58;green,154;blue,216}, fill opacity=0.95] (8.12,9.29) circle (0.094);
\fill[fill={rgb,255:red,58;green,154;blue,215}, fill opacity=0.89] (9.97,11.00) circle (0.087);
\fill[fill={rgb,255:red,58;green,154;blue,216}, fill opacity=0.92] (9.69,10.11) circle (0.092);
\fill[fill={rgb,255:red,55;green,166;blue,190}, fill opacity=0.77] (11.89,10.35) circle (0.072);
\fill[fill={rgb,255:red,63;green,153;blue,211}, fill opacity=0.88] (4.59,6.20) circle (0.086);
\fill[fill={rgb,255:red,58;green,153;blue,216}, fill opacity=0.94] (7.28,8.59) circle (0.094);
\fill[fill={rgb,255:red,58;green,154;blue,216}, fill opacity=0.93] (9.61,9.63) circle (0.092);
\fill[fill={rgb,255:red,58;green,153;blue,216}, fill opacity=0.93] (6.47,7.73) circle (0.093);
\fill[fill={rgb,255:red,58;green,154;blue,216}, fill opacity=0.94] (8.51,9.49) circle (0.094);
\fill[fill={rgb,255:red,61;green,153;blue,213}, fill opacity=0.88] (4.39,7.46) circle (0.087);
\fill[fill={rgb,255:red,58;green,153;blue,216}, fill opacity=0.94] (7.32,7.99) circle (0.094);
\fill[fill={rgb,255:red,53;green,168;blue,187}, fill opacity=0.76] (11.48,12.61) circle (0.072);
\fill[fill={rgb,255:red,59;green,153;blue,215}, fill opacity=0.91] (5.44,7.21) circle (0.090);
\fill[fill={rgb,255:red,58;green,153;blue,216}, fill opacity=0.91] (6.13,8.76) circle (0.090);
\fill[fill={rgb,255:red,60;green,153;blue,214}, fill opacity=0.89] (5.24,8.32) circle (0.088);
\fill[fill={rgb,255:red,58;green,155;blue,212}, fill opacity=0.86] (8.55,11.27) circle (0.084);
\fill[fill={rgb,255:red,55;green,163;blue,197}, fill opacity=0.79] (11.14,12.28) circle (0.075);
\fill[fill={rgb,255:red,58;green,153;blue,216}, fill opacity=0.91] (6.22,8.93) circle (0.090);
\fill[fill={rgb,255:red,76;green,151;blue,200}, fill opacity=0.83] (1.97,6.75) circle (0.080);
\fill[fill={rgb,255:red,52;green,175;blue,170}, fill opacity=0.75] (12.25,10.04) circle (0.070);
\fill[fill={rgb,255:red,60;green,153;blue,214}, fill opacity=0.89] (6.34,6.80) circle (0.088);
\fill[fill={rgb,255:red,58;green,154;blue,216}, fill opacity=0.92] (7.67,9.88) circle (0.091);
\fill[fill={rgb,255:red,43;green,191;blue,142}, fill opacity=0.91] (13.89,10.67) circle (0.091);
\fill[fill={rgb,255:red,58;green,153;blue,216}, fill opacity=0.94] (7.33,8.72) circle (0.094);
\fill[fill={rgb,255:red,58;green,154;blue,216}, fill opacity=0.94] (8.76,9.69) circle (0.094);
\fill[fill={rgb,255:red,59;green,159;blue,201}, fill opacity=0.80] (11.35,9.63) circle (0.076);
\fill[fill={rgb,255:red,58;green,153;blue,216}, fill opacity=0.92] (9.09,8.73) circle (0.092);
\fill[fill={rgb,255:red,49;green,178;blue,168}, fill opacity=0.77] (12.46,12.76) circle (0.073);
\fill[fill={rgb,255:red,58;green,154;blue,216}, fill opacity=0.93] (8.04,9.63) circle (0.093);
\fill[fill={rgb,255:red,57;green,157;blue,209}, fill opacity=0.84] (11.25,10.56) circle (0.081);
\fill[fill={rgb,255:red,57;green,158;blue,206}, fill opacity=0.82] (11.40,10.40) circle (0.079);
\fill[fill={rgb,255:red,58;green,153;blue,216}, fill opacity=0.94] (7.97,8.18) circle (0.093);
\fill[fill={rgb,255:red,58;green,153;blue,216}, fill opacity=0.93] (6.90,8.77) circle (0.093);
\fill[fill={rgb,255:red,87;green,149;blue,191}, fill opacity=0.81] (3.13,4.16) circle (0.078);
\fill[fill={rgb,255:red,58;green,153;blue,216}, fill opacity=0.91] (5.63,7.40) circle (0.091);
\fill[fill={rgb,255:red,60;green,153;blue,214}, fill opacity=0.88] (8.92,8.05) circle (0.087);
\fill[fill={rgb,255:red,58;green,154;blue,216}, fill opacity=0.90] (8.98,10.67) circle (0.089);
\fill[fill={rgb,255:red,58;green,154;blue,216}, fill opacity=0.94] (8.28,9.41) circle (0.094);
\fill[fill={rgb,255:red,58;green,154;blue,215}, fill opacity=0.89] (8.28,10.67) circle (0.088);
\fill[fill={rgb,255:red,92;green,149;blue,185}, fill opacity=0.79] (9.35,7.46) circle (0.075);
\fill[fill={rgb,255:red,67;green,158;blue,194}, fill opacity=0.78] (11.12,8.94) circle (0.074);
\fill[fill={rgb,255:red,50;green,179;blue,163}, fill opacity=0.77] (17.57,15.03) circle (0.073);
\fill[fill={rgb,255:red,62;green,153;blue,212}, fill opacity=0.88] (4.30,6.65) circle (0.086);
\fill[fill={rgb,255:red,54;green,166;blue,191}, fill opacity=0.77] (11.92,10.69) circle (0.073);
\fill[fill={rgb,255:red,60;green,153;blue,214}, fill opacity=0.90] (4.81,7.13) circle (0.088);
\fill[fill={rgb,255:red,58;green,153;blue,216}, fill opacity=0.95] (7.98,9.05) circle (0.095);
\fill[fill={rgb,255:red,58;green,154;blue,216}, fill opacity=0.92] (7.54,9.75) circle (0.091);
\fill[fill={rgb,255:red,58;green,153;blue,216}, fill opacity=0.93] (6.41,8.29) circle (0.093);
\fill[fill={rgb,255:red,57;green,157;blue,209}, fill opacity=0.84] (10.76,11.60) circle (0.081);
\fill[fill={rgb,255:red,58;green,155;blue,213}, fill opacity=0.87] (9.43,11.38) circle (0.084);
\fill[fill={rgb,255:red,54;green,164;blue,195}, fill opacity=0.79] (11.76,11.52) circle (0.075);
\fill[fill={rgb,255:red,59;green,153;blue,215}, fill opacity=0.90] (8.83,8.14) circle (0.088);
\fill[fill={rgb,255:red,73;green,151;blue,203}, fill opacity=0.84] (2.43,7.03) circle (0.081);
\fill[fill={rgb,255:red,58;green,154;blue,216}, fill opacity=0.90] (8.98,10.69) circle (0.089);
\fill[fill={rgb,255:red,152;green,138;blue,135}, fill opacity=0.74] (8.45,6.36) circle (0.069);
\fill[fill={rgb,255:red,58;green,153;blue,216}, fill opacity=0.95] (7.42,8.65) circle (0.095);
\fill[fill={rgb,255:red,58;green,153;blue,216}, fill opacity=0.92] (8.44,8.16) circle (0.091);
\fill[fill={rgb,255:red,58;green,154;blue,216}, fill opacity=0.94] (8.95,9.56) circle (0.094);
\fill[fill={rgb,255:red,61;green,153;blue,214}, fill opacity=0.89] (4.55,7.16) circle (0.087);
\fill[fill={rgb,255:red,59;green,154;blue,214}, fill opacity=0.88] (7.27,10.42) circle (0.086);
\fill[fill={rgb,255:red,58;green,154;blue,216}, fill opacity=0.92] (9.87,9.94) circle (0.091);
\fill[fill={rgb,255:red,48;green,179;blue,164}, fill opacity=0.78] (12.37,13.94) circle (0.073);
\fill[fill={rgb,255:red,60;green,153;blue,214}, fill opacity=0.89] (4.86,8.06) circle (0.087);
\fill[fill={rgb,255:red,58;green,153;blue,216}, fill opacity=0.93] (6.14,7.98) circle (0.092);
\fill[fill={rgb,255:red,58;green,153;blue,216}, fill opacity=0.95] (7.62,8.57) circle (0.095);
\fill[fill={rgb,255:red,58;green,154;blue,216}, fill opacity=0.94] (8.64,9.62) circle (0.094);
\fill[fill={rgb,255:red,65;green,152;blue,210}, fill opacity=0.86] (3.51,6.71) circle (0.084);
\fill[fill={rgb,255:red,58;green,153;blue,216}, fill opacity=0.94] (8.57,8.83) circle (0.094);
\fill[fill={rgb,255:red,58;green,156;blue,210}, fill opacity=0.84] (9.01,11.69) circle (0.081);
\fill[fill={rgb,255:red,58;green,153;blue,216}, fill opacity=0.94] (6.62,8.26) circle (0.093);
\fill[fill={rgb,255:red,58;green,154;blue,215}, fill opacity=0.90] (9.67,8.96) circle (0.089);
\fill[fill={rgb,255:red,58;green,154;blue,216}, fill opacity=0.91] (7.80,10.01) circle (0.090);
\fill[fill={rgb,255:red,58;green,154;blue,216}, fill opacity=0.94] (9.32,9.76) circle (0.093);
\fill[fill={rgb,255:red,58;green,154;blue,216}, fill opacity=0.94] (8.65,9.70) circle (0.094);
\fill[fill={rgb,255:red,47;green,181;blue,161}, fill opacity=0.79] (12.72,12.07) circle (0.075);
\fill[fill={rgb,255:red,57;green,155;blue,212}, fill opacity=0.86] (10.02,11.54) circle (0.083);
\fill[fill={rgb,255:red,58;green,154;blue,214}, fill opacity=0.88] (9.95,11.13) circle (0.086);
\fill[fill={rgb,255:red,66;green,152;blue,209}, fill opacity=0.87] (3.95,5.99) circle (0.084);
\fill[fill={rgb,255:red,58;green,154;blue,216}, fill opacity=0.91] (9.69,10.50) circle (0.091);
\fill[fill={rgb,255:red,58;green,153;blue,216}, fill opacity=0.93] (8.64,8.39) circle (0.092);
\fill[fill={rgb,255:red,62;green,153;blue,211}, fill opacity=0.85] (5.69,9.91) circle (0.083);
\fill[fill={rgb,255:red,58;green,154;blue,216}, fill opacity=0.94] (8.93,9.52) circle (0.094);
\fill[fill={rgb,255:red,57;green,157;blue,210}, fill opacity=0.84] (10.87,11.45) circle (0.082);
\fill[fill={rgb,255:red,51;green,173;blue,177}, fill opacity=0.74] (12.22,10.70) circle (0.069);
\fill[fill={rgb,255:red,66;green,152;blue,210}, fill opacity=0.87] (4.97,5.99) circle (0.085);
\fill[fill={rgb,255:red,64;green,153;blue,210}, fill opacity=0.86] (4.20,8.32) circle (0.084);
\fill[fill={rgb,255:red,58;green,153;blue,216}, fill opacity=0.93] (7.05,9.13) circle (0.092);
\fill[fill={rgb,255:red,73;green,151;blue,203}, fill opacity=0.84] (2.68,5.26) circle (0.081);
\fill[fill={rgb,255:red,60;green,154;blue,212}, fill opacity=0.86] (10.50,9.24) circle (0.083);
\fill[fill={rgb,255:red,59;green,154;blue,214}, fill opacity=0.88] (7.16,10.29) circle (0.086);
\fill[fill={rgb,255:red,61;green,153;blue,213}, fill opacity=0.89] (5.89,6.58) circle (0.087);
\fill[fill={rgb,255:red,58;green,155;blue,214}, fill opacity=0.88] (10.78,10.43) circle (0.086);
\fill[fill={rgb,255:red,58;green,153;blue,216}, fill opacity=0.95] (7.38,8.45) circle (0.095);
\fill[fill={rgb,255:red,58;green,153;blue,216}, fill opacity=0.94] (7.32,9.01) circle (0.094);
\fill[fill={rgb,255:red,58;green,154;blue,216}, fill opacity=0.94] (8.96,9.73) circle (0.094);
\fill[fill={rgb,255:red,58;green,153;blue,216}, fill opacity=0.95] (7.55,8.72) circle (0.095);
\fill[fill={rgb,255:red,62;green,153;blue,212}, fill opacity=0.86] (5.65,9.71) circle (0.084);
\fill[fill={rgb,255:red,58;green,153;blue,216}, fill opacity=0.93] (6.39,8.09) circle (0.093);
\fill[fill={rgb,255:red,48;green,180;blue,164}, fill opacity=0.78] (12.64,12.75) circle (0.074);
\fill[fill={rgb,255:red,43;green,191;blue,142}, fill opacity=0.91] (14.36,11.81) circle (0.091);
\fill[fill={rgb,255:red,59;green,153;blue,215}, fill opacity=0.90] (5.29,8.10) circle (0.089);
\fill[fill={rgb,255:red,65;green,153;blue,209}, fill opacity=0.85] (9.45,8.10) circle (0.083);
\fill[fill={rgb,255:red,58;green,154;blue,216}, fill opacity=0.93] (9.51,9.36) circle (0.092);
\fill[fill={rgb,255:red,65;green,153;blue,208}, fill opacity=0.84] (5.06,9.76) circle (0.082);
\fill[fill={rgb,255:red,58;green,153;blue,216}, fill opacity=0.94] (7.90,8.14) circle (0.093);
\fill[fill={rgb,255:red,56;green,160;blue,204}, fill opacity=0.82] (11.43,11.40) circle (0.078);
\fill[fill={rgb,255:red,58;green,153;blue,216}, fill opacity=0.91] (9.22,8.60) circle (0.090);
\fill[fill={rgb,255:red,58;green,153;blue,216}, fill opacity=0.91] (8.39,8.04) circle (0.090);
\fill[fill={rgb,255:red,62;green,153;blue,212}, fill opacity=0.88] (4.37,7.89) circle (0.086);
\fill[fill={rgb,255:red,60;green,153;blue,214}, fill opacity=0.90] (5.28,6.71) circle (0.088);
\fill[fill={rgb,255:red,58;green,153;blue,216}, fill opacity=0.94] (6.71,8.08) circle (0.093);
\fill[fill={rgb,255:red,58;green,154;blue,216}, fill opacity=0.94] (9.37,9.66) circle (0.093);
\fill[fill={rgb,255:red,48;green,181;blue,161}, fill opacity=0.78] (12.55,10.54) circle (0.074);
\fill[fill={rgb,255:red,58;green,153;blue,216}, fill opacity=0.92] (5.95,7.50) circle (0.091);
\fill[fill={rgb,255:red,58;green,153;blue,216}, fill opacity=0.91] (6.93,9.57) circle (0.090);
\fill[fill={rgb,255:red,44;green,189;blue,146}, fill opacity=0.86] (13.42,11.90) circle (0.083);
\fill[fill={rgb,255:red,58;green,153;blue,216}, fill opacity=0.94] (7.38,9.16) circle (0.093);
\fill[fill={rgb,255:red,58;green,154;blue,215}, fill opacity=0.89] (8.59,10.73) circle (0.088);
\fill[fill={rgb,255:red,59;green,153;blue,215}, fill opacity=0.91] (8.16,7.81) circle (0.090);
\fill[fill={rgb,255:red,57;green,156;blue,211}, fill opacity=0.85] (10.28,11.62) circle (0.082);
\fill[fill={rgb,255:red,80;green,150;blue,197}, fill opacity=0.82] (7.82,6.77) circle (0.079);
\fill[fill={rgb,255:red,63;green,153;blue,212}, fill opacity=0.87] (3.95,7.06) circle (0.086);
\fill[fill={rgb,255:red,63;green,153;blue,212}, fill opacity=0.88] (4.50,6.28) circle (0.086);
\fill[fill={rgb,255:red,58;green,154;blue,216}, fill opacity=0.90] (9.54,8.85) circle (0.089);
\fill[fill={rgb,255:red,83;green,149;blue,194}, fill opacity=0.82] (1.02,4.15) circle (0.079);
\fill[fill={rgb,255:red,69;green,152;blue,207}, fill opacity=0.86] (3.96,5.60) circle (0.083);
\fill[fill={rgb,255:red,65;green,152;blue,209}, fill opacity=0.86] (8.54,7.53) circle (0.084);
\fill[fill={rgb,255:red,58;green,154;blue,215}, fill opacity=0.90] (8.63,10.71) circle (0.088);
\fill[fill={rgb,255:red,54;green,164;blue,195}, fill opacity=0.78] (11.04,12.46) circle (0.074);
\fill[fill={rgb,255:red,58;green,153;blue,216}, fill opacity=0.92] (7.71,7.70) circle (0.091);
\fill[fill={rgb,255:red,58;green,154;blue,216}, fill opacity=0.92] (9.76,9.78) circle (0.092);
\fill[fill={rgb,255:red,58;green,153;blue,216}, fill opacity=0.95] (8.52,8.85) circle (0.094);
\fill[fill={rgb,255:red,58;green,154;blue,215}, fill opacity=0.89] (10.59,10.54) circle (0.087);
\fill[fill={rgb,255:red,56;green,163;blue,195}, fill opacity=0.78] (11.70,10.13) circle (0.074);
\fill[fill={rgb,255:red,58;green,154;blue,216}, fill opacity=0.91] (7.74,9.98) circle (0.090);
\fill[fill={rgb,255:red,49;green,176;blue,171}, fill opacity=0.76] (12.41,12.14) circle (0.071);
\fill[fill={rgb,255:red,58;green,154;blue,214}, fill opacity=0.88] (10.46,9.66) circle (0.086);
\fill[fill={rgb,255:red,59;green,153;blue,216}, fill opacity=0.91] (5.96,7.12) circle (0.090);
\fill[fill={rgb,255:red,68;green,152;blue,207}, fill opacity=0.86] (4.28,5.67) circle (0.083);
\fill[fill={rgb,255:red,58;green,153;blue,216}, fill opacity=0.94] (7.12,8.26) circle (0.094);
\fill[fill={rgb,255:red,58;green,154;blue,214}, fill opacity=0.88] (9.90,11.11) circle (0.087);
\fill[fill={rgb,255:red,58;green,154;blue,216}, fill opacity=0.92] (8.20,9.96) circle (0.092);
\fill[fill={rgb,255:red,58;green,153;blue,216}, fill opacity=0.92] (8.82,8.41) circle (0.091);
\fill[fill={rgb,255:red,58;green,154;blue,216}, fill opacity=0.95] (8.35,9.42) circle (0.094);
\fill[fill={rgb,255:red,61;green,153;blue,213}, fill opacity=0.89] (5.88,6.61) circle (0.088);
\fill[fill={rgb,255:red,66;green,153;blue,209}, fill opacity=0.85] (9.14,7.88) circle (0.083);
\fill[fill={rgb,255:red,60;green,153;blue,215}, fill opacity=0.90] (6.84,7.06) circle (0.089);
\fill[fill={rgb,255:red,45;green,185;blue,153}, fill opacity=0.82] (13.23,12.62) circle (0.079);
\fill[fill={rgb,255:red,75;green,151;blue,201}, fill opacity=0.84] (5.80,5.83) circle (0.081);
\fill[fill={rgb,255:red,56;green,160;blue,203}, fill opacity=0.81] (11.57,10.69) circle (0.077);
\fill[fill={rgb,255:red,56;green,158;blue,206}, fill opacity=0.82] (11.08,11.67) circle (0.079);
\fill[fill={rgb,255:red,58;green,154;blue,214}, fill opacity=0.87] (10.77,10.16) circle (0.085);
\fill[fill={rgb,255:red,58;green,153;blue,216}, fill opacity=0.92] (6.85,7.45) circle (0.092);
\fill[fill={rgb,255:red,58;green,153;blue,216}, fill opacity=0.92] (7.40,7.57) circle (0.091);
\fill[fill={rgb,255:red,50;green,181;blue,159}, fill opacity=0.78] (12.52,10.05) circle (0.074);
\fill[fill={rgb,255:red,72;green,151;blue,204}, fill opacity=0.84] (3.00,5.34) circle (0.082);
\fill[fill={rgb,255:red,58;green,153;blue,216}, fill opacity=0.95] (7.42,8.69) circle (0.095);
\fill[fill={rgb,255:red,58;green,153;blue,216}, fill opacity=0.95] (7.61,8.91) circle (0.094);
\fill[fill={rgb,255:red,79;green,150;blue,198}, fill opacity=0.83] (3.26,4.67) circle (0.080);
\fill[fill={rgb,255:red,45;green,187;blue,150}, fill opacity=0.83] (13.04,11.38) circle (0.080);
\fill[fill={rgb,255:red,54;green,166;blue,192}, fill opacity=0.78] (11.45,12.36) circle (0.073);
\fill[fill={rgb,255:red,60;green,153;blue,214}, fill opacity=0.90] (6.30,6.83) circle (0.088);
\fill[fill={rgb,255:red,51;green,172;blue,179}, fill opacity=0.74] (12.21,11.58) circle (0.069);
\fill[fill={rgb,255:red,150;green,138;blue,138}, fill opacity=0.74] (5.70,4.24) circle (0.069);
\fill[fill={rgb,255:red,59;green,153;blue,215}, fill opacity=0.91] (7.19,7.31) circle (0.090);
\fill[fill={rgb,255:red,64;green,153;blue,211}, fill opacity=0.87] (5.74,6.31) circle (0.086);
\fill[fill={rgb,255:red,61;green,153;blue,213}, fill opacity=0.87] (9.72,8.53) circle (0.085);
\fill[fill={rgb,255:red,57;green,159;blue,204}, fill opacity=0.82] (10.10,12.18) circle (0.078);
\fill[fill={rgb,255:red,58;green,153;blue,216}, fill opacity=0.92] (5.99,7.28) circle (0.091);
\fill[fill={rgb,255:red,62;green,163;blue,188}, fill opacity=0.76] (11.58,9.35) circle (0.071);
\fill[fill={rgb,255:red,58;green,154;blue,216}, fill opacity=0.94] (9.22,9.92) circle (0.093);
\fill[fill={rgb,255:red,58;green,154;blue,216}, fill opacity=0.95] (8.38,9.27) circle (0.095);
\fill[fill={rgb,255:red,56;green,158;blue,207}, fill opacity=0.83] (11.30,11.29) circle (0.080);
\fill[fill={rgb,255:red,58;green,154;blue,215}, fill opacity=0.89] (7.49,10.29) circle (0.088);
\fill[fill={rgb,255:red,71;green,151;blue,204}, fill opacity=0.85] (2.89,5.47) circle (0.082);
\fill[fill={rgb,255:red,58;green,154;blue,215}, fill opacity=0.90] (10.29,10.61) circle (0.088);
\fill[fill={rgb,255:red,58;green,154;blue,216}, fill opacity=0.90] (10.21,9.76) circle (0.089);
\fill[fill={rgb,255:red,58;green,153;blue,216}, fill opacity=0.92] (8.83,8.46) circle (0.091);
\fill[fill={rgb,255:red,58;green,153;blue,216}, fill opacity=0.93] (6.43,7.93) circle (0.093);
\fill[fill={rgb,255:red,58;green,153;blue,216}, fill opacity=0.95] (7.48,8.55) circle (0.095);
\fill[fill={rgb,255:red,58;green,154;blue,214}, fill opacity=0.87] (10.66,9.83) circle (0.085);
\fill[fill={rgb,255:red,58;green,154;blue,216}, fill opacity=0.94] (9.06,9.20) circle (0.094);
\fill[fill={rgb,255:red,58;green,153;blue,216}, fill opacity=0.95] (7.43,8.73) circle (0.094);
\fill[fill={rgb,255:red,53;green,167;blue,188}, fill opacity=0.77] (11.90,11.86) circle (0.072);
\fill[fill={rgb,255:red,51;green,172;blue,179}, fill opacity=0.74] (12.22,11.28) circle (0.069);
\fill[fill={rgb,255:red,78;green,152;blue,195}, fill opacity=0.80] (10.05,8.09) circle (0.076);
\fill[fill={rgb,255:red,58;green,153;blue,216}, fill opacity=0.94] (7.75,9.13) circle (0.094);
\fill[fill={rgb,255:red,58;green,154;blue,216}, fill opacity=0.95] (8.77,9.15) circle (0.094);
\fill[fill={rgb,255:red,58;green,154;blue,216}, fill opacity=0.91] (9.68,9.12) circle (0.090);
\fill[fill={rgb,255:red,58;green,154;blue,216}, fill opacity=0.94] (9.15,9.26) circle (0.094);
\fill[fill={rgb,255:red,58;green,153;blue,216}, fill opacity=0.91] (8.49,8.04) circle (0.090);
\fill[fill={rgb,255:red,58;green,153;blue,216}, fill opacity=0.94] (7.35,8.84) circle (0.094);
\fill[fill={rgb,255:red,58;green,154;blue,216}, fill opacity=0.95] (8.73,9.10) circle (0.094);
\fill[fill={rgb,255:red,70;green,152;blue,206}, fill opacity=0.85] (7.44,6.79) circle (0.082);
\fill[fill={rgb,255:red,58;green,153;blue,216}, fill opacity=0.93] (6.66,8.49) circle (0.093);
\fill[fill={rgb,255:red,58;green,153;blue,216}, fill opacity=0.93] (7.08,8.94) circle (0.093);
\fill[fill={rgb,255:red,43;green,191;blue,141}, fill opacity=0.92] (15.92,9.16) circle (0.091);
\fill[fill={rgb,255:red,43;green,191;blue,142}, fill opacity=0.91] (15.11,12.06) circle (0.090);
\fill[fill={rgb,255:red,47;green,188;blue,144}, fill opacity=0.84] (13.23,9.48) circle (0.081);
\fill[fill={rgb,255:red,43;green,190;blue,143}, fill opacity=0.88] (16.17,12.57) circle (0.086);
\fill[fill={rgb,255:red,43;green,190;blue,144}, fill opacity=0.88] (14.38,12.52) circle (0.086);
\fill[fill={rgb,255:red,70;green,182;blue,131}, fill opacity=0.82] (18.61,8.70) circle (0.078);
\fill[fill={rgb,255:red,135;green,155;blue,114}, fill opacity=0.71] (12.40,7.91) circle (0.065);
\fill[fill={rgb,255:red,47;green,187;blue,146}, fill opacity=0.83] (13.03,9.80) circle (0.080);
\fill[fill={rgb,255:red,82;green,178;blue,126}, fill opacity=0.80] (16.28,7.68) circle (0.077);
\fill[fill={rgb,255:red,45;green,186;blue,151}, fill opacity=0.83] (13.28,12.42) circle (0.080);
\fill[fill={rgb,255:red,45;green,189;blue,145}, fill opacity=0.85] (17.25,12.65) circle (0.083);
\fill[fill={rgb,255:red,46;green,184;blue,156}, fill opacity=0.81] (13.20,13.23) circle (0.077);
\fill[fill={rgb,255:red,46;green,190;blue,141}, fill opacity=0.87] (14.13,8.94) circle (0.085);
\fill[fill={rgb,255:red,44;green,191;blue,141}, fill opacity=0.91] (16.84,9.73) circle (0.090);
\fill[fill={rgb,255:red,43;green,191;blue,142}, fill opacity=0.94] (15.02,11.33) circle (0.094);
\fill[fill={rgb,255:red,44;green,188;blue,149}, fill opacity=0.84] (13.13,11.55) circle (0.081);
\fill[fill={rgb,255:red,43;green,191;blue,142}, fill opacity=0.92] (16.39,11.13) circle (0.091);
\fill[fill={rgb,255:red,43;green,190;blue,143}, fill opacity=0.88] (15.99,12.56) circle (0.086);
\fill[fill={rgb,255:red,43;green,191;blue,142}, fill opacity=0.94] (14.49,10.94) circle (0.094);
\fill[fill={rgb,255:red,43;green,191;blue,142}, fill opacity=0.89] (13.67,10.23) circle (0.088);
\fill[fill={rgb,255:red,78;green,179;blue,128}, fill opacity=0.81] (18.26,8.22) circle (0.077);
\fill[fill={rgb,255:red,43;green,191;blue,142}, fill opacity=0.94] (15.86,10.73) circle (0.094);
\fill[fill={rgb,255:red,124;green,163;blue,109}, fill opacity=0.76] (17.15,7.23) circle (0.071);
\fill[fill={rgb,255:red,46;green,190;blue,140}, fill opacity=0.89] (15.09,8.68) circle (0.087);
\fill[fill={rgb,255:red,188;green,141;blue,84}, fill opacity=0.79] (17.91,6.29) circle (0.075);
\fill[fill={rgb,255:red,43;green,191;blue,143}, fill opacity=0.89] (14.71,12.46) circle (0.088);
\fill[fill={rgb,255:red,43;green,191;blue,141}, fill opacity=0.92] (16.16,9.48) circle (0.092);
\fill[fill={rgb,255:red,61;green,156;blue,206}, fill opacity=0.82] (10.94,9.28) circle (0.079);
\fill[fill={rgb,255:red,43;green,191;blue,142}, fill opacity=0.93] (16.24,10.83) circle (0.092);
\fill[fill={rgb,255:red,45;green,188;blue,148}, fill opacity=0.84] (16.65,13.32) circle (0.081);
\fill[fill={rgb,255:red,44;green,189;blue,145}, fill opacity=0.86] (14.32,12.81) circle (0.084);
\fill[fill={rgb,255:red,43;green,191;blue,142}, fill opacity=0.91] (16.00,11.80) circle (0.090);
\fill[fill={rgb,255:red,43;green,191;blue,141}, fill opacity=0.93] (16.05,9.56) circle (0.092);
\fill[fill={rgb,255:red,46;green,185;blue,153}, fill opacity=0.81] (16.68,13.96) circle (0.078);
\fill[fill={rgb,255:red,43;green,191;blue,142}, fill opacity=0.95] (15.55,10.22) circle (0.094);
\fill[fill={rgb,255:red,44;green,190;blue,143}, fill opacity=0.88] (16.83,12.14) circle (0.086);
\fill[fill={rgb,255:red,44;green,190;blue,143}, fill opacity=0.87] (17.29,11.86) circle (0.085);
\fill[fill={rgb,255:red,47;green,190;blue,140}, fill opacity=0.88] (16.42,8.68) circle (0.087);
\fill[fill={rgb,255:red,43;green,191;blue,142}, fill opacity=0.91] (16.65,10.93) circle (0.091);
\fill[fill={rgb,255:red,52;green,170;blue,182}, fill opacity=0.75] (12.14,10.95) circle (0.070);
\fill[fill={rgb,255:red,45;green,190;blue,141}, fill opacity=0.89] (17.01,9.26) circle (0.088);
\fill[fill={rgb,255:red,79;green,178;blue,128}, fill opacity=0.80] (14.34,7.97) circle (0.076);
\fill[fill={rgb,255:red,43;green,190;blue,145}, fill opacity=0.87] (14.67,12.82) circle (0.085);
\fill[fill={rgb,255:red,43;green,191;blue,142}, fill opacity=0.93] (14.16,10.96) circle (0.092);
\fill[fill={rgb,255:red,43;green,190;blue,144}, fill opacity=0.87] (13.76,12.11) circle (0.085);
\fill[fill={rgb,255:red,46;green,184;blue,156}, fill opacity=0.81] (12.83,11.60) circle (0.077);
\fill[fill={rgb,255:red,45;green,191;blue,141}, fill opacity=0.90] (16.77,9.24) circle (0.089);
\fill[fill={rgb,255:red,44;green,191;blue,141}, fill opacity=0.90] (17.16,9.95) circle (0.089);
\fill[fill={rgb,255:red,43;green,191;blue,142}, fill opacity=0.91] (13.90,10.13) circle (0.090);
\fill[fill={rgb,255:red,43;green,191;blue,142}, fill opacity=0.95] (14.86,10.34) circle (0.095);
\fill[fill={rgb,255:red,43;green,190;blue,145}, fill opacity=0.87] (13.41,11.55) circle (0.085);
\fill[fill={rgb,255:red,44;green,191;blue,142}, fill opacity=0.88] (17.49,10.91) circle (0.087);
\fill[fill={rgb,255:red,49;green,180;blue,163}, fill opacity=0.78] (15.01,16.32) circle (0.073);
\fill[fill={rgb,255:red,43;green,191;blue,142}, fill opacity=0.94] (15.59,10.86) circle (0.094);
\fill[fill={rgb,255:red,46;green,184;blue,155}, fill opacity=0.81] (13.25,13.10) circle (0.078);
\fill[fill={rgb,255:red,46;green,190;blue,140}, fill opacity=0.89] (16.81,8.95) circle (0.087);
\fill[fill={rgb,255:red,43;green,191;blue,142}, fill opacity=0.91] (13.77,10.55) circle (0.090);
\fill[fill={rgb,255:red,220;green,129;blue,71}, fill opacity=0.84] (15.73,6.32) circle (0.081);
\fill[fill={rgb,255:red,43;green,191;blue,143}, fill opacity=0.89] (13.58,10.99) circle (0.088);
\fill[fill={rgb,255:red,43;green,191;blue,142}, fill opacity=0.94] (15.85,10.49) circle (0.094);
\fill[fill={rgb,255:red,43;green,191;blue,142}, fill opacity=0.92] (14.17,11.31) circle (0.092);
\fill[fill={rgb,255:red,43;green,191;blue,142}, fill opacity=0.90] (13.63,10.58) circle (0.088);
\fill[fill={rgb,255:red,45;green,191;blue,141}, fill opacity=0.89] (14.58,8.95) circle (0.088);
\fill[fill={rgb,255:red,44;green,189;blue,145}, fill opacity=0.86] (13.89,12.49) circle (0.084);
\fill[fill={rgb,255:red,44;green,191;blue,141}, fill opacity=0.91] (15.92,8.92) circle (0.089);
\fill[fill={rgb,255:red,43;green,191;blue,142}, fill opacity=0.92] (16.57,10.81) circle (0.091);
\fill[fill={rgb,255:red,179;green,144;blue,87}, fill opacity=0.78] (16.17,6.78) circle (0.074);
\fill[fill={rgb,255:red,51;green,186;blue,144}, fill opacity=0.82] (19.08,12.25) circle (0.078);
\fill[fill={rgb,255:red,43;green,191;blue,143}, fill opacity=0.89] (14.02,12.10) circle (0.087);
\fill[fill={rgb,255:red,46;green,186;blue,151}, fill opacity=0.82] (17.06,13.71) circle (0.078);
\fill[fill={rgb,255:red,43;green,191;blue,141}, fill opacity=0.94] (15.24,9.79) circle (0.094);
\fill[fill={rgb,255:red,44;green,189;blue,147}, fill opacity=0.85] (14.45,13.12) circle (0.083);
\fill[fill={rgb,255:red,43;green,191;blue,142}, fill opacity=0.91] (13.89,11.24) circle (0.090);
\fill[fill={rgb,255:red,48;green,181;blue,161}, fill opacity=0.78] (15.05,15.50) circle (0.074);
\fill[fill={rgb,255:red,43;green,191;blue,142}, fill opacity=0.91] (15.52,12.17) circle (0.089);
\fill[fill={rgb,255:red,118;green,165;blue,112}, fill opacity=0.76] (18.84,7.60) circle (0.071);
\fill[fill={rgb,255:red,43;green,191;blue,142}, fill opacity=0.93] (16.20,10.61) circle (0.093);
\fill[fill={rgb,255:red,43;green,191;blue,142}, fill opacity=0.92] (14.08,11.41) circle (0.091);
\fill[fill={rgb,255:red,87;green,176;blue,124}, fill opacity=0.79] (18.06,7.93) circle (0.076);
\fill[fill={rgb,255:red,44;green,191;blue,141}, fill opacity=0.91] (15.95,8.94) circle (0.089);
\fill[fill={rgb,255:red,43;green,191;blue,141}, fill opacity=0.92] (15.75,9.17) circle (0.091);
\fill[fill={rgb,255:red,133;green,159;blue,106}, fill opacity=0.74] (19.93,7.65) circle (0.069);
\fill[fill={rgb,255:red,43;green,191;blue,142}, fill opacity=0.89] (16.47,12.02) circle (0.088);
\fill[fill={rgb,255:red,47;green,181;blue,161}, fill opacity=0.79] (12.86,14.02) circle (0.075);
\fill[fill={rgb,255:red,190;green,140;blue,83}, fill opacity=0.79] (16.62,6.56) circle (0.076);
\fill[fill={rgb,255:red,64;green,184;blue,133}, fill opacity=0.83] (16.34,7.96) circle (0.080);
\fill[fill={rgb,255:red,65;green,180;blue,143}, fill opacity=0.78] (12.78,8.81) circle (0.073);
\fill[fill={rgb,255:red,47;green,187;blue,146}, fill opacity=0.83] (13.05,9.72) circle (0.080);
\fill[fill={rgb,255:red,43;green,191;blue,142}, fill opacity=0.95] (14.87,10.96) circle (0.095);
\fill[fill={rgb,255:red,43;green,191;blue,141}, fill opacity=0.94] (15.71,9.92) circle (0.094);
\fill[fill={rgb,255:red,43;green,191;blue,141}, fill opacity=0.91] (15.33,9.02) circle (0.090);
\fill[fill={rgb,255:red,100;green,171;blue,119}, fill opacity=0.78] (17.55,7.57) circle (0.074);
\fill[fill={rgb,255:red,48;green,180;blue,163}, fill opacity=0.78] (15.49,15.62) circle (0.074);
\fill[fill={rgb,255:red,45;green,187;blue,149}, fill opacity=0.84] (14.57,13.51) circle (0.081);
\fill[fill={rgb,255:red,47;green,190;blue,141}, fill opacity=0.86] (18.09,10.71) circle (0.084);
\fill[fill={rgb,255:red,43;green,191;blue,142}, fill opacity=0.95] (15.39,10.38) circle (0.095);
\fill[fill={rgb,255:red,43;green,191;blue,141}, fill opacity=0.94] (15.89,10.12) circle (0.094);
\fill[fill={rgb,255:red,43;green,191;blue,141}, fill opacity=0.93] (15.86,9.70) circle (0.093);
\fill[fill={rgb,255:red,53;green,188;blue,138}, fill opacity=0.86] (16.34,8.30) circle (0.084);
\fill[fill={rgb,255:red,48;green,180;blue,162}, fill opacity=0.78] (12.35,14.46) circle (0.074);
\fill[fill={rgb,255:red,43;green,191;blue,142}, fill opacity=0.93] (16.38,10.62) circle (0.092);
\fill[fill={rgb,255:red,44;green,188;blue,147}, fill opacity=0.84] (13.08,10.73) circle (0.082);
\fill[fill={rgb,255:red,43;green,191;blue,142}, fill opacity=0.91] (16.82,11.04) circle (0.090);
\fill[fill={rgb,255:red,43;green,191;blue,142}, fill opacity=0.94] (15.20,11.13) circle (0.094);
\fill[fill={rgb,255:red,48;green,178;blue,167}, fill opacity=0.77] (12.53,12.62) circle (0.073);
\fill[fill={rgb,255:red,43;green,191;blue,142}, fill opacity=0.94] (14.64,10.20) circle (0.094);
\fill[fill={rgb,255:red,57;green,157;blue,210}, fill opacity=0.84] (11.01,11.32) circle (0.082);
\fill[fill={rgb,255:red,44;green,191;blue,141}, fill opacity=0.90] (15.60,8.87) circle (0.089);
\fill[fill={rgb,255:red,43;green,191;blue,142}, fill opacity=0.90] (15.30,12.27) circle (0.089);
\fill[fill={rgb,255:red,55;green,162;blue,199}, fill opacity=0.80] (11.61,11.50) circle (0.076);
\fill[fill={rgb,255:red,43;green,191;blue,142}, fill opacity=0.91] (14.06,11.58) circle (0.090);
\fill[fill={rgb,255:red,44;green,191;blue,141}, fill opacity=0.91] (16.94,9.88) circle (0.089);
\fill[fill={rgb,255:red,204;green,135;blue,77}, fill opacity=0.81] (13.98,6.99) circle (0.077);
\fill[fill={rgb,255:red,56;green,186;blue,137}, fill opacity=0.83] (14.00,8.49) circle (0.080);
\fill[fill={rgb,255:red,43;green,191;blue,141}, fill opacity=0.91] (16.71,9.76) circle (0.090);
\fill[fill={rgb,255:red,44;green,189;blue,144}, fill opacity=0.86] (16.99,12.57) circle (0.084);
\fill[fill={rgb,255:red,52;green,187;blue,139}, fill opacity=0.84] (13.73,8.76) circle (0.081);
\fill[fill={rgb,255:red,50;green,189;blue,139}, fill opacity=0.87] (15.22,8.36) circle (0.084);
\fill[fill={rgb,255:red,52;green,188;blue,138}, fill opacity=0.86] (17.53,8.82) circle (0.083);
\fill[fill={rgb,255:red,141;green,156;blue,104}, fill opacity=0.73] (13.75,7.56) circle (0.068);
\fill[fill={rgb,255:red,43;green,191;blue,142}, fill opacity=0.92] (14.35,11.65) circle (0.091);
\fill[fill={rgb,255:red,43;green,191;blue,142}, fill opacity=0.94] (15.07,11.26) circle (0.094);
\fill[fill={rgb,255:red,43;green,191;blue,142}, fill opacity=0.93] (14.63,11.68) circle (0.092);
\fill[fill={rgb,255:red,44;green,188;blue,147}, fill opacity=0.85] (14.33,13.10) circle (0.083);
\fill[fill={rgb,255:red,43;green,191;blue,142}, fill opacity=0.93] (15.41,11.61) circle (0.092);
\fill[fill={rgb,255:red,77;green,179;blue,128}, fill opacity=0.81] (17.08,7.84) circle (0.078);
\fill[fill={rgb,255:red,43;green,191;blue,142}, fill opacity=0.94] (16.00,10.28) circle (0.093);
\fill[fill={rgb,255:red,63;green,184;blue,135}, fill opacity=0.82] (19.22,9.70) circle (0.079);
\fill[fill={rgb,255:red,44;green,190;blue,142}, fill opacity=0.88] (13.74,9.59) circle (0.087);
\fill[fill={rgb,255:red,120;green,164;blue,111}, fill opacity=0.76] (17.27,7.29) circle (0.071);
\fill[fill={rgb,255:red,43;green,191;blue,142}, fill opacity=0.93] (14.61,9.71) circle (0.092);
\fill[fill={rgb,255:red,43;green,191;blue,142}, fill opacity=0.94] (14.47,10.10) circle (0.093);
\fill[fill={rgb,255:red,43;green,191;blue,142}, fill opacity=0.91] (13.77,10.83) circle (0.090);
\fill[fill={rgb,255:red,43;green,191;blue,142}, fill opacity=0.90] (16.68,11.39) circle (0.089);
\fill[fill={rgb,255:red,51;green,175;blue,172}, fill opacity=0.75] (10.79,14.09) circle (0.070);
\fill[fill={rgb,255:red,44;green,191;blue,141}, fill opacity=0.90] (14.95,8.94) circle (0.089);
\fill[fill={rgb,255:red,45;green,187;blue,150}, fill opacity=0.83] (12.94,10.84) circle (0.080);
\fill[fill={rgb,255:red,43;green,191;blue,141}, fill opacity=0.93] (15.99,9.67) circle (0.093);
\fill[fill={rgb,255:red,66;green,183;blue,133}, fill opacity=0.83] (15.47,7.94) circle (0.080);
\fill[fill={rgb,255:red,130;green,160;blue,109}, fill opacity=0.73] (13.37,7.74) circle (0.068);
\fill[fill={rgb,255:red,88;green,175;blue,124}, fill opacity=0.80] (16.76,7.63) circle (0.076);
\fill[fill={rgb,255:red,43;green,191;blue,142}, fill opacity=0.95] (14.97,10.63) circle (0.095);
\fill[fill={rgb,255:red,145;green,155;blue,101}, fill opacity=0.75] (17.12,7.01) circle (0.069);
\fill[fill={rgb,255:red,43;green,191;blue,143}, fill opacity=0.89] (15.06,12.45) circle (0.088);
\fill[fill={rgb,255:red,43;green,191;blue,142}, fill opacity=0.92] (14.38,9.57) circle (0.091);
\fill[fill={rgb,255:red,43;green,191;blue,142}, fill opacity=0.91] (13.85,11.10) circle (0.090);
\fill[fill={rgb,255:red,77;green,179;blue,128}, fill opacity=0.81] (15.91,7.75) circle (0.078);
\fill[fill={rgb,255:red,45;green,186;blue,151}, fill opacity=0.83] (13.79,13.28) circle (0.080);
\fill[fill={rgb,255:red,44;green,187;blue,149}, fill opacity=0.84] (13.29,12.12) circle (0.081);
\fill[fill={rgb,255:red,44;green,189;blue,145}, fill opacity=0.87] (13.42,11.62) circle (0.084);
\fill[fill={rgb,255:red,49;green,186;blue,146}, fill opacity=0.82] (18.73,12.59) circle (0.078);
\fill[fill={rgb,255:red,43;green,191;blue,143}, fill opacity=0.90] (13.77,11.41) circle (0.089);
\fill[fill={rgb,255:red,43;green,191;blue,142}, fill opacity=0.94] (15.92,10.70) circle (0.094);
\fill[fill={rgb,255:red,45;green,187;blue,151}, fill opacity=0.83] (13.21,12.20) circle (0.080);
\fill[fill={rgb,255:red,125;green,163;blue,109}, fill opacity=0.76] (17.31,7.24) circle (0.071);
\fill[fill={rgb,255:red,43;green,191;blue,142}, fill opacity=0.95] (15.52,10.67) circle (0.095);
\fill[fill={rgb,255:red,44;green,188;blue,148}, fill opacity=0.84] (14.09,13.14) circle (0.082);
\fill[fill={rgb,255:red,45;green,187;blue,150}, fill opacity=0.83] (12.96,10.93) circle (0.080);
\fill[fill={rgb,255:red,47;green,184;blue,153}, fill opacity=0.81] (16.91,14.02) circle (0.077);
\fill[fill={rgb,255:red,43;green,191;blue,142}, fill opacity=0.91] (14.71,12.13) circle (0.090);
\fill[fill={rgb,255:red,45;green,189;blue,144}, fill opacity=0.85] (17.45,12.48) circle (0.083);
\fill[fill={rgb,255:red,46;green,190;blue,140}, fill opacity=0.89] (16.46,8.82) circle (0.088);
\fill[fill={rgb,255:red,43;green,191;blue,142}, fill opacity=0.91] (15.74,12.04) circle (0.090);
\fill[fill={rgb,255:red,43;green,191;blue,141}, fill opacity=0.93] (16.35,10.06) circle (0.092);
\fill[fill={rgb,255:red,43;green,191;blue,141}, fill opacity=0.94] (15.91,10.17) circle (0.094);
\fill[fill={rgb,255:red,44;green,191;blue,141}, fill opacity=0.89] (14.13,9.24) circle (0.087);
\fill[fill={rgb,255:red,43;green,191;blue,142}, fill opacity=0.91] (16.53,11.32) circle (0.090);
\fill[fill={rgb,255:red,47;green,190;blue,140}, fill opacity=0.88] (16.35,8.63) circle (0.086);
\fill[fill={rgb,255:red,44;green,189;blue,146}, fill opacity=0.86] (13.19,10.70) circle (0.083);
\fill[fill={rgb,255:red,43;green,191;blue,142}, fill opacity=0.90] (13.83,9.87) circle (0.088);
\fill[fill={rgb,255:red,48;green,182;blue,159}, fill opacity=0.79] (16.84,14.69) circle (0.075);
\fill[fill={rgb,255:red,48;green,180;blue,163}, fill opacity=0.78] (12.50,13.87) circle (0.074);
\fill[fill={rgb,255:red,43;green,191;blue,142}, fill opacity=0.91] (14.59,11.95) circle (0.091);
\fill[fill={rgb,255:red,44;green,190;blue,143}, fill opacity=0.87] (17.21,12.17) circle (0.085);
\fill[fill={rgb,255:red,58;green,186;blue,136}, fill opacity=0.84] (18.44,9.23) circle (0.081);
\fill[fill={rgb,255:red,44;green,190;blue,142}, fill opacity=0.88] (17.24,11.77) circle (0.086);
\fill[fill={rgb,255:red,43;green,191;blue,142}, fill opacity=0.93] (14.31,10.24) circle (0.093);
\fill[fill={rgb,255:red,43;green,191;blue,142}, fill opacity=0.92] (15.97,11.71) circle (0.091);
\fill[fill={rgb,255:red,43;green,191;blue,142}, fill opacity=0.93] (15.97,11.06) circle (0.093);
\fill[fill={rgb,255:red,143;green,155;blue,104}, fill opacity=0.73] (13.39,7.63) circle (0.068);
\fill[fill={rgb,255:red,43;green,190;blue,144}, fill opacity=0.88] (15.11,12.77) circle (0.086);
\fill[fill={rgb,255:red,43;green,191;blue,142}, fill opacity=0.95] (15.29,10.86) circle (0.095);
\fill[fill={rgb,255:red,44;green,189;blue,146}, fill opacity=0.86] (13.48,11.92) circle (0.084);
\fill[fill={rgb,255:red,43;green,191;blue,141}, fill opacity=0.91] (16.08,9.13) circle (0.090);
\fill[fill={rgb,255:red,52;green,169;blue,184}, fill opacity=0.75] (12.11,11.32) circle (0.071);
\fill[fill={rgb,255:red,45;green,185;blue,153}, fill opacity=0.82] (13.05,12.02) circle (0.079);
\fill[fill={rgb,255:red,52;green,188;blue,139}, fill opacity=0.85] (18.41,9.80) circle (0.082);
\fill[fill={rgb,255:red,46;green,190;blue,141}, fill opacity=0.87] (13.89,9.15) circle (0.085);
\fill[fill={rgb,255:red,43;green,191;blue,142}, fill opacity=0.90] (17.04,10.79) circle (0.089);
\fill[fill={rgb,255:red,44;green,187;blue,149}, fill opacity=0.84] (13.31,12.22) circle (0.081);
\fill[fill={rgb,255:red,43;green,191;blue,142}, fill opacity=0.91] (15.77,11.92) circle (0.090);
\fill[fill={rgb,255:red,43;green,191;blue,141}, fill opacity=0.91] (16.68,9.75) circle (0.090);
\fill[fill={rgb,255:red,43;green,191;blue,142}, fill opacity=0.94] (14.77,9.98) circle (0.094);
\fill[fill={rgb,255:red,43;green,191;blue,141}, fill opacity=0.93] (16.12,9.74) circle (0.092);
\fill[fill={rgb,255:red,43;green,191;blue,142}, fill opacity=0.94] (14.65,10.06) circle (0.094);
\fill[fill={rgb,255:red,44;green,191;blue,142}, fill opacity=0.89] (17.35,10.84) circle (0.088);
\fill[fill={rgb,255:red,43;green,191;blue,142}, fill opacity=0.91] (16.23,11.74) circle (0.090);
\fill[fill={rgb,255:red,44;green,189;blue,147}, fill opacity=0.85] (14.58,13.19) circle (0.083);
\fill[fill={rgb,255:red,50;green,175;blue,173}, fill opacity=0.75] (11.41,13.73) circle (0.070);
\fill[fill={rgb,255:red,43;green,191;blue,142}, fill opacity=0.90] (13.95,11.65) circle (0.089);
\fill[fill={rgb,255:red,48;green,181;blue,161}, fill opacity=0.79] (15.13,15.37) circle (0.074);
\fill[fill={rgb,255:red,208;green,134;blue,75}, fill opacity=0.82] (16.72,6.22) circle (0.079);
\fill[fill={rgb,255:red,43;green,191;blue,143}, fill opacity=0.89] (15.53,12.44) circle (0.088);
\fill[fill={rgb,255:red,43;green,191;blue,142}, fill opacity=0.92] (14.35,11.76) circle (0.091);
\fill[fill={rgb,255:red,43;green,191;blue,142}, fill opacity=0.91] (15.63,11.95) circle (0.090);
\fill[fill={rgb,255:red,43;green,191;blue,142}, fill opacity=0.90] (16.79,11.37) circle (0.089);
\fill[fill={rgb,255:red,143;green,156;blue,102}, fill opacity=0.74] (19.13,7.19) circle (0.069);
\fill[fill={rgb,255:red,49;green,189;blue,139}, fill opacity=0.86] (17.95,9.63) circle (0.084);
\fill[fill={rgb,255:red,43;green,191;blue,141}, fill opacity=0.91] (14.58,9.31) circle (0.090);
\fill[fill={rgb,255:red,43;green,191;blue,142}, fill opacity=0.90] (13.77,10.34) circle (0.089);
\fill[fill={rgb,255:red,126;green,162;blue,109}, fill opacity=0.75] (14.61,7.48) circle (0.070);
\fill[fill={rgb,255:red,43;green,191;blue,142}, fill opacity=0.93] (16.06,10.96) circle (0.093);
\fill[fill={rgb,255:red,43;green,190;blue,144}, fill opacity=0.88] (13.89,12.19) circle (0.086);
\fill[fill={rgb,255:red,43;green,190;blue,144}, fill opacity=0.88] (14.06,12.39) circle (0.086);
\fill[fill={rgb,255:red,43;green,191;blue,142}, fill opacity=0.93] (16.24,10.91) circle (0.092);
\fill[fill={rgb,255:red,44;green,189;blue,145}, fill opacity=0.86] (13.59,12.07) circle (0.084);
\fill[fill={rgb,255:red,43;green,191;blue,141}, fill opacity=0.94] (15.51,9.96) circle (0.094);
\fill[fill={rgb,255:red,43;green,191;blue,142}, fill opacity=0.92] (14.16,9.84) circle (0.091);
\fill[fill={rgb,255:red,43;green,189;blue,145}, fill opacity=0.87] (13.56,11.96) circle (0.084);
\fill[fill={rgb,255:red,43;green,191;blue,142}, fill opacity=0.91] (13.76,10.58) circle (0.089);
\fill[fill={rgb,255:red,239;green,123;blue,63}, fill opacity=0.92] (11.62,5.85) circle (0.091);
\fill[fill={rgb,255:red,68;green,182;blue,133}, fill opacity=0.81] (13.97,8.23) circle (0.077);
\fill[fill={rgb,255:red,169;green,141;blue,108}, fill opacity=0.74] (11.35,7.57) circle (0.069);
\fill[fill={rgb,255:red,238;green,123;blue,64}, fill opacity=0.90] (9.78,4.27) circle (0.089);
\fill[fill={rgb,255:red,236;green,124;blue,64}, fill opacity=0.90] (14.11,6.24) circle (0.088);
\fill[fill={rgb,255:red,239;green,123;blue,63}, fill opacity=0.94] (12.68,5.56) circle (0.094);
\fill[fill={rgb,255:red,234;green,124;blue,65}, fill opacity=0.88] (14.16,3.41) circle (0.086);
\fill[fill={rgb,255:red,229;green,126;blue,67}, fill opacity=0.86] (14.02,2.88) circle (0.084);
\fill[fill={rgb,255:red,239;green,123;blue,63}, fill opacity=0.93] (11.76,4.46) circle (0.093);
\fill[fill={rgb,255:red,237;green,124;blue,64}, fill opacity=0.90] (15.34,5.58) circle (0.089);
\fill[fill={rgb,255:red,238;green,123;blue,63}, fill opacity=0.91] (12.21,3.63) circle (0.089);
\fill[fill={rgb,255:red,201;green,136;blue,78}, fill opacity=0.81] (18.95,4.99) circle (0.077);
\fill[fill={rgb,255:red,192;green,139;blue,82}, fill opacity=0.79] (19.92,5.25) circle (0.076);
\fill[fill={rgb,255:red,223;green,126;blue,75}, fill opacity=0.83] (7.72,3.50) circle (0.081);
\fill[fill={rgb,255:red,237;green,123;blue,64}, fill opacity=0.90] (10.93,3.53) circle (0.089);
\fill[fill={rgb,255:red,239;green,123;blue,63}, fill opacity=0.94] (11.94,5.00) circle (0.094);
\fill[fill={rgb,255:red,238;green,123;blue,63}, fill opacity=0.90] (13.00,6.26) circle (0.089);
\fill[fill={rgb,255:red,174;green,134;blue,117}, fill opacity=0.76] (6.95,4.98) circle (0.072);
\fill[fill={rgb,255:red,239;green,123;blue,63}, fill opacity=0.93] (12.64,4.26) circle (0.092);
\fill[fill={rgb,255:red,239;green,123;blue,63}, fill opacity=0.93] (10.97,4.96) circle (0.093);
\fill[fill={rgb,255:red,238;green,123;blue,63}, fill opacity=0.91] (10.10,5.10) circle (0.090);
\fill[fill={rgb,255:red,215;green,127;blue,81}, fill opacity=0.81] (7.37,2.28) circle (0.078);
\fill[fill={rgb,255:red,223;green,128;blue,69}, fill opacity=0.84] (13.74,6.76) circle (0.082);
\fill[fill={rgb,255:red,239;green,123;blue,63}, fill opacity=0.93] (14.20,4.70) circle (0.092);
\fill[fill={rgb,255:red,237;green,123;blue,63}, fill opacity=0.90] (14.02,3.93) circle (0.089);
\fill[fill={rgb,255:red,89;green,166;blue,145}, fill opacity=0.72] (12.09,8.50) circle (0.066);
\fill[fill={rgb,255:red,235;green,124;blue,66}, fill opacity=0.88] (9.23,3.64) circle (0.086);
\fill[fill={rgb,255:red,234;green,125;blue,65}, fill opacity=0.88] (15.90,5.42) circle (0.087);
\fill[fill={rgb,255:red,237;green,123;blue,64}, fill opacity=0.89] (9.65,4.06) circle (0.088);
\fill[fill={rgb,255:red,186;green,141;blue,84}, fill opacity=0.79] (15.07,6.94) circle (0.075);
\fill[fill={rgb,255:red,238;green,123;blue,63}, fill opacity=0.91] (10.05,4.67) circle (0.090);
\fill[fill={rgb,255:red,221;green,129;blue,70}, fill opacity=0.84] (16.91,4.11) circle (0.082);
\fill[fill={rgb,255:red,222;green,129;blue,70}, fill opacity=0.84] (15.61,3.07) circle (0.081);
\fill[fill={rgb,255:red,239;green,123;blue,63}, fill opacity=0.92] (10.48,4.41) circle (0.091);
\fill[fill={rgb,255:red,152;green,142;blue,124}, fill opacity=0.72] (10.96,7.63) circle (0.067);
\fill[fill={rgb,255:red,234;green,124;blue,66}, fill opacity=0.87] (9.14,3.57) circle (0.085);
\fill[fill={rgb,255:red,204;green,129;blue,90}, fill opacity=0.79] (9.78,6.57) circle (0.076);
\fill[fill={rgb,255:red,235;green,124;blue,64}, fill opacity=0.89] (15.56,4.68) circle (0.088);
\fill[fill={rgb,255:red,236;green,124;blue,64}, fill opacity=0.90] (13.53,3.59) circle (0.088);
\fill[fill={rgb,255:red,221;green,127;blue,76}, fill opacity=0.82] (8.35,1.88) circle (0.079);
\fill[fill={rgb,255:red,223;green,128;blue,69}, fill opacity=0.85] (16.66,5.69) circle (0.082);
\fill[fill={rgb,255:red,239;green,123;blue,63}, fill opacity=0.93] (13.63,5.84) circle (0.093);
\fill[fill={rgb,255:red,239;green,123;blue,63}, fill opacity=0.95] (12.02,5.14) circle (0.094);
\fill[fill={rgb,255:red,229;green,125;blue,69}, fill opacity=0.85] (9.43,2.36) circle (0.082);
\fill[fill={rgb,255:red,239;green,123;blue,63}, fill opacity=0.94] (13.01,4.64) circle (0.094);
\fill[fill={rgb,255:red,239;green,123;blue,63}, fill opacity=0.94] (12.41,5.74) circle (0.093);
\fill[fill={rgb,255:red,233;green,125;blue,65}, fill opacity=0.88] (16.07,5.14) circle (0.086);
\fill[fill={rgb,255:red,237;green,123;blue,64}, fill opacity=0.89] (10.17,5.60) circle (0.088);
\fill[fill={rgb,255:red,52;green,188;blue,138}, fill opacity=0.86] (14.84,8.37) circle (0.083);
\fill[fill={rgb,255:red,176;green,134;blue,115}, fill opacity=0.76] (5.37,2.67) circle (0.072);
\fill[fill={rgb,255:red,88;green,159;blue,165}, fill opacity=0.72] (11.46,8.51) circle (0.066);
\fill[fill={rgb,255:red,238;green,123;blue,64}, fill opacity=0.90] (10.22,5.45) circle (0.089);
\fill[fill={rgb,255:red,239;green,123;blue,63}, fill opacity=0.94] (13.68,4.95) circle (0.094);
\fill[fill={rgb,255:red,239;green,123;blue,63}, fill opacity=0.94] (13.70,4.94) circle (0.094);
\fill[fill={rgb,255:red,239;green,123;blue,63}, fill opacity=0.92] (10.45,4.70) circle (0.091);
\fill[fill={rgb,255:red,237;green,123;blue,63}, fill opacity=0.90] (14.58,4.23) circle (0.089);
\fill[fill={rgb,255:red,239;green,123;blue,63}, fill opacity=0.93] (13.21,4.44) circle (0.093);
\fill[fill={rgb,255:red,219;green,130;blue,72}, fill opacity=0.83] (13.33,6.90) circle (0.080);
\fill[fill={rgb,255:red,238;green,123;blue,63}, fill opacity=0.91] (11.58,5.99) circle (0.090);
\fill[fill={rgb,255:red,239;green,123;blue,63}, fill opacity=0.94] (11.69,5.44) circle (0.093);
\fill[fill={rgb,255:red,239;green,123;blue,63}, fill opacity=0.95] (13.37,5.20) circle (0.095);
\fill[fill={rgb,255:red,238;green,123;blue,63}, fill opacity=0.90] (10.52,3.84) circle (0.089);
\fill[fill={rgb,255:red,225;green,125;blue,74}, fill opacity=0.84] (8.38,4.97) circle (0.082);
\fill[fill={rgb,255:red,237;green,124;blue,64}, fill opacity=0.90] (15.49,5.26) circle (0.089);
\fill[fill={rgb,255:red,234;green,124;blue,65}, fill opacity=0.89] (15.14,6.03) circle (0.087);
\fill[fill={rgb,255:red,142;green,157;blue,102}, fill opacity=0.74] (15.12,7.27) circle (0.069);
\fill[fill={rgb,255:red,219;green,129;blue,71}, fill opacity=0.84] (15.90,3.06) circle (0.081);
\fill[fill={rgb,255:red,225;green,128;blue,69}, fill opacity=0.85] (15.69,3.36) circle (0.082);
\fill[fill={rgb,255:red,223;green,128;blue,70}, fill opacity=0.84] (12.85,6.86) circle (0.081);
\fill[fill={rgb,255:red,238;green,123;blue,64}, fill opacity=0.90] (10.94,3.58) circle (0.089);
\fill[fill={rgb,255:red,60;green,185;blue,135}, fill opacity=0.84] (15.16,8.08) circle (0.081);
\fill[fill={rgb,255:red,239;green,123;blue,63}, fill opacity=0.94] (11.82,5.39) circle (0.094);
\fill[fill={rgb,255:red,60;green,153;blue,214}, fill opacity=0.89] (8.93,8.10) circle (0.087);
\fill[fill={rgb,255:red,226;green,125;blue,74}, fill opacity=0.84] (8.51,5.12) circle (0.082);
\fill[fill={rgb,255:red,238;green,123;blue,64}, fill opacity=0.90] (11.60,6.13) circle (0.089);
\fill[fill={rgb,255:red,237;green,123;blue,64}, fill opacity=0.90] (15.13,4.49) circle (0.089);
\fill[fill={rgb,255:red,239;green,123;blue,63}, fill opacity=0.93] (12.87,4.33) circle (0.093);
\fill[fill={rgb,255:red,239;green,123;blue,63}, fill opacity=0.94] (13.08,5.68) circle (0.094);
\fill[fill={rgb,255:red,156;green,137;blue,132}, fill opacity=0.74] (5.16,3.42) circle (0.069);
\fill[fill={rgb,255:red,239;green,123;blue,63}, fill opacity=0.95] (13.33,5.38) circle (0.095);
\fill[fill={rgb,255:red,228;green,125;blue,69}, fill opacity=0.85] (11.25,6.63) circle (0.082);
\fill[fill={rgb,255:red,238;green,123;blue,63}, fill opacity=0.92] (14.68,4.73) circle (0.091);
\fill[fill={rgb,255:red,204;green,135;blue,77}, fill opacity=0.81] (18.57,5.12) circle (0.078);
\fill[fill={rgb,255:red,239;green,123;blue,63}, fill opacity=0.93] (11.08,5.34) circle (0.092);
\fill[fill={rgb,255:red,229;green,126;blue,67}, fill opacity=0.86] (16.37,5.52) circle (0.084);
\fill[fill={rgb,255:red,238;green,123;blue,63}, fill opacity=0.91] (12.70,6.14) circle (0.090);
\fill[fill={rgb,255:red,238;green,123;blue,63}, fill opacity=0.92] (14.93,5.48) circle (0.091);
\fill[fill={rgb,255:red,199;green,137;blue,79}, fill opacity=0.80] (19.20,4.92) circle (0.077);
\fill[fill={rgb,255:red,219;green,127;blue,78}, fill opacity=0.82] (7.82,2.08) circle (0.079);
\fill[fill={rgb,255:red,239;green,123;blue,63}, fill opacity=0.92] (13.99,5.95) circle (0.091);
\fill[fill={rgb,255:red,239;green,123;blue,63}, fill opacity=0.93] (14.29,4.69) circle (0.092);
\fill[fill={rgb,255:red,239;green,123;blue,63}, fill opacity=0.95] (13.10,5.55) circle (0.094);
\fill[fill={rgb,255:red,234;green,124;blue,65}, fill opacity=0.88] (13.48,3.23) circle (0.086);
\fill[fill={rgb,255:red,239;green,123;blue,63}, fill opacity=0.93] (13.53,4.53) circle (0.093);
\fill[fill={rgb,255:red,236;green,124;blue,64}, fill opacity=0.90] (14.70,6.06) circle (0.088);
\fill[fill={rgb,255:red,230;green,125;blue,69}, fill opacity=0.85] (10.59,6.35) circle (0.083);
\fill[fill={rgb,255:red,205;green,135;blue,77}, fill opacity=0.81] (14.41,6.90) circle (0.078);
\fill[fill={rgb,255:red,167;green,136;blue,122}, fill opacity=0.75] (8.76,6.41) circle (0.070);
\fill[fill={rgb,255:red,226;green,125;blue,73}, fill opacity=0.84] (8.24,4.72) circle (0.082);
\fill[fill={rgb,255:red,230;green,126;blue,66}, fill opacity=0.86] (13.34,6.63) circle (0.084);
\fill[fill={rgb,255:red,187;green,141;blue,84}, fill opacity=0.79] (17.45,6.40) circle (0.075);
\fill[fill={rgb,255:red,239;green,123;blue,63}, fill opacity=0.92] (14.74,5.31) circle (0.092);
\fill[fill={rgb,255:red,200;green,136;blue,78}, fill opacity=0.81] (18.90,5.21) circle (0.077);
\fill[fill={rgb,255:red,231;green,125;blue,66}, fill opacity=0.86] (12.81,6.64) circle (0.084);
\fill[fill={rgb,255:red,237;green,123;blue,64}, fill opacity=0.89] (9.58,4.98) circle (0.088);
\fill[fill={rgb,255:red,239;green,123;blue,63}, fill opacity=0.93] (13.67,4.58) circle (0.093);
\fill[fill={rgb,255:red,239;green,123;blue,63}, fill opacity=0.94] (11.58,5.04) circle (0.094);
\fill[fill={rgb,255:red,158;green,140;blue,123}, fill opacity=0.73] (10.56,7.43) circle (0.068);
\fill[fill={rgb,255:red,237;green,123;blue,63}, fill opacity=0.90] (13.25,3.68) circle (0.089);
\fill[fill={rgb,255:red,229;green,126;blue,67}, fill opacity=0.86] (16.03,4.11) circle (0.084);
\fill[fill={rgb,255:red,83;green,177;blue,126}, fill opacity=0.80] (16.70,7.69) circle (0.077);
\fill[fill={rgb,255:red,237;green,123;blue,63}, fill opacity=0.90] (12.38,3.54) circle (0.089);
\fill[fill={rgb,255:red,211;green,129;blue,83}, fill opacity=0.80] (10.49,6.80) circle (0.077);
\fill[fill={rgb,255:red,239;green,123;blue,63}, fill opacity=0.93] (14.22,4.85) circle (0.093);
\fill[fill={rgb,255:red,238;green,123;blue,63}, fill opacity=0.91] (13.98,6.03) circle (0.091);
\fill[fill={rgb,255:red,217;green,128;blue,78}, fill opacity=0.81] (10.54,6.71) circle (0.078);
\fill[fill={rgb,255:red,153;green,138;blue,134}, fill opacity=0.74] (8.62,6.45) circle (0.069);
\fill[fill={rgb,255:red,193;green,131;blue,101}, fill opacity=0.79] (8.20,5.74) circle (0.074);
\fill[fill={rgb,255:red,195;green,138;blue,81}, fill opacity=0.80] (18.06,6.06) circle (0.076);
\fill[fill={rgb,255:red,239;green,123;blue,63}, fill opacity=0.94] (11.90,4.91) circle (0.094);
\fill[fill={rgb,255:red,238;green,123;blue,64}, fill opacity=0.91] (9.89,4.67) circle (0.089);
\fill[fill={rgb,255:red,78;green,179;blue,128}, fill opacity=0.81] (15.77,7.74) circle (0.077);
\fill[fill={rgb,255:red,238;green,123;blue,63}, fill opacity=0.91] (10.62,5.50) circle (0.090);
\fill[fill={rgb,255:red,208;green,129;blue,87}, fill opacity=0.80] (9.60,6.43) circle (0.076);
\fill[fill={rgb,255:red,229;green,126;blue,67}, fill opacity=0.85] (12.15,2.32) circle (0.083);
\fill[fill={rgb,255:red,238;green,123;blue,63}, fill opacity=0.91] (15.24,5.27) circle (0.090);
\fill[fill={rgb,255:red,217;green,130;blue,72}, fill opacity=0.84] (16.17,6.22) circle (0.081);
\fill[fill={rgb,255:red,235;green,124;blue,66}, fill opacity=0.88] (8.96,4.39) circle (0.086);
\fill[fill={rgb,255:red,48;green,189;blue,139}, fill opacity=0.88] (16.37,8.56) circle (0.086);
\fill[fill={rgb,255:red,237;green,123;blue,64}, fill opacity=0.90] (14.62,4.20) circle (0.089);
\fill[fill={rgb,255:red,239;green,123;blue,63}, fill opacity=0.92] (12.02,5.95) circle (0.091);
\fill[fill={rgb,255:red,197;green,137;blue,80}, fill opacity=0.80] (15.13,6.83) circle (0.076);
\fill[fill={rgb,255:red,238;green,123;blue,63}, fill opacity=0.91] (14.69,4.60) circle (0.091);
\fill[fill={rgb,255:red,236;green,124;blue,64}, fill opacity=0.89] (12.16,6.40) circle (0.087);
\fill[fill={rgb,255:red,232;green,125;blue,66}, fill opacity=0.88] (14.97,6.22) circle (0.086);
\fill[fill={rgb,255:red,58;green,182;blue,147}, fill opacity=0.78] (12.74,9.08) circle (0.074);
\fill[fill={rgb,255:red,104;green,170;blue,118}, fill opacity=0.77] (15.05,7.57) circle (0.073);
\fill[fill={rgb,255:red,239;green,123;blue,63}, fill opacity=0.95] (13.09,5.39) circle (0.095);
\fill[fill={rgb,255:red,228;green,125;blue,70}, fill opacity=0.85] (10.32,6.31) circle (0.082);
\fill[fill={rgb,255:red,222;green,126;blue,76}, fill opacity=0.83] (8.58,5.37) circle (0.080);
\fill[fill={rgb,255:red,238;green,123;blue,63}, fill opacity=0.92] (15.00,5.13) circle (0.091);
\fill[fill={rgb,255:red,237;green,123;blue,64}, fill opacity=0.90] (11.08,6.02) circle (0.088);
\fill[fill={rgb,255:red,207;green,128;blue,89}, fill opacity=0.80] (8.12,5.42) circle (0.077);
\fill[fill={rgb,255:red,129;green,142;blue,154}, fill opacity=0.75] (9.08,6.93) circle (0.070);
\fill[fill={rgb,255:red,238;green,123;blue,63}, fill opacity=0.90] (12.59,3.61) circle (0.089);
\fill[fill={rgb,255:red,220;green,129;blue,70}, fill opacity=0.84] (17.10,4.27) circle (0.081);
\fill[fill={rgb,255:red,223;green,126;blue,76}, fill opacity=0.83] (7.68,3.76) circle (0.081);
\fill[fill={rgb,255:red,213;green,127;blue,83}, fill opacity=0.82] (8.65,5.70) circle (0.078);
\fill[fill={rgb,255:red,98;green,147;blue,182}, fill opacity=0.80] (6.96,6.00) circle (0.076);
\fill[fill={rgb,255:red,239;green,123;blue,63}, fill opacity=0.95] (13.27,5.33) circle (0.095);
\fill[fill={rgb,255:red,239;green,123;blue,63}, fill opacity=0.94] (12.58,5.71) circle (0.094);
\fill[fill={rgb,255:red,238;green,123;blue,63}, fill opacity=0.92] (14.68,5.73) circle (0.091);
\fill[fill={rgb,255:red,222;green,126;blue,76}, fill opacity=0.83] (7.76,2.69) circle (0.080);
\fill[fill={rgb,255:red,239;green,123;blue,63}, fill opacity=0.93] (12.17,4.29) circle (0.093);
\fill[fill={rgb,255:red,234;green,124;blue,65}, fill opacity=0.88] (12.12,6.49) circle (0.086);
\fill[fill={rgb,255:red,234;green,124;blue,65}, fill opacity=0.88] (12.96,6.51) circle (0.086);
\fill[fill={rgb,255:red,105;green,170;blue,117}, fill opacity=0.78] (18.06,7.62) circle (0.073);
\fill[fill={rgb,255:red,221;green,129;blue,70}, fill opacity=0.84] (14.82,6.57) circle (0.082);
\fill[fill={rgb,255:red,227;green,127;blue,68}, fill opacity=0.85] (12.71,2.24) circle (0.082);
\fill[fill={rgb,255:red,230;green,126;blue,66}, fill opacity=0.87] (15.40,6.14) circle (0.085);
\fill[fill={rgb,255:red,196;green,130;blue,98}, fill opacity=0.79] (7.62,5.19) circle (0.075);
\fill[fill={rgb,255:red,238;green,123;blue,63}, fill opacity=0.91] (12.26,3.66) circle (0.090);
\fill[fill={rgb,255:red,65;green,152;blue,210}, fill opacity=0.87] (7.78,7.13) circle (0.084);
\fill[fill={rgb,255:red,239;green,123;blue,63}, fill opacity=0.95] (12.83,5.24) circle (0.095);
\fill[fill={rgb,255:red,123;green,163;blue,109}, fill opacity=0.76] (18.19,7.35) circle (0.071);
\fill[fill={rgb,255:red,73;green,151;blue,203}, fill opacity=0.84] (3.69,5.20) circle (0.082);
\fill[fill={rgb,255:red,213;green,127;blue,84}, fill opacity=0.81] (8.55,5.63) circle (0.078);
\fill[fill={rgb,255:red,86;green,176;blue,124}, fill opacity=0.80] (15.98,7.63) circle (0.076);
\fill[fill={rgb,255:red,215;green,131;blue,73}, fill opacity=0.83] (17.08,5.82) circle (0.080);
\fill[fill={rgb,255:red,234;green,124;blue,66}, fill opacity=0.87] (11.81,6.50) circle (0.085);
\fill[fill={rgb,255:red,239;green,123;blue,63}, fill opacity=0.92] (11.43,5.82) circle (0.091);
\fill[fill={rgb,255:red,233;green,125;blue,65}, fill opacity=0.88] (15.67,4.27) circle (0.086);
\fill[fill={rgb,255:red,239;green,123;blue,63}, fill opacity=0.94] (12.80,4.62) circle (0.094);
\fill[fill={rgb,255:red,217;green,130;blue,72}, fill opacity=0.83] (17.50,4.55) circle (0.081);
\fill[fill={rgb,255:red,231;green,125;blue,66}, fill opacity=0.87] (14.85,3.45) circle (0.085);
\fill[fill={rgb,255:red,213;green,132;blue,74}, fill opacity=0.83] (17.70,4.11) circle (0.079);
\fill[fill={rgb,255:red,239;green,123;blue,63}, fill opacity=0.92] (12.21,5.95) circle (0.092);
\fill[fill={rgb,255:red,201;green,130;blue,93}, fill opacity=0.79] (6.39,2.77) circle (0.075);
\fill[fill={rgb,255:red,184;green,142;blue,85}, fill opacity=0.78] (15.03,6.97) circle (0.074);
\fill[fill={rgb,255:red,239;green,123;blue,63}, fill opacity=0.95] (13.68,5.22) circle (0.094);
\fill[fill={rgb,255:red,207;green,128;blue,89}, fill opacity=0.81] (7.20,4.34) circle (0.077);
\fill[fill={rgb,255:red,222;green,128;blue,70}, fill opacity=0.85] (17.00,4.41) circle (0.082);
\fill[fill={rgb,255:red,239;green,123;blue,63}, fill opacity=0.93] (11.71,5.76) circle (0.092);
\fill[fill={rgb,255:red,216;green,131;blue,72}, fill opacity=0.83] (17.67,4.92) circle (0.080);
\fill[fill={rgb,255:red,237;green,123;blue,64}, fill opacity=0.90] (11.00,6.00) circle (0.088);
\fill[fill={rgb,255:red,218;green,130;blue,71}, fill opacity=0.84] (14.95,6.59) circle (0.081);
\fill[fill={rgb,255:red,235;green,124;blue,65}, fill opacity=0.88] (11.27,3.06) circle (0.087);
\fill[fill={rgb,255:red,237;green,123;blue,63}, fill opacity=0.90] (13.87,6.18) circle (0.089);
\fill[fill={rgb,255:red,226;green,127;blue,68}, fill opacity=0.85] (15.18,3.13) circle (0.083);
\fill[fill={rgb,255:red,230;green,126;blue,66}, fill opacity=0.86] (13.72,6.59) circle (0.084);
\fill[fill={rgb,255:red,220;green,127;blue,77}, fill opacity=0.82] (7.94,2.10) circle (0.079);
\fill[fill={rgb,255:red,139;green,140;blue,146}, fill opacity=0.74] (4.29,2.95) circle (0.069);
\fill[fill={rgb,255:red,187;green,132;blue,105}, fill opacity=0.77] (9.47,6.63) circle (0.073);
\fill[fill={rgb,255:red,236;green,123;blue,65}, fill opacity=0.88] (9.56,5.27) circle (0.087);
\fill[fill={rgb,255:red,166;green,148;blue,92}, fill opacity=0.77] (16.70,6.83) circle (0.072);
\fill[fill={rgb,255:red,216;green,130;blue,72}, fill opacity=0.83] (15.72,2.72) circle (0.080);
\fill[fill={rgb,255:red,73;green,151;blue,203}, fill opacity=0.84] (3.26,5.13) circle (0.081);
\fill[fill={rgb,255:red,209;green,128;blue,87}, fill opacity=0.81] (6.93,3.58) circle (0.077);
\fill[fill={rgb,255:red,73;green,151;blue,203}, fill opacity=0.84] (7.69,6.84) circle (0.081);
\fill[fill={rgb,255:red,212;green,128;blue,85}, fill opacity=0.81] (7.67,4.80) circle (0.078);
\fill[fill={rgb,255:red,216;green,131;blue,72}, fill opacity=0.83] (16.97,5.85) circle (0.080);
\fill[fill={rgb,255:red,189;green,140;blue,83}, fill opacity=0.79] (20.58,4.41) circle (0.075);
\fill[fill={rgb,255:red,239;green,123;blue,63}, fill opacity=0.93] (12.75,4.26) circle (0.092);
\fill[fill={rgb,255:red,239;green,123;blue,63}, fill opacity=0.95] (13.46,5.47) circle (0.094);
\fill[fill={rgb,255:red,67;green,152;blue,208}, fill opacity=0.86] (6.18,6.27) circle (0.084);
\fill[fill={rgb,255:red,213;green,131;blue,75}, fill opacity=0.81] (12.84,7.05) circle (0.078);
\fill[fill={rgb,255:red,236;green,124;blue,64}, fill opacity=0.89] (15.57,5.46) circle (0.088);
\fill[fill={rgb,255:red,237;green,123;blue,64}, fill opacity=0.89] (10.04,3.58) circle (0.087);
\fill[fill={rgb,255:red,238;green,123;blue,63}, fill opacity=0.91] (14.05,4.01) circle (0.090);
\fill[fill={rgb,255:red,180;green,141;blue,93}, fill opacity=0.76] (12.23,7.50) circle (0.071);
\fill[fill={rgb,255:red,230;green,126;blue,66}, fill opacity=0.87] (15.38,6.15) circle (0.085);
\fill[fill={rgb,255:red,231;green,125;blue,66}, fill opacity=0.87] (16.10,4.51) circle (0.085);
\fill[fill={rgb,255:red,71;green,181;blue,130}, fill opacity=0.82] (16.81,7.89) circle (0.079);
\fill[fill={rgb,255:red,86;green,149;blue,192}, fill opacity=0.82] (4.87,5.06) circle (0.078);
\fill[fill={rgb,255:red,205;green,135;blue,77}, fill opacity=0.81] (15.18,6.73) circle (0.078);
\fill[fill={rgb,255:red,226;green,127;blue,68}, fill opacity=0.85] (13.88,6.67) circle (0.083);
\fill[fill={rgb,255:red,239;green,123;blue,63}, fill opacity=0.93] (11.87,5.74) circle (0.092);
\fill[fill={rgb,255:red,239;green,123;blue,63}, fill opacity=0.93] (10.95,4.54) circle (0.092);
\fill[fill={rgb,255:red,234;green,125;blue,65}, fill opacity=0.88] (14.54,3.57) circle (0.086);
\fill[fill={rgb,255:red,236;green,124;blue,64}, fill opacity=0.89] (14.62,3.96) circle (0.088);
\fill[fill={rgb,255:red,239;green,123;blue,63}, fill opacity=0.93] (13.54,4.40) circle (0.092);
\fill[fill={rgb,255:red,239;green,123;blue,63}, fill opacity=0.92] (12.48,6.05) circle (0.091);
\fill[fill={rgb,255:red,213;green,132;blue,73}, fill opacity=0.83] (17.80,5.16) circle (0.080);
\fill[fill={rgb,255:red,173;green,146;blue,89}, fill opacity=0.77] (15.94,6.88) circle (0.073);
\fill[fill={rgb,255:red,115;green,144;blue,167}, fill opacity=0.77] (7.01,5.82) circle (0.073);
\fill[fill={rgb,255:red,147;green,154;blue,101}, fill opacity=0.74] (13.82,7.50) circle (0.069);
\fill[fill={rgb,255:red,75;green,180;blue,129}, fill opacity=0.81] (17.74,8.07) circle (0.078);
\fill[fill={rgb,255:red,212;green,132;blue,74}, fill opacity=0.82] (15.56,2.37) circle (0.079);
\fill[fill={rgb,255:red,65;green,152;blue,210}, fill opacity=0.87] (3.72,6.29) circle (0.085);
\fill[fill={rgb,255:red,215;green,131;blue,72}, fill opacity=0.83] (16.34,6.20) circle (0.080);
\fill[fill={rgb,255:red,239;green,123;blue,63}, fill opacity=0.94] (13.60,5.75) circle (0.093);
\fill[fill={rgb,255:red,209;green,133;blue,75}, fill opacity=0.82] (18.00,5.35) circle (0.079);
\fill[fill={rgb,255:red,239;green,123;blue,63}, fill opacity=0.92] (11.28,5.59) circle (0.092);
\fill[fill={rgb,255:red,231;green,124;blue,69}, fill opacity=0.86] (8.85,5.09) circle (0.084);
\fill[fill={rgb,255:red,150;green,138;blue,137}, fill opacity=0.74] (5.11,3.58) circle (0.069);
\fill[fill={rgb,255:red,235;green,124;blue,65}, fill opacity=0.88] (12.34,3.03) circle (0.086);
\fill[fill={rgb,255:red,81;green,178;blue,126}, fill opacity=0.80] (16.89,7.74) circle (0.077);
\fill[fill={rgb,255:red,198;green,137;blue,79}, fill opacity=0.80] (15.98,6.61) circle (0.077);
\fill[fill={rgb,255:red,239;green,123;blue,63}, fill opacity=0.93] (13.27,4.26) circle (0.092);
\fill[fill={rgb,255:red,239;green,123;blue,63}, fill opacity=0.94] (14.15,5.05) circle (0.093);
\draw[white, line width=0.8pt] (8.07,8.87) circle (0.16);
\fill[white] (8.07,8.87) circle (0.045);
\draw[white, line width=0.8pt] (12.86,5.19) circle (0.16);
\fill[white] (12.86,5.19) circle (0.045);
\draw[white, line width=0.8pt] (15.10,10.57) circle (0.16);
\fill[white] (15.10,10.57) circle (0.045);
\end{scope}
    \end{scope}
    \fill[fcnavy, path fading=artfade, fading angle=180] (0,22.6) rectangle (21,25.4);
    \node[anchor=north west, inner sep=0, text width=15.2cm, text=white, align=justify,
          font=\fontsize{11.5}{16}\selectfont] at (2.9,20.6) {%
      {\fcbold\fontsize{15}{19}\selectfont From the c-means functional to model identification\par}
      \vspace{10pt}
      Fuzzy clustering assigns every observation a degree of membership in every group instead of forcing it into one. This book develops the main families of fuzzy clustering algorithms in one notation and shows, for each of them, how to run it in Python and how its parameters change the result.\par\vspace{6pt}
      The core chapters treat fuzzy c-means and the Gustafson--Kessel and Gath--Geva algorithms, cluster validity, visualization by fuzzy Sammon mapping, and possibilistic, kernel, shell, and graph-based variants. Later chapters cover Takagi--Sugeno model identification, time-series segmentation, relational and mixed data, missing values, very large data, image segmentation, partial supervision, co-clustering, collaborative clustering, and cluster ensembles.\par\vspace{6pt}
      Every method is implemented in {\plexmono fuzzyclust}, an open-source Python package with a scikit-learn interface and the successor of the \emph{Fuzzy Clustering and Data Analysis Toolbox} for MATLAB. Each figure in the book is produced by an example script that generates its own data; the same scripts run as Jupyter notebooks.};
    \node[anchor=south west, inner sep=0, text=fclight, align=left, font=\fcmedium\fontsize{10}{14}\selectfont] at (2.9,2.0)
      {{\plexmono pip install fuzzyclust}\\github.com/ComplexSystemsIntelligence/fuzzyclust};
    \node[anchor=south east, inner sep=0] at (18.1,1.9) {\fclogo[0.42]};
  \end{scope}
\end{tikzpicture}
\end{document}
